% Redacción : carta abierta / petición ; presentación oral de un país hispanohablante — Espagnol Tle A — Cours signé OnBuch+
% Programme officiel MINESEC. Compiler avec : tectonic E19-carta-abierta-peticion-pais-hispanohablante.tex
\newcommand{\DOCMATIERE}{Espagnol}
\newcommand{\DOCNIVEAU}{Tle A}
\newcommand{\DOCMODULE}{Módulo 4 : Activités économiques et monde du travail}
\newcommand{\DOCLECON}{E19}
\newcommand{\DOCTITRE}{Redacción : carta abierta / petición ; presentación oral de un país hispanohablante}
% OnBuch+ — préambule « cours signés » ENRICHI (compilé avec tectonic / XeLaTeX)
% Système visuel « Pop » d'OnBuch+ : fond crème (jamais blanc), bordures encre
% épaisses, ombres « dures » décalées, titres Archivo Black en capitales.
% Version MATHÉMATIQUES : graphes (pgfplots/tikz), boîtes pédagogiques
% (définition, propriété, méthode, attention, le savais-tu, exercice résolu,
% exercice, TP, bilan), page de garde et signature OnBuch+ ancrée en bas.
%
% Macros attendues dans main.tex AVANT \input{preamble} :
%   \DOCMATIERE  \DOCNIVEAU  \DOCMODULE  \DOCLECON (numéro)  \DOCTITRE
\documentclass[11pt]{article}

\usepackage{fontspec}
\usepackage[spanish,french]{babel} % français = langue principale ; l'espagnol via \esp{..} / espagnol
\usepackage[a4paper,margin=2cm,top=3.4cm,bottom=2.8cm,headheight=1.4cm,footskip=1.3cm]{geometry}
\usepackage{amsmath,amssymb,amsfonts}
\usepackage{mathtools} % \xmapsto, \coloneqq... souvent utilisés par les modèles
\usepackage{cancel} % simplifications barrées (bilans de forces)
% \abs{z} (module, valeur absolue) et \norm{u} (norme), utilisés par les modèles
\providecommand{\abs}{}\let\abs\relax\DeclarePairedDelimiter{\abs}{\lvert}{\rvert}
\providecommand{\norm}{}\let\norm\relax\DeclarePairedDelimiter{\norm}{\lVert}{\rVert}
\usepackage[table]{xcolor} % [table] : fournit aussi \rowcolors (alternance de lignes)
\usepackage{pagecolor}
\usepackage{tikz}
\usetikzlibrary{arrows.meta,positioning,calc,shapes.geometric,shapes.misc,shapes.arrows,decorations.pathmorphing,decorations.pathreplacing,decorations.markings,patterns,shadows,mindmap,trees,fit,backgrounds,matrix,intersections,angles,quotes}
\usepackage{pgfplots}
\usepackage[version=4]{mhchem} % équations nucléaires \ce{^{235}_{92}U}
\usepackage[european,siunitx]{circuitikz} % schémas électriques
\pgfplotsset{compat=1.18}
\usepgfplotslibrary{fillbetween,groupplots} % aires entre courbes (calcul intégral)
\usepackage{enumitem}
\usepackage{fancyhdr}
% [most] charge toutes les bibliothèques tcolorbox (listings, minted, raster,
% documentation...) et gonfle inutilement la mémoire TeX sur un document long
% (~300 boîtes) : on ne charge que ce qui est réellement utilisé.
\usepackage[skins,breakable]{tcolorbox}
\usepackage{graphicx}
\usepackage{colortbl}
\usepackage{booktabs}
\usepackage{tabularx}
\usepackage{diagbox} % case d'en-tête diagonale des tableaux à double entrée
% Colonnes extensibles centrée / gauche / droite (C, L, R), souvent utilisées par les modèles
\newcolumntype{C}{>{\centering\arraybackslash}X}
\newcolumntype{L}{>{\raggedright\arraybackslash}X}
\newcolumntype{R}{>{\raggedleft\arraybackslash}X}
\usepackage{array}
\usepackage{multicol}
\usepackage{siunitx}
% parse-numbers=false : accepte aussi \SI{2,5 \times 10^{-3}}{...}
\sisetup{locale=FR,output-decimal-marker={,},group-minimum-digits=4,parse-numbers=false}
\DeclareSIUnit\ohm{\ensuremath{\symup{\Omega}}} % U+2126 absent de la police
\DeclareSIUnit\division{div} % réglages d'oscilloscope (ms/div, V/div)
\DeclareSIUnit\tr{tr} % tours (vitesse de rotation, ex. \SI{1500}{\tr\per\minute})
\DeclareSIUnit\molar{M} % concentration molaire (ex. \SI{3}{\molar}, \SI{1}{\micro\molar})
\DeclareSIUnit\an{an} % année (le modèle confond parfois avec \year, un compteur TeX) — ex. \SI{5}{\kilo\meter\per\an}
\DeclareSIUnit\FCFA{FCFA} % franc CFA (ex. \SI{500}{\FCFA\per\kilo\gram})
\DeclareSIUnit\degreeBrix{\textdegree Brix} % teneur en sucre d'un sirop/jus (réfractométrie)
\DeclareSIUnit\month{mois} % le modèle utilise parfois \month, un compteur TeX, comme unité de durée
\usepackage{qrcode}
% \degree : commande fréquemment utilisée par les modèles pour un angle ou une
% température en dehors de \SI{}{} (non fournie par les paquets chargés).
\providecommand{\degree}{^{\circ}}
% \qed (fin de démonstration), utilisé par les modèles sans amsthm
\providecommand{\qed}{\hfill$\square$}
% \barre : double barre des valeurs interdites dans un tableau de variations (tkz-tab)
\providecommand{\barre}{\ensuremath{\|}}
% environnement proof (amsthm non chargé) utilisé par les modèles
\ifdefined\proof\else\newenvironment{proof}[1][Démonstration]{\par\noindent\textit{#1.}\ }{\qed\par}\fi
\usepackage{lastpage}
\usepackage{hyperref}
\hypersetup{hidelinks}

% ============================================================
% Palette « Pop » OnBuch+ — mêmes hex que la classe P (plus_ui.dart)
% ============================================================
\definecolor{poporange}{HTML}{F59321}
\definecolor{popdark}{HTML}{E07A0C}
\definecolor{popcream}{HTML}{FFF6E8}
\definecolor{popink}{HTML}{1C1714}
\definecolor{poppurple}{HTML}{7A5AE0}
\definecolor{popgreen}{HTML}{2E9E6A}
\definecolor{poppink}{HTML}{E0457B}
\definecolor{popblue}{HTML}{2F7DE1}
\definecolor{popgold}{HTML}{C98A12}
\definecolor{popmuted}{HTML}{8A7F76}
\definecolor{popline}{HTML}{EADFD2}
\definecolor{popgray}{HTML}{8A7F76}
\colorlet{popgrey}{popgray}
% Alias tolérants (le modèle nomme parfois les couleurs différemment)
\colorlet{popwhite}{white}
\colorlet{popblack}{popink}
\colorlet{popred}{poppink}
\colorlet{popyellow}{popgold}
% Teintes claires (fonds de tableaux/encadrés) — le modèle nomme parfois
% les couleurs avec un suffixe L (ex. popblueL) sans les avoir définies.
\colorlet{poporangeL}{poporange!20}
\colorlet{popdarkL}{popdark!20}
\colorlet{poppurpleL}{poppurple!20}
\colorlet{popgreenL}{popgreen!20}
\colorlet{poppinkL}{poppink!20}
\colorlet{popblueL}{popblue!20}
\colorlet{popgoldL}{popgold!20}
\colorlet{popredL}{popred!20}
\colorlet{popgrayL}{popgray!20}
% Teintes claires pour fonds de tableaux / zones
\colorlet{popblueL}{popblue!10!white}
\colorlet{poporangeL}{poporange!14!white}
\colorlet{popgreenL}{popgreen!12!white}
\colorlet{poppurpleL}{poppurple!12!white}
\colorlet{poppinkL}{poppink!12!white}
\colorlet{popgoldL}{popgold!14!white}

\pagecolor{popcream}

% ============================================================
% Typographie
% ============================================================
\defaultfontfeatures{Ligatures=TeX}
\setmainfont{PlusJakartaSans-Medium.ttf}[Path=../../../fonts/,BoldFont=PlusJakartaSans-Bold.ttf]
% Maths en sans-serif (Fira Math) avec chiffres et lettres droites de Plus
% Jakarta, pour que les formules se fondent dans le texte.
\usepackage[math-style=ISO,bold-style=ISO]{unicode-math}
\setmathfont{FiraMath-Regular.otf}[Path=../../../fonts/]
\setmathfont{PlusJakartaSans-Medium.ttf}[Path=../../../fonts/,range={up/num,up/{Latin,latin},\mathup}]
\usepackage{icomma} % virgule décimale sans espace en mode math
% Caractères mathématiques Unicode absents des polices de texte (Plus Jakarta,
% Archivo Black) : rendus par leur équivalent mathématique, en texte comme en
% formule — sinon ils s'affichent en carré vide (ex. « Arithmétique dans ℤ »).
\usepackage{newunicodechar}
\newunicodechar{ℝ}{\ensuremath{\mathbb{R}}}
\newunicodechar{ℤ}{\ensuremath{\mathbb{Z}}}
\newunicodechar{ℂ}{\ensuremath{\mathbb{C}}}
\newunicodechar{ℕ}{\ensuremath{\mathbb{N}}}
\newunicodechar{ℚ}{\ensuremath{\mathbb{Q}}}
\newunicodechar{𝔻}{\ensuremath{\mathbb{D}}}
\newunicodechar{α}{\ensuremath{\alpha}}
\newunicodechar{β}{\ensuremath{\beta}}
\newunicodechar{γ}{\ensuremath{\gamma}}
\newunicodechar{δ}{\ensuremath{\delta}}
\newunicodechar{ε}{\ensuremath{\varepsilon}}
\newunicodechar{θ}{\ensuremath{\theta}}
\newunicodechar{λ}{\ensuremath{\lambda}}
\newunicodechar{μ}{\ensuremath{\mu}}
\newunicodechar{σ}{\ensuremath{\sigma}}
\newunicodechar{τ}{\ensuremath{\tau}}
\newunicodechar{φ}{\ensuremath{\varphi}}
\newunicodechar{ψ}{\ensuremath{\psi}}
\newunicodechar{ω}{\ensuremath{\omega}}
\newunicodechar{Σ}{\ensuremath{\Sigma}}
% majuscules produites par \MakeUppercase dans les titres (α-aminés -> Α)
\newunicodechar{Α}{\ensuremath{\alpha}}
\newunicodechar{Β}{\ensuremath{\beta}}
\newunicodechar{Γ}{\ensuremath{\Gamma}}
\newunicodechar{Δ}{\ensuremath{\Delta}}
\newunicodechar{Ε}{\ensuremath{\varepsilon}}
\newunicodechar{Θ}{\ensuremath{\Theta}}
\newunicodechar{Λ}{\ensuremath{\Lambda}}
\newunicodechar{Μ}{\ensuremath{\mu}}
\newunicodechar{Τ}{\ensuremath{\tau}}
\newunicodechar{Φ}{\ensuremath{\Phi}}
\newunicodechar{Ψ}{\ensuremath{\Psi}}
\newunicodechar{Ω}{\ensuremath{\Omega}}
\newunicodechar{↦}{\ensuremath{\mapsto}}
\newunicodechar{∈}{\ensuremath{\in}}
\newunicodechar{∉}{\ensuremath{\notin}}
\newunicodechar{∧}{\ensuremath{\wedge}}
\newunicodechar{⇔}{\ensuremath{\Leftrightarrow}}
\newunicodechar{⟺}{\ensuremath{\Longleftrightarrow}}
\newunicodechar{⇒}{\ensuremath{\Rightarrow}}
\newunicodechar{⟹}{\ensuremath{\Longrightarrow}}
\newunicodechar{∘}{\ensuremath{\circ}}
\newunicodechar{∪}{\ensuremath{\cup}}
\newunicodechar{∩}{\ensuremath{\cap}}
\newunicodechar{≡}{\ensuremath{\equiv}}
\newunicodechar{⊂}{\ensuremath{\subset}}
\newunicodechar{⊥}{\ensuremath{\perp}}
\newunicodechar{∃}{\ensuremath{\exists}}
\newunicodechar{∀}{\ensuremath{\forall}}
\newunicodechar{∛}{\ensuremath{\sqrt[3]{\,}}}
\newunicodechar{∜}{\ensuremath{\sqrt[4]{\,}}}
\newunicodechar{⃗}{\ensuremath{{}^{\to}}}
\newunicodechar{ⁿ}{\ifmmode^{n}\else\textsuperscript{\ensuremath{n}}\fi}
\newunicodechar{ˣ}{\ifmmode^{x}\else\textsuperscript{\ensuremath{x}}\fi}
\newunicodechar{ᵇ}{\ifmmode^{b}\else\textsuperscript{\ensuremath{b}}\fi}
\newunicodechar{ʳ}{\ifmmode^{r}\else\textsuperscript{\ensuremath{r}}\fi}
\newunicodechar{ᵘ}{\ifmmode^{u}\else\textsuperscript{\ensuremath{u}}\fi}
\newunicodechar{ʸ}{\ifmmode^{y}\else\textsuperscript{\ensuremath{y}}\fi}
\newunicodechar{ᵃ}{\ifmmode^{a}\else\textsuperscript{\ensuremath{a}}\fi}
\newunicodechar{ᵉ}{\ifmmode^{e}\else\textsuperscript{\ensuremath{e}}\fi}
\newunicodechar{ᵏ}{\ifmmode^{k}\else\textsuperscript{\ensuremath{k}}\fi}
\newunicodechar{ᵖ}{\ifmmode^{p}\else\textsuperscript{\ensuremath{p}}\fi}
\newunicodechar{ᴺ}{\ifmmode^{N}\else\textsuperscript{\ensuremath{N}}\fi}
\newunicodechar{ᶜ}{\ifmmode^{c}\else\textsuperscript{\ensuremath{c}}\fi}
\newunicodechar{ʰ}{\ifmmode^{h}\else\textsuperscript{\ensuremath{h}}\fi}
\newunicodechar{ᵠ}{\ifmmode^{q}\else\textsuperscript{\ensuremath{q}}\fi}
\newunicodechar{ₙ}{\ifmmode_{n}\else\textsubscript{\ensuremath{n}}\fi}
\newunicodechar{ₐ}{\ifmmode_{a}\else\textsubscript{\ensuremath{a}}\fi}
\newunicodechar{ᵢ}{\ifmmode_{i}\else\textsubscript{\ensuremath{i}}\fi}
\newunicodechar{ᵣ}{\ifmmode_{r}\else\textsubscript{\ensuremath{r}}\fi}
\newunicodechar{ₖ}{\ifmmode_{k}\else\textsubscript{\ensuremath{k}}\fi}
\newunicodechar{ᵦ}{\ifmmode_{\beta}\else\textsubscript{\ensuremath{\beta}}\fi}
\newunicodechar{ₓ}{\ifmmode_{x}\else\textsubscript{\ensuremath{x}}\fi}
\newfontfamily\popdisplay{ArchivoBlack-Regular.ttf}[Path=../../../fonts/]
\newfontfamily\poplogo{SpaceGrotesk-Bold.ttf}[Path=../../../fonts/]
\newfontfamily\popsemi{PlusJakartaSans-SemiBold.ttf}[Path=../../../fonts/]
\newfontfamily\popxbold{PlusJakartaSans-ExtraBold.ttf}[Path=../../../fonts/]
\newfontfamily\popxboldls{PlusJakartaSans-ExtraBold.ttf}[Path=../../../fonts/,LetterSpace=3.0]

\setlist[enumerate]{leftmargin=1.7em,itemsep=5pt,topsep=4pt}
\setlist[itemize]{leftmargin=1.7em,itemsep=4pt,topsep=4pt,label={\color{poporange}\textbullet}}
\setlength{\parindent}{0pt}
\setlength{\parskip}{5pt}
\renewcommand{\arraystretch}{1.25}

% pgfplots : style Pop par défaut
\pgfplotsset{
  every axis/.append style={axis line style={popink,line width=1pt},tick style={popink},
    grid style={popline,line width=0.5pt},label style={font=\small},tick label style={font=\footnotesize}},
  popcurve/.style={poporange,line width=1.8pt,no markers,smooth},
}
% Même style disponible dans un \draw TikZ simple (hors pgfplots)
\tikzset{popcurve/.style={poporange,line width=1.8pt}}
\tikzset{popgrid/.style={popline,very thin}}
\colorlet{popcurve}{poporange} % le nom du style est parfois utilisé comme couleur

% ============================================================
% Pill / squiggle
% ============================================================
\newcommand{\poppill}[3]{%
  \tikz[baseline=-0.55ex]\node[draw=popink,line width=1.2pt,rounded corners=2.6pt,%
    fill=#2,inner xsep=6.5pt,inner ysep=3pt]{%
    \popxboldls\fontsize{7}{7}\selectfont\color{#3}\MakeUppercase{#1}};%
}
\newcommand{\popsquiggle}[1]{%
  \tikz[baseline=-0.4ex]{
    \draw[#1,line width=2.6pt,line cap=round]
      plot [smooth] coordinates {(0,0.09) (0.34,0.01) (0.68,0.21) (1.02,0.01) (1.36,0.10)};
  }%
}
% Mot-clé surligné (marqueur orange sous le texte)
\newcommand{\cle}[1]{{\popxbold\color{popdark}#1}}

% ============================================================
% En-tête / pied
% ============================================================
\pagestyle{fancy}
\fancyhf{}
\renewcommand{\headrulewidth}{0pt}
\renewcommand{\footrulewidth}{0pt}
\lhead{\raisebox{-0.15ex}{\poplogo\fontsize{13}{13}\selectfont\color{popink}OnBuch\color{poporange}+}}
\rhead{\poppill{\DOCMATIERE\ \textbullet\ \DOCNIVEAU\ \textbullet\ Leçon \DOCLECON}{white}{popink}}
\lfoot{\popxbold\fontsize{7.6}{7.6}\selectfont\color{popmuted}\MakeUppercase{Cours signé OnBuch+}\ \textbullet\ {\popsemi\DOCTITRE}}
\rfoot{\tikz[baseline=-0.6ex]\node[rounded corners=8.5pt,fill=popink,minimum height=17pt,minimum width=17pt,inner xsep=5pt,inner sep=0]{\popxbold\fontsize{8}{8}\selectfont\color{popcream}\thepage\,/\,\pageref*{LastPage}};}

% ============================================================
% Titres
% ============================================================
\newcommand{\chaptitre}[2]{%
  \begin{tcolorbox}[enhanced,colback=poporange,colframe=popink,boxrule=2.6pt,arc=11pt,
      drop shadow=popink,left=15pt,right=15pt,top=12pt,bottom=11pt,before skip=3pt,after skip=18pt]
    \poppill{#2}{popink}{popcream}\par\vspace{9pt}
    {\raggedright\hyphenpenalty=10000\exhyphenpenalty=10000\popdisplay\fontsize{22}{24}\selectfont\color{popcream}\MakeUppercase{#1}\par}
  \end{tcolorbox}
}
% Section numérotée : pastille orange avec le numéro + titre Archivo
\newcounter{popsec}
\newcommand{\coursec}[1]{%
  \stepcounter{popsec}\setcounter{popsub}{0}%
  \par\vspace{16pt}\noindent
  \begin{minipage}{\linewidth}
  \tikz[baseline=-0.7ex]\node[draw=popink,line width=1.6pt,fill=poporange,rounded corners=4pt,minimum size=22pt,inner sep=0]{\popdisplay\fontsize{12}{12}\selectfont\color{popcream}\thepopsec};%
  \hspace{8pt}{\popdisplay\fontsize{15}{17}\selectfont\color{popink}\MakeUppercase{#1}}\par
  \vspace{4pt}\noindent{\color{poporange}\rule{36pt}{3.6pt}}
  \end{minipage}\par\nopagebreak\vspace{8pt}
}
\newcounter{popsub}
\newcommand{\courssub}[1]{%
  \stepcounter{popsub}%
  \par\vspace{9pt}\noindent{\popxbold\fontsize{11.5}{13}\selectfont\color{popink}{\color{poporange}\thepopsec.\thepopsub}\hspace{6pt}#1}\par\nopagebreak\vspace{4pt}
}

% ============================================================
% Boîtes pédagogiques — même recette : fond blanc, bordure épaisse colorée,
% ombre dure encre, badge plein. Titre optionnel [#1].
% ============================================================
\tcbset{popbase/.style={breakable,enhanced,colback=white,boxrule=2.3pt,arc=9pt,
  shadow={3pt}{-3pt}{0pt}{popink},left=14pt,right=14pt,top=11pt,bottom=11pt,before skip=11pt,after skip=15pt,
  fontupper=\popsemi\fontsize{10.6}{14.5}\selectfont\color{popink},
  fontlower=\popsemi\fontsize{10.6}{14.5}\selectfont\color{popink}}}
% Coche dessinée (le glyphe ✓ n'existe pas dans les polices du document)
\newcommand{\popcheck}[1]{\tikz[baseline=-0.5ex]\draw[#1,line width=1.6pt,line cap=round,line join=round](-0.13,0.0)--(-0.04,-0.09)--(0.13,0.1);}
\providecommand{\coche}{\popcheck{popgreen}}
\providecommand{\resultat}[1]{\par\smallskip\noindent\fcolorbox{popgreen}{popgreenL}{\parbox{\dimexpr\linewidth-2\fboxsep-2\fboxrule\relax}{\textbf{Résultat.} #1}}\par\smallskip}
\newcommand{\popboxhead}[3]{\poppill{#1}{#2}{white}\ifx\relax#3\relax\else\hspace{6pt}{\popxbold\fontsize{10.6}{12}\selectfont\color{popink}#3}\fi\par\vspace{7pt}}

\newtcolorbox{aretenir}[1][]{popbase,colframe=popgreen,before upper={\popboxhead{À retenir}{popgreen}{#1}}}
% Texte en espagnol (document de lecture, dialogue, poème, extrait) : cadre
% d'étude. Un titre optionnel (source, auteur) s'affiche dans l'en-tête.
\newtcolorbox{textebox}[1][]{popbase,colframe=popblue,before upper={\popboxhead{Texto}{popblue}{#1}\begin{otherlanguage}{spanish}},after upper={\end{otherlanguage}}}
% Marques ✓ / ✗ (forme correcte / fautive) souvent utilisées par les modèles
\providecommand{\cross}{\textcolor{poppink}{\ensuremath{\times}}}
\providecommand{\xmark}{\cross}\providecommand{\crossmark}{\cross}\providecommand{\cmark}{\ensuremath{\checkmark}}
% Ligne de réponse / justification à compléter (inventée par les modèles)
\providecommand{\justification}[1]{\par\noindent\textit{Justification~:} \dotfill\par}
% \textipa (package tipa non chargé) : la police ne couvre pas l'API -> simple passage
\providecommand{\textipa}[1]{#1}
% Soulignement (ulem non chargé) : \uline, \ul
\providecommand{\uline}[1]{\underline{#1}}\providecommand{\ul}[1]{\underline{#1}}
% Mot ou phrase en espagnol à l'intérieur d'un paragraphe français
\newcommand{\esp}[1]{\ifmmode\text{\textit{#1}}\else\begingroup\selectlanguage{spanish}\itshape #1\endgroup\fi}
\newtcolorbox{exemplebox}[1][]{popbase,colframe=poporange,before upper={\popboxhead{Exemple}{poporange}{#1}}}
\newtcolorbox{definition}[1][]{popbase,colframe=popblue,before upper={\popboxhead{Définition}{popblue}{#1}}}
\newtcolorbox{propriete}[1][]{popbase,colframe=poppurple,before upper={\popboxhead{Propriété}{poppurple}{#1}}}
\newtcolorbox{methode}[1][]{popbase,colframe=popgold,before upper={\popboxhead{Méthode}{popgold}{#1}}}
\newtcolorbox{attention}[1][]{popbase,colframe=poppink,before upper={\popboxhead{Attention}{poppink}{#1}}}
\newtcolorbox{experience}[1][]{popbase,colframe=popblue,colback=popblueL,before upper={\popboxhead{Expérience}{popblue}{#1}}}
% « Le savais-tu ? » : carte violette pleine (contexte, culture, Cameroun)
\newtcolorbox{savaistu}[1][]{popbase,colback=poppurple,colframe=popink,
  fontupper=\popsemi\fontsize{10.4}{14.2}\selectfont\color{white},
  before upper={\poppill{Le savais-tu\,?}{white}{poppurple}\ifx\relax#1\relax\else\hspace{6pt}{\popxbold\color{white}#1}\fi\par\vspace{7pt}}}
% Exercice résolu : énoncé en haut, \tcblower puis la solution sur fond crème
\newcounter{popexo}\newcounter{popexr}
\newtcolorbox{exoresolu}[1][]{popbase,colframe=poporange,colbacklower=popcream,
  segmentation style={popink,line width=1.2pt,dashed},
  before upper={\stepcounter{popexr}\popboxhead{Exercice résolu \thepopexr}{poporange}{#1}},
  before lower={\poppill{Solution}{popink}{popcream}\par\vspace{6pt}}}
% Exercice (non résolu en place) — la correction est dans la partie Corrigés
\newtcolorbox{exercice}[1][]{popbase,colframe=popink,
  before upper={\stepcounter{popexo}\popboxhead{Exercice \thepopexo}{popink}{#1}}}
% Corrigé
\newtcolorbox{corrige}[1][]{popbase,colframe=popgreen,colback=popgreenL,
  before upper={\popboxhead{Corrigé}{popgreen}{#1}}}
% Objectifs / prérequis / bilan
\newtcolorbox{objectifs}{popbase,colframe=popink,colback=poporangeL,
  before upper={\popboxhead{Objectifs de la leçon}{popink}{}}}
\newtcolorbox{prerequis}{popbase,colframe=popink,colback=popblueL,
  before upper={\popboxhead{Prérequis}{popblue}{}}}
\newtcolorbox{bilan}{popbase,colframe=popink,colback=white,boxrule=2.8pt,
  before upper={\popboxhead{Fiche bilan}{poporange}{L'essentiel à connaître pour l'examen}}}
% Figure encadrée avec légende
\newtcolorbox{popfigure}[1][]{enhanced,breakable=false,colback=white,colframe=popline,boxrule=1.4pt,arc=8pt,
  left=10pt,right=10pt,top=10pt,bottom=8pt,before skip=10pt,after skip=12pt,halign=center,
  lower separated=false,halign lower=center,
  fontlower=\popsemi\fontsize{9}{11.5}\selectfont\color{popmuted},
  #1}
\newcommand{\legende}[1]{\par\vspace{6pt}{\popsemi\fontsize{9}{11.5}\selectfont\color{popmuted}#1}}

% ============================================================
% Signature OnBuch+
% ============================================================
\newcommand{\poplogoinline}[1]{{\poplogo\fontsize{#1}{#1}\selectfont\color{popink}OnBuch\color{poporange}+}}
\newcommand{\signatureonbuch}{%
  \begin{tcolorbox}[enhanced,colback=popink,colframe=popink,boxrule=2.2pt,arc=10pt,
      drop shadow=poporange,left=14pt,right=14pt,top=10pt,bottom=10pt,before skip=0pt,after skip=0pt]
    \tikz[baseline=-0.7ex]\node[circle,fill=popgreen,draw=popcream,line width=1.3pt,minimum size=22pt,inner sep=0]{\popcheck{white}};%
    \hspace{9pt}\begin{minipage}[c]{0.62\linewidth}
      {\popxbold\fontsize{12}{13}\selectfont\color{popcream}Signé\ \textbullet\ {\poplogo OnBuch\color{poporange}+}}\par\vspace{2pt}
      {\raggedright\popsemi\fontsize{8}{10.5}\selectfont\color{popline}\MakeUppercase{Équipe pédagogique OnBuch+}\\\MakeUppercase{Conforme au programme officiel MINESEC}\par}
    \end{minipage}\hfill
    {\poplogo\fontsize{22}{22}\selectfont\color{popcream}OnBuch\color{poporange}+}
  \end{tcolorbox}%
}

% ============================================================
% Page de garde
% #1 titre · #2 matière · #3 niveau · #4 module · #5 n° leçon · #6 durée ·
% #7 description / objectifs (texte ou itemize)
% ============================================================
\newcommand{\pagegarde}[7]{%
  \thispagestyle{empty}
  \newgeometry{a4paper,margin=1.7cm,top=1.6cm,bottom=1.4cm}%
  \begin{tikzpicture}[remember picture,overlay]
    \fill[poporange!18] ([xshift=-3.2cm,yshift=-3.6cm]current page.north east) circle (4.6cm);
    \fill[poppurple!14] ([xshift=2.4cm,yshift=7.6cm]current page.south west) circle (3.8cm);
  \end{tikzpicture}%
  {\poplogo\fontsize{30}{30}\selectfont\color{popink}OnBuch\color{poporange}+}\hfill
  \raisebox{0.6ex}{\poppill{Cours signé \textbullet\ Premium}{popink}{popcream}}\par
  \vspace{1pt}\popsquiggle{poppurple}\par
  \vspace{8pt}
  \begin{tcolorbox}[enhanced,colback=poporange,colframe=popink,boxrule=2.8pt,arc=13pt,
      drop shadow=popink,left=18pt,right=18pt,top=12pt,bottom=13pt,after skip=10pt]
    \poppill{#2 \textbullet\ #3}{popink}{popcream}\hspace{5pt}\poppill{#4}{white}{popink}\par\vspace{8pt}
    {\popdisplay\fontsize{14}{14}\selectfont\color{popink}LEÇON #5}\par\vspace{4pt}
    {\raggedright\hyphenpenalty=10000\exhyphenpenalty=10000\popdisplay\fontsize{25}{27}\selectfont\color{popcream}\MakeUppercase{#1}\par}
    \vspace{8pt}
    {\popxbold\fontsize{9.5}{9.5}\selectfont\color{popink}\MakeUppercase{Durée conseillée : #6}}
  \end{tcolorbox}
  \begin{tcolorbox}[enhanced,colback=white,colframe=popink,boxrule=2.2pt,arc=10pt,
      drop shadow=popink,left=16pt,right=16pt,top=10pt,bottom=10pt,after skip=8pt]
    \poppill{Dans cette leçon}{poppurple}{white}\par\vspace{5pt}
    \popsemi\fontsize{9.3}{12.6}\selectfont\color{popink}#7
  \end{tcolorbox}
  \noindent\begin{minipage}[t]{0.487\linewidth}
    \begin{tcolorbox}[enhanced,colback=popgreen,colframe=popink,boxrule=2.2pt,arc=10pt,
        drop shadow=popink,left=11pt,right=11pt,top=10pt,bottom=10pt,height=3.1cm,valign=center]
      \tcbox[colback=white,colframe=popink,boxrule=1.3pt,arc=5pt,boxsep=0pt,nobeforeafter,
        left=3pt,right=3pt,top=3pt,bottom=3pt,valign=center]{\qrcode[height=1.9cm]{https://play.google.com/store/apps/details?id=com.luvvix.onbuch&hl=fr}}%
      \hspace{8pt}\begin{minipage}[c]{0.5\linewidth}
        {\popdisplay\fontsize{11}{12.5}\selectfont\color{white}\MakeUppercase{Télécharge OnBuch}}\par\vspace{4pt}
        {\raggedright\popsemi\fontsize{8.4}{11}\selectfont\color{white}Gratuit \textbullet\ Google Play\\Tuteur IA, annales, concours, résultats\par}
      \end{minipage}
    \end{tcolorbox}
  \end{minipage}\hfill
  \begin{minipage}[t]{0.487\linewidth}
    \begin{tcolorbox}[enhanced,colback=poppurple,colframe=popink,boxrule=2.2pt,arc=10pt,
        drop shadow=popink,left=11pt,right=11pt,top=10pt,bottom=10pt,height=3.1cm,valign=center]
      \tikz[baseline=-0.7ex]\node[circle,fill=white,draw=popink,line width=1.3pt,minimum size=1.6cm,inner sep=0]{\popdisplay\fontsize{24}{24}\selectfont\color{poppurple}+};%
      \hspace{8pt}\begin{minipage}[c]{0.6\linewidth}
        {\popdisplay\fontsize{11}{12.5}\selectfont\color{white}\MakeUppercase{Passe à OnBuch+}}\par\vspace{4pt}
        {\raggedright\popsemi\fontsize{8.4}{11}\selectfont\color{white}Tous les cours signés, Tuteur IA illimité, fiches PDF — onglet OnBuch+ de l'app\par}
      \end{minipage}
    \end{tcolorbox}
  \end{minipage}
  \vspace{8pt}
  \popcontenu
  \vfill
  \signatureonbuch
  \restoregeometry
  \newpage
}

% Rangée « Ce cours contient » (page de garde)
\newcommand{\poptile}[3]{% #1 couleur #2 gros texte #3 légende
  \begin{tcolorbox}[enhanced,colback=#1,colframe=popink,boxrule=2pt,arc=9pt,shadow={2.5pt}{-2.5pt}{0pt}{popink},
      left=6pt,right=6pt,top=9pt,bottom=9pt,halign=center,valign=center,height=1.75cm,width=\linewidth,nobeforeafter]
    {\popdisplay\fontsize{12.5}{13}\selectfont\color{white}\MakeUppercase{#2}}\par\vspace{3pt}
    {\centering\popsemi\fontsize{7.8}{9.5}\selectfont\color{white}#3\par}
  \end{tcolorbox}}
\newcommand{\popcontenu}{%
  \par\noindent{\popxboldls\fontsize{7.5}{7.5}\selectfont\color{popmuted}CE COURS SIGNÉ CONTIENT}\par\vspace{7pt}
  \noindent\begin{minipage}[t]{0.235\linewidth}\poptile{popblue}{Cours}{complet, illustré, pas à pas}\end{minipage}\hfill
  \begin{minipage}[t]{0.235\linewidth}\poptile{poporange}{Méthodes}{et exercices résolus}\end{minipage}\hfill
  \begin{minipage}[t]{0.235\linewidth}\poptile{poppink}{Exercices}{type examen + corrigés}\end{minipage}\hfill
  \begin{minipage}[t]{0.235\linewidth}\poptile{popgreen}{Fiche bilan}{l'essentiel à retenir}\end{minipage}\par}

% Sommaire
\newtcolorbox{sommaire}{enhanced,colback=white,colframe=popink,boxrule=2.2pt,arc=10pt,
  drop shadow=popink,left=18pt,right=18pt,top=13pt,bottom=13pt,before skip=6pt,after skip=20pt,
  before upper={\poppill{Sommaire}{poppurple}{white}\par\vspace{9pt}\popsemi\fontsize{10.6}{14.5}\selectfont\color{popink}}}

% Signature de fin de cours — poussée tout en bas de la dernière page
\newcommand{\findecours}{%
  \par\vfill\nopagebreak
  \begin{center}\popsquiggle{poporange}\end{center}\vspace{4pt}
  \signatureonbuch
}
\AtBeginDocument{\def\llbracket{\mathopen{[\mkern-3.5mu[}}\def\rrbracket{\mathclose{]\mkern-3.5mu]}}} % intervalles d'entiers (glyphe ⟦ absent de la police)
% Glyphes absents de Fira Math : redéfinis à partir de glyphes disponibles
\AtBeginDocument{%
  \def\setminus{\mathbin{\backslash}}\let\smallsetminus\setminus
  \def\vdots{\mathord{\vbox{\baselineskip3.2pt\lineskiplimit0pt\kern4pt\hbox{.}\hbox{.}\hbox{.}}}}%
  \def\ddots{\mathinner{\mkern1mu\raise6.5pt\hbox{.}\mkern2mu\raise3.6pt\hbox{.}\mkern2mu\raise0.7pt\hbox{.}\mkern1mu}}%
  \def\triangle{\mathord{\scalebox{0.9}{$\symup{\Delta}$}}}%
  \def\nabla{\mathord{\raisebox{\depth}{\scalebox{1}[-1]{$\symup{\Delta}$}}}}%
  \def\checkmark{\popcheck{popdark}}%
  \def\star{\mathbin{\ast}}\def\bigstar{\mathord{\ast}}%
  \def\diamond{\mathbin{\scalebox{0.6}{$\Diamond$}}}%
  \def\bigcup{\mathop{\mathchoice{\raisebox{-0.3em}{\scalebox{1.7}{$\cup$}}}{\scalebox{1.25}{$\cup$}}{\cup}{\cup}}\displaylimits}%
  \def\bigcap{\mathop{\mathchoice{\raisebox{-0.3em}{\scalebox{1.7}{$\cap$}}}{\scalebox{1.25}{$\cap$}}{\cap}{\cap}}\displaylimits}%
  \def\lesseqqgtr{\mathrel{\lessgtr}}\def\bigoplus{\mathop{\oplus}\displaylimits}%
}
\providecommand{\ssi}{si et seulement si} % abréviation fréquente
\AtBeginDocument{\providecommand{\wideparen}{\overparen}} % arc de cercle

%% FIGFIX-BEGIN v5 — correctifs globaux des figures (docs/onbuch-generation/tools/figfix)
\usepackage{adjustbox}\usepackage{etoolbox}\tcbuselibrary{raster}
% toute figure TikZ/pgfplots plus large que la ligne est réduite à la largeur disponible
\BeforeBeginEnvironment{tikzpicture}{\begin{adjustbox}{max width=\linewidth}}
\AfterEndEnvironment{tikzpicture}{\end{adjustbox}}
% légendes pgfplots lisibles par-dessus les courbes
\pgfplotsset{every axis/.append style={legend style={fill=white,fill opacity=0.92,text opacity=1,draw=black!15,inner sep=3pt}}}
% étiquettes d'arêtes (syntaxe edge["..."]) avec fond blanc
\tikzset{every edge quotes/.append style={fill=white,inner sep=1.2pt}}
%% FIGFIX-END
\begin{document}

\pagegarde{Redacción : carta abierta / petición ; presentación oral de un país hispanohablante}{Espagnol}{Tle A}{Módulo 4 : Activités économiques et monde du travail}{E19}{2 h}{Cette leçon mixte apprend à lire et commenter un texte revendicatif, à rédiger une lettre ouverte ou une pétition en espagnol, et à présenter oralement un pays hispanophone (géographie, histoire, culture, économie). Elle combine lexique thématique, grammaire de l'argumentation et expression écrite et orale.
\vspace{4pt}
\begin{multicols}{2}\small
\begin{itemize}
\item À la fin de cette leçon, je sais identifier la structure et les procédés d'un texte revendicatif
\item je sais utiliser le lexique de la revendication, des droits et de la vie économique
\item je sais rédiger une lettre ouverte conforme aux conventions hispaniques
\item je sais rédiger une pétition avec formules d'appel et de signature
\item je sais argumenter une revendication avec connecteurs et subjonctif
\item je sais présenter oralement un pays hispanophone (géographie, histoire, culture, économie)
\end{itemize}\end{multicols}}

\begin{sommaire}
\begin{multicols}{2}
\begin{enumerate}
\item Texte support : une lettre ouverte revendicative
\item Lexique thématique : revendication, droits et monde du travail
\item Grammaire de l'argumentation : exprimer l'exigence et l'opinion
\item La lettre ouverte et la pétition : structure et formules
\item Présentation orale d'un pays hispanophone
\item Production guidée et entraînement au Bac
\item Activité d'intégration
\item Méthodes et exercices résolus
\item Exercices
\item Corrigés des exercices
\item Fiche bilan
\end{enumerate}
\end{multicols}
\end{sommaire}

\begin{objectifs}
\begin{itemize}
\item À la fin de cette leçon, je sais identifier la structure et les procédés d'un texte revendicatif
\item je sais utiliser le lexique de la revendication, des droits et de la vie économique
\item je sais rédiger une lettre ouverte conforme aux conventions hispaniques
\item je sais rédiger une pétition avec formules d'appel et de signature
\item je sais argumenter une revendication avec connecteurs et subjonctif
\item je sais présenter oralement un pays hispanophone (géographie, histoire, culture, économie)
\item je sais employer les tournures de l'opinion et de l'exigence (es necesario que + subjuntivo...)
\item je sais éviter les pièges francophones dans les formules de lettre formelle
\end{itemize}
\end{objectifs}

\begin{prerequis}
\begin{itemize}
\item Le subjonctif présent après les verbes d'opinion et de volonté (Première)
\item Les connecteurs logiques de base (pero, además, por eso, sin embargo)
\item La structure générale d'une lettre formelle (Première)
\item Les gentilés et noms de pays hispanophones
\item L'imparfait et le passé simple pour raconter un contexte historique
\end{itemize}
\end{prerequis}

\newpage

\coursec{Texte support : une lettre ouverte revendicative}

\courssub{Situation d'accroche et problématique}

Imaginez la scène : à Yaoundé, les élèves d'un lycée découvrent que leur bibliothèque, fermée depuis des mois, n'a plus ni livres récents ni tables en bon état. Plutôt que de se plaindre dans les couloirs, ils décident de rédiger une \cle{pétition} adressée au proviseur et publiée dans le journal du lycée. Mais comment formuler une revendication claire, respectueuse et convaincante en espagnol ? De l'autre côté de l'Atlantique, des jeunes mexicains ou chiliens signent chaque jour des \esp{cartas abiertas} publiées dans la presse pour défendre leurs droits. Et si, pour convaincre un jury ou un public, vous deviez aussi présenter en quelques minutes les richesses du Mexique ou du Pérou ? Cette leçon vous donne les outils linguistiques et méthodologiques pour écrire, revendiquer et présenter.

\textbf{Problématique :} comment une lettre ouverte réussit-elle à être à la fois ferme et respectueuse, et quels procédés linguistiques rendent une revendication efficace ?

Pour répondre à cette question, nous allons lire un texte authentique de type revendicatif, l'observer méthodiquement, dégager son lexique et ses procédés, puis nous entraîner à le comprendre en profondeur.

\courssub{Le genre de la lettre ouverte}

Avant de lire le texte, définissons précisément le genre auquel il appartient.

\begin{definition}[La carta abierta]
La \cle{carta abierta} (lettre ouverte) est une lettre adressée officiellement à une autorité (président, ministre, maire, recteur...) mais \textbf{publiée} dans la presse ou sur les réseaux sociaux afin d'informer et de mobiliser l'opinion publique. Elle possède donc un \emph{double destinataire} : l'autorité visée, qui doit répondre, et les lecteurs, qu'elle cherche à convaincre et à rallier à la cause.
\end{definition}

\begin{savaistu}[Une tradition latino-américaine]
Les \esp{cartas abiertas} sont une tradition journalistique et citoyenne très forte en Amérique latine. Des écrivains célèbres comme Gabriel García Márquez en ont publié dans la presse colombienne pour interpeller les gouvernants sur des questions de société. En Guinée équatoriale, seul pays hispanophone d'Afrique, les élèves et les associations utilisent aussi ce genre pour s'adresser aux autorités éducatives. Maîtriser la lettre ouverte, c'est donc entrer dans une pratique réelle du monde hispanophone.
\end{savaistu}

\courssub{Lecture du texte}

\begin{methode}[Avant de lire : la stratégie du premier repérage]
Ne plongez jamais dans un texte argumentatif sans repères. Procédez en trois étapes :
\begin{enumerate}
\item Lisez d'abord \textbf{l'en-tête et la signature} pour identifier l'émetteur (qui écrit ?) et le destinataire (à qui ?).
\item Repérez ensuite la \textbf{revendication centrale}, souvent introduite par un verbe comme \esp{exigir}, \esp{pedir} ou \esp{reivindicar}.
\item Cherchez enfin les \textbf{arguments} qui justifient la demande, signalés par des connecteurs (\esp{además}, \esp{por eso}...).
\end{enumerate}
Cette démarche, exigée au baccalauréat, vous fera gagner un temps précieux.
\end{methode}

\begin{textebox}[Carta abierta a las autoridades : por una educación digna]
Señoras y señores representantes del gobierno :

Nos dirigimos a ustedes, los jóvenes estudiantes de esta región, para exigir que se respete nuestro derecho a una educación digna. Nuestras escuelas carecen de libros, de laboratorios y de profesores suficientes. Es urgente que el Estado aumente el presupuesto destinado a la educación pública. Además, pedimos que se construyan nuevas aulas y que se reparen las que están en ruinas. No se trata de un capricho : se trata de nuestro futuro y del futuro del país. Por eso, hacemos un llamamiento a la ciudadanía para que apoye esta petición con su firma. Esperamos una respuesta concreta y rápida.

Atentamente,

Los estudiantes de la región.
\end{textebox}

\textbf{Traduction :} « Mesdames et Messieurs les représentants du gouvernement : Nous nous adressons à vous, les jeunes étudiants de cette région, pour exiger que soit respecté notre droit à une éducation digne. Nos écoles manquent de livres, de laboratoires et de professeurs en nombre suffisant. Il est urgent que l'État augmente le budget destiné à l'éducation publique. De plus, nous demandons que de nouvelles salles de classe soient construites et que soient réparées celles qui sont en ruines. Il ne s'agit pas d'un caprice : il s'agit de notre avenir et de l'avenir du pays. C'est pourquoi nous lançons un appel à la citoyenneté pour qu'elle soutienne cette pétition par sa signature. Nous espérons une réponse concrète et rapide. Cordialement, les étudiants de la région. »

\courssub{Glossaire : le vocabulaire de la revendication}

\begin{center}
\begin{tabularx}{\linewidth}{|X|X|X|}
\hline
\rowcolor{popblueL} \textbf{Español} & \textbf{Français} & \textbf{Exemple en contexte} \\
\hline
\esp{exigir} & exiger & \esp{Exigimos que se respete el derecho.} \\
\hline
\esp{reivindicar} & revendiquer & \esp{Los sindicatos reivindican mejores salarios.} \\
\hline
\esp{la huelga} & la grève & \esp{Los profesores están en huelga.} \\
\hline
\esp{el derecho} & le droit & \esp{La educación es un derecho.} \\
\hline
\esp{la desigualdad} & l'inégalité & \esp{Luchamos contra la desigualdad.} \\
\hline
\esp{los recursos} & les ressources & \esp{Faltan recursos en las escuelas.} \\
\hline
\esp{digno, digna} & digne & \esp{Una educación digna para todos.} \\
\hline
\esp{el presupuesto} & le budget & \esp{El presupuesto de educación es bajo.} \\
\hline
\esp{comprometerse} & s'engager & \esp{El gobierno debe comprometerse.} \\
\hline
\esp{la ciudadanía} & la citoyenneté & \esp{La ciudadanía apoya la petición.} \\
\hline
\end{tabularx}
\end{center}

\begin{attention}[Faux amis et pièges lexicaux]
\begin{itemize}
\item \esp{el derecho} signifie « le droit » (juridique) ; attention à ne pas confondre avec l'adverbe \esp{derecho} (« tout droit ») ou l'adjectif \esp{a la derecha} (« à droite »).
\item \esp{exigir} se conjugue avec un changement \textbf{g $\rightarrow$ j} devant \textbf{a} et \textbf{o} : \esp{yo exijo}, \esp{que yo exija}. Ne dites jamais \esp{exigo}.
\item \esp{comprometerse} est pronominal : on dit \esp{me comprometo}, « je m'engage ». Sans pronom, \esp{comprometer} signifie « compromettre ».
\item « Une pétition » se dit \esp{una petición} ou \esp{una petitoria} ; \esp{una petición} désigne aussi toute demande.
\end{itemize}
\end{attention}

\courssub{Analyse du texte : émetteur, destinataire, revendication}

Observons maintenant le texte de près, comme au commentaire de texte du baccalauréat.

\textbf{1. L'émetteur.} La signature (\esp{Los estudiantes de la región}) et la formule \esp{Nos dirigimos a ustedes} indiquent un émetteur \textbf{collectif} : les jeunes étudiants. Le « nous » (\esp{nos}, \esp{nuestro derecho}, \esp{nuestras escuelas}) donne du poids à la parole : ce n'est pas une plainte individuelle mais la voix d'un groupe.

\textbf{2. Le destinataire.} L'en-tête \esp{Señoras y señores representantes del gobierno} désigne les autorités. Mais comme la lettre est \emph{ouverte}, le vrai destinataire est double : le gouvernement, qui doit agir, et l'opinion publique, appelée à signer (\esp{hacemos un llamamiento a la ciudadanía}).

\textbf{3. La revendication centrale.} Elle apparaît dès la première phrase : \esp{exigir que se respete nuestro derecho a una educación digna}. Le verbe \esp{exigir} marque la fermeté ; le complément \esp{que se respete...} précise l'objet de l'exigence.

\textbf{4. Les arguments.} On en compte trois, organisés en progression :
\begin{itemize}
\item \textbf{Argument 1 (constat) :} \esp{Nuestras escuelas carecen de libros, de laboratorios y de profesores suficientes} — la description d'une carence concrète.
\item \textbf{Argument 2 (urgence budgétaire) :} \esp{Es urgente que el Estado aumente el presupuesto} — la solution passe par une décision politique.
\item \textbf{Argument 3 (enjeu collectif) :} \esp{se trata de nuestro futuro y del futuro del país} — l'élargissement : l'école concerne toute la nation.
\end{itemize}

\textbf{5. L'appel final.} La lettre se termine par un appel à l'action (\esp{hacemos un llamamiento... para que apoye esta petición}) et une attente ferme (\esp{Esperamos una respuesta concreta y rápida}).

\begin{popfigure}
\begin{tikzpicture}[
mindmap,
grow cyclic,
every node/.style={concept, text=white, align=center, font=\small},
concept color=popblue,
level 1/.append style={level distance=3.4cm, sibling angle=72},
level 2/.append style={level distance=2.6cm, sibling angle=50}
]
\node[concept color=popdark, font=\small\bfseries] {Carta abierta}
child[concept color=poporange] { node[align=center]{Emisor\\ \esp{los estudiantes}} }
child[concept color=poppurple] { node[align=center]{Destinatario\\ \esp{el gobierno}\\ + opinión pública} }
child[concept color=popgreen] { node[align=center]{Reivindicación\\ \esp{educación digna}} }
child[concept color=poppink] { node {Argumentos}
  child[concept color=poppink!70] { node[font=\scriptsize, align=center]{falta de\\ recursos} }
  child[concept color=poppink!70] { node[font=\scriptsize, align=center]{aumentar el\\ presupuesto} }
  child[concept color=poppink!70] { node[font=\scriptsize, align=center]{futuro\\ del país} }
}
child[concept color=popgold] { node[align=center]{Llamamiento\\ \esp{firma y apoyo}} };
\end{tikzpicture}
\legende{Carte mentale du texte : de l'émetteur à l'appel final.}
\end{popfigure}

\courssub{Le ton et le registre}

\begin{propriete}[Ton solennel, ferme mais respectueux]
Une lettre ouverte efficace évite l'insulte : elle combine \textbf{fermeté} dans le fond et \textbf{respect} dans la forme. Trois indices le montrent ici :
\begin{itemize}
\item Le \textbf{vouvoiement soutenu} : \esp{Nos dirigimos a ustedes} (et non \esp{vosotros}) avec la majuscule de politesse implicite dans la formule d'appel.
\item Les \textbf{formules protocolaires} : \esp{Señoras y señores representantes}, \esp{Atentamente}.
\item L'absence d'agressivité personnelle : on critique une \emph{situation} (\esp{las escuelas carecen de...}), jamais une personne.
\end{itemize}
Le registre est \cle{soutenu} : vocabulaire abstrait (\esp{derecho}, \esp{presupuesto}, \esp{ciudadanía}), phrases construites, subjonctifs corrects.
\end{propriete}

\begin{aretenir}[Les procédés linguistiques de la revendication]
Trois procédés dominent dans ce texte :
\begin{enumerate}
\item \textbf{Le subjonctif d'exigence} : après \esp{exigir que}, \esp{pedir que}, \esp{es urgente que}, le subordonné est au subjonctif : \esp{exigimos que se \textbf{respete}}, \esp{pedimos que se \textbf{construyan}}. (Cette règle sera étudiée en détail dans la section grammaire.)
\item \textbf{Les connecteurs d'organisation} : \esp{además} (ajout), \esp{por eso} (conséquence), \esp{no se trata de... se trata de...} (correction et recentrage).
\item \textbf{Le champ lexical du droit} : \esp{derecho}, \esp{educación digna}, \esp{ciudadanía}, \esp{petición}, \esp{firma} — un vocabulaire qui inscrit la demande dans le cadre légitime de la loi, et non dans le caprice.
\end{enumerate}
\end{aretenir}

\courssub{Questions de compréhension}

\textbf{Compréhension globale}
\begin{enumerate}
\item ¿Quiénes escriben esta carta y a quiénes se dirigen ?
\item ¿Cuál es la reivindicación principal del texto ?
\item ¿Qué tono general tiene la carta : agresivo, irónico o firme y respetuoso ?
\end{enumerate}

\textbf{Compréhension détaillée}
\begin{enumerate}
\item ¿De qué carecen las escuelas de la región ? Cita tres elementos.
\item ¿Qué debe hacer el Estado según los estudiantes ?
\item ¿Qué piden los estudiantes además del aumento del presupuesto ?
\item Explica la frase : \esp{No se trata de un capricho : se trata de nuestro futuro}.
\item ¿Qué esperan los estudiantes de la ciudadanía y del gobierno al final de la carta ?
\end{enumerate}

\begin{exoresolu}[Repérage des procédés argumentatifs]
Énoncé : relevez dans la lettre (a) un verbe de revendication suivi du subjonctif, (b) un connecteur d'addition, (c) un élément du champ lexical du droit, et justifiez leur effet.
\tcblower
Solution détaillée :
\begin{itemize}
\item (a) \esp{pedimos que se \textbf{construyan} nuevas aulas} : le verbe \esp{pedir} exprime une exigence ; le subjonctif \esp{construyan} marque que l'action n'est pas encore réalisée — c'est une demande, pas un fait. Effet : la revendication reste ferme mais grammaticalement correcte et respectueuse.
\item (b) \esp{\textbf{Además}, pedimos que...} : ce connecteur d'addition organise les arguments en escalier et donne au texte une structure claire, facile à suivre pour le lecteur.
\item (c) \esp{nuestro \textbf{derecho} a una educación \textbf{digna}} : en parlant de « droit » et non de « souhait », les étudiants transforment leur demande en exigence légitime que l'autorité ne peut ignorer sans se discréditer.
\end{itemize}
\end{exoresolu}

\begin{exercice}[À votre tour]
Répondez en espagnol par des phrases complètes :
\begin{enumerate}
\item Subraya en el texto todos los verbos seguidos de subjuntivo y explica por qué se usa este modo.
\item ¿Qué diferencia hay entre \esp{exigir} y \esp{pedir} ? ¿Cuál de los dos verbos es más fuerte ?
\item Imagina que eres alumno de un liceo de Douala : escribe dos frases de revendication con \esp{exigir que} et \esp{pedir que} au subjonctif.
\end{enumerate}
\end{exercice}

\begin{corrige}[Exercice : À votre tour]
\begin{enumerate}
\item Les verbes suivis du subjonctif : \esp{exigir que se respete}, \esp{es urgente que el Estado aumente}, \esp{pedimos que se construyan}, \esp{que se reparen}, \esp{para que apoye}. Le subjonctif s'emploie car ces verbes expriment une volonté, une exigence ou un but portant sur des actions non réalisées.
\item \esp{Pedir} signifie « demander » (neutre, poli) ; \esp{exigir} signifie « exiger », c'est-à-dire demander avec autorité ce qu'on considère comme un droit. \esp{Exigir} est donc le verbe le plus fort.
\item Exemples attendus : \esp{Exigimos que se repare la biblioteca del liceo.} « Nous exigeons que la bibliothèque du lycée soit réparée. » ; \esp{Pedimos que haya más libros en español.} « Nous demandons qu'il y ait plus de livres en espagnol. »
\end{enumerate}
\end{corrige}

\begin{aretenir}[L'essentiel de la section]
\begin{itemize}
\item La \esp{carta abierta} vise un double destinataire : l'autorité et l'opinion publique.
\item Sa structure type : en-tête solennel $\rightarrow$ présentation de l'émetteur $\rightarrow$ revendication centrale $\rightarrow$ arguments $\rightarrow$ appel final $\rightarrow$ formule de politesse.
\item Son efficacité repose sur un ton ferme mais respectueux, le subjonctif d'exigence, des connecteurs clairs et un champ lexical du droit.
\item Stratégie de lecture : en-tête et signature d'abord, revendication ensuite, arguments enfin.
\end{itemize}
\end{aretenir}

\coursec{Lexique thématique : revendication, droits et monde du travail}

Dans la section précédente, nous avons lu et commenté une lettre ouverte dans laquelle des travailleurs réclamaient de meilleures conditions de travail. Ce texte regorgeait de mots comme \esp{exigimos}, \esp{derechos}, \esp{huelga} ou \esp{salario}. Pour pouvoir, à votre tour, rédiger une lettre ouverte ou une pétition — et pour présenter un pays hispanophone à l'oral — il faut d'abord maîtriser ce vocabulaire. C'est l'objet de cette section : nous allons construire, champ par champ, le lexique de la revendication, des droits, du monde du travail et de la présentation d'un pays.

\courssub{Observer avant d'apprendre : le lexique en contexte}

Avant tout tableau, lisons ce court extrait de pétition et repérons les mots importants.

\begin{textebox}[Extrait d'une pétition (texte original)]
Los trabajadores de la empresa exigimos un aumento de salario y denunciamos las malas condiciones de trabajo. Tenemos derecho a un contrato justo y a la libertad de expresión. Por eso hacemos un llamamiento a las autoridades y pedimos a todos los ciudadanos que firmen esta petición. Queremos poner fin a la desigualdad y dar respuesta a las necesidades de las familias. Si la empresa no responde, estaremos en huelga la próxima semana.
\end{textebox}

\textbf{Glossaire :} \esp{el aumento} = l'augmentation ; \esp{las condiciones} = les conditions ; \esp{justo} = juste, équitable ; \esp{las autoridades} = les autorités ; \esp{ciudadanos} = citoyens ; \esp{las necesidades} = les besoins ; \esp{si... no responde} = si... ne répond pas.

Observez : chaque verbe exprime une action précise de la revendication (\esp{exigir}, \esp{denunciar}, \esp{pedir}, \esp{hacer un llamamiento}), et chaque nom désigne une réalité du monde du travail ou des droits (\esp{salario}, \esp{contrato}, \esp{huelga}, \esp{desigualdad}). Organisons maintenant ce vocabulaire en quatre champs lexicaux.

\courssub{Champ 1 — La revendication}

\begin{definition}[Les verbes de la revendication]
Revendiquer, c'est demander avec force quelque chose que l'on considère comme un droit. L'espagnol dispose de plusieurs verbes, du plus neutre au plus combatif :
\begin{itemize}
\item \esp{pedir} = demander (neutre, courant) : \esp{Pedimos más información.} « Nous demandons plus d'informations. »
\item \esp{solicitar} = solliciter, demander officiellement (soutenu) : \esp{Solicitan una entrevista con el ministro.} « Ils sollicitent un entretien avec le ministre. »
\item \esp{reclamar} = réclamer (insister sur ce qui nous est dû) : \esp{Los vecinos reclaman agua potable.} « Les habitants réclament de l'eau potable. »
\item \esp{exigir} = exiger (le plus fort) : \esp{Exigimos respuestas.} « Nous exigeons des réponses. »
\item \esp{reivindicar} = revendiquer (défendre une cause, un droit) : \esp{El sindicato reivindica mejores salarios.} « Le syndicat revendique de meilleurs salaires. »
\item \esp{protestar} = protester : \esp{Los estudiantes protestan contra el aumento de las tasas.} « Les étudiants protestent contre l'augmentation des frais. »
\item \esp{denunciar} = dénoncer : \esp{Denuncian la corrupción.} « Ils dénoncent la corruption. »
\end{itemize}
\end{definition}

\begin{propriete}[Constructions à connaître]
\begin{itemize}
\item \esp{exigir / pedir / reclamar + que + subjonctif} : \esp{Exigimos que se respeten nuestros derechos.} « Nous exigeons que nos droits soient respectés. » (Nous reverrons ce subjonctif dans la section 3.)
\item \esp{protestar contra / por + nom} : \esp{protestar contra la injusticia} « protester contre l'injustice ».
\item \esp{denunciar + nom} (sans préposition) : \esp{denunciar un abuso} « dénoncer un abus ».
\item Expressions nominales : \esp{hacer un llamamiento (a)} = lancer un appel (à) ; \esp{firmar una petición} = signer une pétition ; \esp{una manifestación} = une manifestation ; \esp{un manifiesto} = un manifeste.
\end{itemize}
\end{propriete}

\begin{definition}[La \cle{petición}]
Une \esp{petición} est un texte signé collectivement pour demander une mesure concrète à une autorité (maire, ministre, directeur d'entreprise). Elle contient toujours : l'objet de la demande, les arguments, la mesure attendue et la liste des signatures. Exemple : \esp{Los abajo firmantes pedimos la construcción de un pozo de agua en el barrio.} « Les soussignés demandent la construction d'un puits d'eau dans le quartier. »
\end{definition}

\courssub{Champ 2 — Les droits et la justice}

\begin{definition}[Vocabulaire des \cle{derechos}]
\begin{itemize}
\item \esp{el derecho} = le droit ; \esp{los derechos humanos} = les droits de l'homme ; \esp{el derecho al trabajo / a la educación / a la salud} = le droit au travail / à l'éducation / à la santé.
\item \esp{la justicia} = la justice ; \esp{la injusticia} = l'injustice.
\item \esp{la igualdad} = l'égalité ; \esp{la desigualdad} = l'inégalité : \esp{la igualdad entre hombres y mujeres} « l'égalité entre les hommes et les femmes ».
\item \esp{la libertad de expresión} = la liberté d'expression ; \esp{la libertad de prensa} = la liberté de la presse.
\item \esp{la dignidad} = la dignité : \esp{una vida digna} « une vie digne » ; \esp{un salario digno} « un salaire décent ».
\end{itemize}
\end{definition}

\begin{exemplebox}[Exemples traduits]
\esp{Exigimos que se respeten nuestros derechos.} « Nous exigeons que nos droits soient respectés. »\par
\esp{Todos los ciudadanos tienen derecho a la educación.} « Tous les citoyens ont droit à l'éducation. »\par
\esp{La libertad de expresión es un derecho fundamental.} « La liberté d'expression est un droit fondamental. »\par
\esp{Luchamos por la igualdad y contra la desigualdad.} « Nous luttons pour l'égalité et contre l'inégalité. »
\end{exemplebox}

\courssub{Champ 3 — Le travail et l'économie}

\begin{definition}[Vocabulaire du \cle{monde du travail}]
\begin{itemize}
\item \esp{el empleo} = l'emploi ; \esp{el desempleo / el paro} = le chômage ; \esp{un desempleado / un parado} = un chômeur.
\item \esp{el trabajador / la trabajadora} = le/la travailleur(se) ; \esp{el obrero} = l'ouvrier ; \esp{el empleado} = l'employé.
\item \esp{el salario / el sueldo} = le salaire ; \esp{el salario mínimo} = le salaire minimum.
\item \esp{la empresa} = l'entreprise ; \esp{el empresario} = le chef d'entreprise ; \esp{el jefe} = le chef, le patron.
\item \esp{el sindicato} = le syndicat ; \esp{el sindicalista} = le syndicaliste.
\item \esp{la huelga} = la grève ; \esp{el contrato} = le contrat ; \esp{el presupuesto} = le budget.
\end{itemize}
\end{definition}

\begin{exemplebox}[Exemples traduits]
\esp{Los trabajadores están en huelga.} « Les travailleurs sont en grève. »\par
\esp{El sindicato negocia un nuevo contrato con la empresa.} « Le syndicat négocie un nouveau contrat avec l'entreprise. »\par
\esp{El desempleo juvenil es un problema grave.} « Le chômage des jeunes est un problème grave. »\par
\esp{El gobierno aumentará el presupuesto de la educación.} « Le gouvernement augmentera le budget de l'éducation. »
\end{exemplebox}

\courssub{Champ 4 — La présentation d'un pays}

Pour la présentation orale d'un pays hispanophone, il faut un lexique descriptif précis.

\begin{definition}[Vocabulaire de la \cle{présentation d'un pays}]
\begin{itemize}
\item \esp{la capital} = la capitale ; \esp{la población} = la population ; \esp{los habitantes} = les habitants.
\item \esp{la superficie} = la superficie : \esp{México tiene una superficie de casi dos millones de kilómetros cuadrados.} « Le Mexique a une superficie de près de deux millions de kilomètres carrés. »
\item \esp{la moneda} = la monnaie : \esp{La moneda de Argentina es el peso.} « La monnaie de l'Argentine est le peso. »
\item \esp{el idioma oficial} = la langue officielle ; \esp{las lenguas indígenas} = les langues indigènes.
\item \esp{las exportaciones} = les exportations ; \esp{exportar petróleo, café, cacao} = exporter du pétrole, du café, du cacao.
\item \esp{el turismo} = le tourisme ; \esp{los recursos naturales} = les ressources naturelles.
\item \esp{la independencia} = l'indépendance : \esp{Guinea Ecuatorial obtuvo su independencia en 1968.} « La Guinée équatoriale a obtenu son indépendance en 1968. »
\end{itemize}
\end{definition}

\begin{savaistu}[Le Mexique, géant hispanophone ; la Guinée équatoriale, voisine africaine]
Le Mexique est le premier pays hispanophone par sa population : plus de 130 millions d'habitants (2024), devant la Colombie, l'Argentine et l'Espagne. Sa capitale, \esp{Ciudad de México}, est l'une des plus grandes villes du monde. À l'inverse, la Guinée équatoriale, seul pays d'Afrique subsaharienne de langue officielle espagnole, compte environ 1,7 million d'habitants — un sujet de présentation orale idéal pour un élève camerounais, par sa proximité géographique et culturelle !
\end{savaistu}

\courssub{Tableau récapitulatif des quatre champs lexicaux}

\begin{center}
\begin{tabularx}{\linewidth}{|l|X|X|}
\hline
\rowcolor{popblueL} \textbf{Champ} & \textbf{Español} & \textbf{Français} \\
\hline
Revendication & exigir, reclamar, reivindicar, protestar, denunciar, hacer un llamamiento, firmar una petición & exiger, réclamer, revendiquer, protester, dénoncer, lancer un appel, signer une pétition \\
\hline
Droits et justice & el derecho, los derechos humanos, la justicia, la igualdad, la desigualdad, la libertad de expresión, la dignidad & le droit, les droits de l'homme, la justice, l'égalité, l'inégalité, la liberté d'expression, la dignité \\
\hline
Travail et économie & el empleo, el desempleo, el salario, el trabajador, la empresa, el sindicato, la huelga, el contrato, el presupuesto & l'emploi, le chômage, le salaire, le travailleur, l'entreprise, le syndicat, la grève, le contrat, le budget \\
\hline
Présentation d'un pays & la capital, la población, la superficie, la moneda, el idioma oficial, las exportaciones, el turismo, la independencia & la capitale, la population, la superficie, la monnaie, la langue officielle, les exportations, le tourisme, l'indépendance \\
\hline
\end{tabularx}
\end{center}

\courssub{Carte mentale du champ lexical}

\begin{popfigure}
\begin{center}
\begin{tikzpicture}[mindmap, grow cyclic, every node/.style={concept}, concept color=popblue!60,
level 1/.append style={level distance=3.4cm, sibling angle=90},
level 2/.append style={level distance=2.1cm, sibling angle=45}]
\node[concept color=poporange!70, font=\bfseries] {la reivindicación}
child[concept color=popblue!50] { node {revendiquer}
  child { node[concept color=popblue!25] {\esp{exigir}} }
  child { node[concept color=popblue!25] {\esp{denunciar}} }
  child { node[concept color=popblue!25] {\esp{firmar una petición}} }
}
child[concept color=popgreen!50] { node {droits}
  child { node[concept color=popgreen!25] {\esp{los derechos humanos}} }
  child { node[concept color=popgreen!25] {\esp{la igualdad}} }
  child { node[concept color=popgreen!25] {\esp{la dignidad}} }
}
child[concept color=popgold!60] { node {travail}
  child { node[concept color=popgold!30] {\esp{el salario}} }
  child { node[concept color=popgold!30] {\esp{la huelga}} }
  child { node[concept color=popgold!30] {\esp{el sindicato}} }
}
child[concept color=poppurple!45] { node {pays}
  child { node[concept color=poppurple!20] {\esp{la capital}} }
  child { node[concept color=poppurple!20] {\esp{la moneda}} }
  child { node[concept color=poppurple!20] {\esp{la independencia}} }
};
\end{tikzpicture}
\end{center}
\legende{Carte mentale du lexique : quatre branches autour de la revendication.}
\end{popfigure}

\courssub{Expressions figées à mémoriser}

\begin{aretenir}[Cinq expressions indispensables]
\begin{itemize}
\item \esp{luchar por + nom/infinitif} = lutter pour : \esp{Luchamos por nuestros derechos.} « Nous luttons pour nos droits. »
\item \esp{tener derecho a + nom/infinitif} = avoir le droit de : \esp{Tenemos derecho a protestar.} « Nous avons le droit de protester. »
\item \esp{estar en huelga} = être en grève : \esp{Los obreros están en huelga desde el lunes.} « Les ouvriers sont en grève depuis lundi. »
\item \esp{poner fin a + nom} = mettre fin à : \esp{Hay que poner fin a la injusticia.} « Il faut mettre fin à l'injustice. »
\item \esp{dar respuesta a + nom} = répondre à : \esp{El gobierno debe dar respuesta a las demandas.} « Le gouvernement doit répondre aux revendications. »
\end{itemize}
\end{aretenir}

\courssub{Faux amis et pièges pour les francophones}

\begin{attention}[« demandar » ne veut pas dire « demander » !]
\esp{demandar} signifie « réclamer, exiger » (sens fort) ou « poursuivre en justice ». Pour « demander » poliment, on dit \esp{pedir} ou \esp{solicitar}.\par
\esp{Los sindicatos demandan mejores condiciones.} « Les syndicats réclament de meilleures conditions. »\par
\esp{Pido un aumento de salario.} « Je demande une augmentation de salaire. »\par
De même, \esp{actual} signifie « actuel » (du moment présent), jamais « réel » : \esp{el presidente actual} = « le président actuel ». Pour « réel », dites \esp{real} ou \esp{verdadero}. Et attention : \esp{demanda} (nom) = « revendication » ou « procès », pas « demande » au sens courant.
\end{attention}

\begin{attention}[Autres pièges fréquents]
\begin{itemize}
\item \esp{la empresa} = l'entreprise (pas « l'emprise ») ; \esp{el empresario} = le chef d'entreprise.
\item \esp{el sueldo} et \esp{el salario} sont synonymes ; ne dites pas \esp{el salario mínimo vital y móvil} hors d'Argentine, préférez simplement \esp{el salario mínimo}.
\item On dit \esp{estar en huelga} (avec \esp{estar}), jamais \esp{ser en huelga}.
\item \esp{la moneda} = la monnaie (système) ; « une pièce de monnaie » = \esp{una moneda} aussi, mais « de la monnaie (rendue) » = \esp{el cambio} ou \esp{la vuelta}.
\end{itemize}
\end{attention}

\courssub{Exercice résolu}

\begin{exoresolu}[Complétez avec le mot juste]
Complétez les phrases avec : \esp{exigir}, \esp{huelga}, \esp{derechos humanos}, \esp{salario}, \esp{petición}, \esp{desempleo}.\par
1. Los ciudadanos firman una \ldots{} para pedir un hospital en el barrio.\par
2. Los trabajadores \ldots{} un aumento de sueldo.\par
3. Si la empresa no negocia, los obreros estarán en \ldots{}.\par
4. La libertad de expresión forma parte de los \ldots{}.\par
5. Un \ldots{} digno permite vivir bien.\par
6. El \ldots{} juvenil preocupa a los gobiernos africanos.
\tcblower
1. \esp{petición} — on signe une pétition (\esp{firmar una petición}).\par
2. \esp{exigen} — verbe \esp{exigir}, 3\textsuperscript{e} personne du pluriel ; sens : « réclament avec force ».\par
3. \esp{huelga} — expression figée \esp{estar en huelga} = être en grève.\par
4. \esp{derechos humanos} — la liberté d'expression est un droit de l'homme.\par
5. \esp{salario} — \esp{un salario digno} = un salaire décent.\par
6. \esp{desempleo} — \esp{el desempleo juvenil} = le chômage des jeunes.
\end{exoresolu}

\begin{exercice}[À vous de jouer]
Traduisez en espagnol : 1. Les travailleurs réclament un meilleur salaire. 2. Nous luttons pour l'égalité. 3. Les citoyens ont signé une pétition contre l'injustice. 4. La capitale du Mexique compte des millions d'habitants. 5. Il faut mettre fin au chômage des jeunes.
\end{exercice}

\begin{corrige}[Exercice « À vous de jouer »]
1. \esp{Los trabajadores reclaman un mejor salario.} 2. \esp{Luchamos por la igualdad.} 3. \esp{Los ciudadanos han firmado una petición contra la injusticia.} 4. \esp{La capital de México cuenta con millones de habitantes.} 5. \esp{Hay que poner fin al desempleo juvenil.} (Attention : \esp{poner fin a + el} donne la contraction \esp{al}.)
\end{corrige}

Ce lexique est désormais votre boîte à outils. Dans la section suivante, nous verrons comment la grammaire — en particulier le subjonctif après les verbes d'exigence et d'opinion — permet de donner à ces mots toute leur force argumentative.

\coursec{Grammaire de l'argumentation : exprimer l'exigence et l'opinion}

Dans la section précédente, nous avons constitué le lexique de la revendication et du monde du travail : \esp{los derechos}, \esp{la huelga}, \esp{el salario mínimo}, \esp{reivindicar}... Mais un vocabulaire riche ne suffit pas pour écrire une lettre ouverte efficace. Encore faut-il savoir \emph{construire} la phrase revendicative : comment dire « nous exigeons que... », « il est urgent que... », « nous ne croyons pas que... » ? C'est précisément l'objet de cette section : la grammaire de l'argumentation, c'est-à-dire l'articulation entre les verbes de volonté, les tournures impersonnelles et le \cle{subjonctif}, mode roi de la revendication en espagnol.

\courssub{Observation : que remarque-t-on dans le texte ?}

Reprenons trois phrases tirées de la lettre ouverte étudiée dans la section 1 :

\begin{exemplebox}[Phrases à observer]
\begin{itemize}
\item \esp{Exigimos que se \textbf{respete} el salario mínimo.} « Nous exigeons que le salaire minimum \textbf{soit} respecté. »
\item \esp{Es urgente que el Estado \textbf{aumente} el presupuesto de educación.} « Il est urgent que l'État \textbf{augmente} le budget de l'éducation. »
\item \esp{Pedimos que se \textbf{construyan} más escuelas en los barrios populares.} « Nous demandons que l'on \textbf{construise} davantage d'écoles dans les quartiers populaires. »
\end{itemize}
\end{exemplebox}

Observons attentivement. Après \esp{exigimos que}, \esp{es urgente que}, \esp{pedimos que}, on attendrait peut-être l'indicatif (\esp{respeta}, \esp{aumenta}, \esp{construyen}). Or les verbes sont \esp{respete}, \esp{aumente}, \esp{construyan} : c'est le \cle{subjonctif présent}. Remarquons aussi que le français fait exactement la même chose : « que le salaire \emph{soit} respecté », « qu'il \emph{augmente} ». Bonne nouvelle : pour une fois, les deux langues fonctionnent de la même manière. Il faut seulement connaître les déclencheurs et les formes.

\courssub{Les verbes de volonté et d'exigence + que + subjonctif}

\begin{propriete}[Exigir, pedir, querer, desear + que + subjonctif]
Après un verbe exprimant la \cle{volonté}, l'\cle{exigence}, le \cle{souhait} ou l'\cle{interdiction}, suivi de \esp{que}, le verbe de la subordonnée se met au \textbf{subjonctif} (obligatoirement, sans exception) :
\begin{center}
\esp{Exigir / pedir / querer / desear / prohibir / permitir} + \esp{que} + \textbf{SUBJONCTIF}
\end{center}
\esp{Pedimos que se reparen las aulas.} « Nous demandons que les salles soient réparées. »
\end{propriete}

\textbf{Emplois et exemples :}

\begin{itemize}
\item \esp{Exigimos que se respeten los derechos de los trabajadores.} « Nous exigeons que soient respectés les droits des travailleurs. »
\item \esp{Quiero que me escuches.} « Je veux que tu m'écoutes. »
\item \esp{Deseamos que el diálogo continúe.} « Nous souhaitons que le dialogue continue. »
\item \esp{El sindicato pide que se negocie un nuevo convenio.} « Le syndicat demande que l'on négocie une nouvelle convention. »
\item \esp{La ley prohíbe que los menores trabajen de noche.} « La loi interdit que les mineurs travaillent la nuit. »
\end{itemize}

\textbf{Exception fondamentale} : si le sujet des deux verbes est \emph{le même}, on n'utilise pas \esp{que} + subjonctif, mais l'\textbf{infinitif} :
\esp{Quiero trabajar.} « Je veux travailler. » (même sujet) — mais \esp{Quiero que tú trabajes.} « Je veux que \emph{tu} travailles. » (sujets différents). Comparez : \esp{Pedimos mejorar nuestras condiciones} (« Nous demandons à améliorer nos conditions ») / \esp{Pedimos que el Estado mejore nuestras condiciones} (« Nous demandons que l'État améliore nos conditions »).

\courssub{Les tournures impersonnelles + que + subjonctif}

\begin{propriete}[Es necesario / urgente / justo / importante que + subjonctif]
Les tournures \textbf{impersonnelles} exprimant la nécessité, l'urgence, l'importance ou la justice sont suivies de \esp{que} + \textbf{subjonctif} :
\begin{center}
\esp{Es necesario que / es urgente que / es importante que / es justo que / es preciso que / conviene que} + \textbf{SUBJONCTIF}
\end{center}
\esp{Es justo que todos tengan acceso a la educación.} « Il est juste que tous aient accès à l'éducation. »
\end{propriete}

\begin{itemize}
\item \esp{Es urgente que el gobierno actúe.} « Il est urgent que le gouvernement agisse. »
\item \esp{Es necesario que se creen empleos para los jóvenes.} « Il est nécessaire que des emplois soient créés pour les jeunes. »
\item \esp{Conviene que los alumnos conozcan sus derechos.} « Il convient que les élèves connaissent leurs droits. »
\item \esp{Es importante que la prensa sea libre.} « Il est important que la presse soit libre. »
\end{itemize}

\begin{attention}[L'erreur n° 1 des francophones : l'indicatif après « es importante que »]
En français, après « il est important que », beaucoup d'élèves hésitent et écrivent parfois l'indicatif à l'oral (« il est important qu'il \emph{vient} »). En espagnol, cette hésitation n'existe pas : le subjonctif est \textbf{obligatoire} après toutes ces tournures. On ne dira \emph{jamais} \esp{es importante que la prensa \textbf{es} libre}, mais \esp{es importante que la prensa \textbf{sea} libre}. De même : \esp{es urgente que el Estado \textbf{aumente}} et non \esp{aumenta}.
\end{attention}

\courssub{Rappel : les formes du subjonctif présent}

Pour employer ces structures, il faut maîtriser parfaitement les formes. Rappel de formation : on part de la \textbf{1\textsuperscript{re} personne du singulier de l'indicatif présent}, on retire \esp{-o}, et on ajoute les terminaisons « inversées » : \esp{-e} pour les verbes en \esp{-ar}, \esp{-a} pour les verbes en \esp{-er/-ir}.

\begin{center}
\begin{tabular}{|l|l|l|l|l|l|}
\hline
\rowcolor{popblueL} Personne & hablar & comer & vivir & ser & ir \\
\hline
yo & hable & coma & viva & sea & vaya \\
\hline
tú & hables & comas & vivas & seas & vayas \\
\hline
él/ella/usted & hable & coma & viva & sea & vaya \\
\hline
nosotros & hablemos & comamos & vivamos & seamos & vayamos \\
\hline
vosotros & habléis & comáis & viváis & seáis & vayáis \\
\hline
ellos/ustedes & hablen & coman & vivan & sean & vayan \\
\hline
\end{tabular}
\end{center}

\begin{center}
\begin{tabular}{|l|l|l|l|}
\hline
\rowcolor{popblueL} Personne & estar & haber & saber / dar \\
\hline
yo & esté & haya & sepa / dé \\
\hline
tú & estés & hayas & sepas / des \\
\hline
él/ella/usted & esté & haya & sepa / dé \\
\hline
nosotros & estemos & hayamos & sepamos / demos \\
\hline
vosotros & estéis & hayáis & sepáis / deis \\
\hline
ellos/ustedes & estén & hayan & sepan / den \\
\hline
\end{tabular}
\end{center}

\begin{aretenir}[Les six irréguliers à connaître par cœur]
\esp{ser → sea} ; \esp{ir → vaya} ; \esp{haber → haya} ; \esp{estar → esté} ; \esp{saber → sepa} ; \esp{dar → dé} (avec accent pour le distinguer de la préposition \esp{de}). Attention aussi aux radicaux irréguliers : \esp{tener → tenga}, \esp{hacer → haga}, \esp{poder → pueda}, \esp{decir → diga}, \esp{pedir → pidan} (changement \esp{e→i}). Pour les verbes en \esp{-uir} (\esp{construir}, \esp{huir}, \esp{incluir}), le \esp{y} se maintient à toutes les personnes : \esp{construya, construyas, construya, construyamos, construyáis, construyan}.
\end{aretenir}

\courssub{Exprimer l'opinion : indicatif ou subjonctif ?}

C'est le point le plus délicat de la leçon. La règle est logique :

\begin{propriete}[Creer / pensar que : le mode dépend de la certitude]
\begin{itemize}
\item \textbf{Opinion affirmative} (certitude, conviction) : \esp{creer que / pensar que / estar seguro de que} + \textbf{INDICATIF}.\\
\esp{Creemos que el gobierno responderá.} « Nous croyons que le gouvernement répondra. »
\item \textbf{Opinion niée ou doutée} : \esp{no creer que / no pensar que / dudar que} + \textbf{SUBJONCTIF}.\\
\esp{No creo que sea suficiente.} « Je ne crois pas que ce soit suffisant. »
\end{itemize}
\end{propriete}

Autres exemples :
\begin{itemize}
\item \esp{Pienso que la situación es grave.} « Je pense que la situation est grave. » (indicatif)
\item \esp{No pensamos que la situación sea justa.} « Nous ne pensons pas que la situation soit juste. » (subjonctif)
\item \esp{¿Crees que el salario mínimo sea suficiente?} « Crois-tu que le salaire minimum soit suffisant ? » (à l'interrogative, le subjonctif exprime le doute)
\item \esp{Dudo que haya una solución rápida.} « Je doute qu'il y ait une solution rapide. »
\end{itemize}

\begin{attention}[Pièges à éviter absolument]
\begin{itemize}
\item \textbf{Jamais de subjonctif après une certitude affirmée} : \esp{Es verdad que hay pobreza.} « Il est vrai qu'il y a de la pauvreté. » (indicatif, car on affirme un fait). De même : \esp{es cierto que}, \esp{es evidente que}, \esp{es obvio que} + indicatif. Mais à la forme négative, le subjonctif revient : \esp{No es verdad que todo vaya bien.} « Il n'est pas vrai que tout aille bien. »
\item \textbf{\esp{Para que} + subjonctif, toujours} : \esp{Luchamos para que nuestros hijos tengan un futuro mejor.} « Nous luttons pour que nos enfants aient un avenir meilleur. » Ne confondez pas avec \esp{para} + infinitif (même sujet) : \esp{Trabajamos para vivir.} « Nous travaillons pour vivre. »
\item Ne dites pas \esp{es necesario de que} : la construction est \esp{es necesario que}, sans préposition.
\end{itemize}
\end{attention}

\courssub{Schéma récapitulatif : quel mode après quel déclencheur ?}

\begin{popfigure}
\begin{tikzpicture}[node distance=0.5cm and 1.1cm,
  trig/.style={rectangle, rounded corners, draw=popblue, fill=popblueL, align=center, font=\small, inner sep=5pt},
  mode/.style={rectangle, rounded corners, draw=poporange, fill=poporangeL, align=center, font=\small\bfseries, inner sep=5pt},
  lbl/.style={font=\scriptsize\itshape, text=popmuted}]
\node[trig, align=center] (t1){exigir / pedir / querer\\desear / prohibir que};
\node[trig, below=of t1, align=center] (t2){es necesario / urgente\\justo / importante que};
\node[trig, below=of t2, align=center] (t3){no creer / no pensar\\dudar que};
\node[trig, below=of t3] (t4) {para que};
\node[mode, right=of t2, xshift=1.6cm, align=center] (subj){SUBJONCTIF\\ \esp{Exigimos que se respete...}};
\node[trig, below=1.6cm of t4, align=center] (t5){creer / pensar que\\es verdad / cierto que};
\node[mode, right=of t5, xshift=1.6cm, align=center] (ind){INDICATIF\\ \esp{Creemos que responderá.}};
\draw[-{Stealth}, thick, popblue] (t1.east) -- (subj.north west);
\draw[-{Stealth}, thick, popblue] (t2.east) -- (subj.west);
\draw[-{Stealth}, thick, popblue] (t3.east) -- (subj.south west);
\draw[-{Stealth}, thick, popblue] (t4.east) |- (subj.south);
\draw[-{Stealth}, thick, popgreen] (t5.east) -- (ind.west);
\end{tikzpicture}
\legende{Carte mentale : le déclencheur détermine le mode de la subordonnée.}
\end{popfigure}

\begin{center}
\begin{tabularx}{\linewidth}{|X|X|X|}
\hline
\rowcolor{popblueL} Français & Español & Remarque \\
\hline
Nous exigeons que le salaire soit augmenté. & Exigimos que el salario sea aumentado. & Subjonctif dans les deux langues. \\
\hline
Il faut que les élèves étudient. & Es necesario que los alumnos estudien. & « Il faut que » = \esp{es necesario que}. \\
\hline
Je ne crois pas que ce soit juste. & No creo que sea justo. & Négation → subjonctif partout. \\
\hline
Il est vrai que la situation est grave. & Es verdad que la situación es grave. & Certitude → indicatif dans les deux langues. \\
\hline
Il est important qu'il vienne. & Es importante que venga. & En espagnol, aucune hésitation possible : subjonctif obligatoire. \\
\hline
\end{tabularx}
\end{center}

\courssub{Les connecteurs de l'argumentation}

Une lettre ouverte n'est pas une liste d'exigences : c'est un raisonnement organisé. Les connecteurs donnent la charpente du texte :

\begin{center}
\begin{tabularx}{\linewidth}{|X|X|}
\hline
\rowcolor{popblueL} Español & Français \\
\hline
en primer lugar & en premier lieu, d'abord \\
\hline
además & de plus, en outre \\
\hline
por otra parte & d'autre part \\
\hline
sin embargo & cependant, pourtant \\
\hline
por lo tanto / por eso & par conséquent, c'est pourquoi \\
\hline
en conclusión / para terminar & en conclusion, pour terminer \\
\hline
\end{tabularx}
\end{center}

Exemple de chaîne argumentative : \esp{En primer lugar, exigimos que se respete el convenio. Además, es urgente que se aumenten los salarios. Sin embargo, no creemos que el diálogo sea imposible. Por lo tanto, pedimos que se abran negociaciones. En conclusión, es justo que todos tengan un trabajo digno.} « En premier lieu, nous exigeons que la convention soit respectée. De plus, il est urgent que les salaires soient augmentés. Cependant, nous ne croyons pas que le dialogue soit impossible. Par conséquent, nous demandons que des négociations soient ouvertes. En conclusion, il est juste que tous aient un travail digne. »

\begin{methode}[Traduire « il faut que + subjonctif »]
\begin{enumerate}
\item \textbf{Identifier le sujet} de la subordonnée française : est-il précis ou indéfini ?
\item \textbf{Sujet précis} (une personne, un groupe déterminé) : traduire par \esp{es necesario que} + subjonctif. « Il faut que les jeunes trouvent du travail » → \esp{Es necesario que los jóvenes encuentren trabajo.}
\item \textbf{Sujet indéfini} (on, tout le monde, personne en particulier) : préférer \esp{hay que} + infinitif, plus naturel. « Il faut lutter contre le chômage » → \esp{Hay que luchar contra el desempleo.}
\item \textbf{Vérifier la forme} du subjonctif (radical irrégulier ? accent ?) avant de passer à la phrase suivante.
\end{enumerate}
\end{methode}

\begin{exoresolu}[Thème : phrases de revendication]
Traduisez en espagnol : 1. Nous exigeons que les droits des ouvriers soient respectés. 2. Il est urgent que l'État construise plus d'écoles. 3. Je ne crois pas que cette mesure soit suffisante. 4. Nous croyons que le dialogue est possible. 5. Il faut lutter pour que tous aient un salaire digne.
\tcblower
\textbf{Solution détaillée :}
\begin{enumerate}
\item \esp{Exigimos que se respeten los derechos de los obreros.} — Verbe d'exigence → subjonctif \esp{respeten} ; tournure passive avec \esp{se}.
\item \esp{Es urgente que el Estado construya más escuelas.} — Tournure impersonnelle → subjonctif ; \esp{construir} a un radical en \esp{-y-} : \esp{construya}.
\item \esp{No creo que esta medida sea suficiente.} — \esp{No creer que} → subjonctif \esp{sea} (irrégulier de \esp{ser}).
\item \esp{Creemos que el diálogo es posible.} — Opinion affirmative → \textbf{indicatif} \esp{es}. Piège classique : ne pas mettre \esp{sea} ici.
\item \esp{Hay que luchar para que todos tengan un salario digno.} — Sujet indéfini → \esp{hay que} + infinitif ; \esp{para que} → subjonctif \esp{tengan} (de \esp{tener}).
\end{enumerate}
\end{exoresolu}

\begin{exercice}[À vous de jouer]
Mettez le verbe entre parenthèses à la forme correcte (indicatif ou subjonctif présent) : 1. Pedimos que se \esp{(repartir)} mejor la riqueza. 2. Es verdad que \esp{(haber)} mucho desempleo. 3. No pensamos que la propuesta \esp{(ser)} justa. 4. Es importante que los jóvenes \esp{(participar)} en la vida pública. 5. Creo que el gobierno \esp{(escuchar)} nuestras demandas.
\end{exercice}

\begin{corrige}[Exercice « À vous de jouer »]
1. \esp{reparta} (pedir que + subjonctif). 2. \esp{hay} (\esp{es verdad que} affirmatif + indicatif). 3. \esp{sea} (\esp{no pensar que} + subjonctif). 4. \esp{participen} (tournure impersonnelle + subjonctif). 5. \esp{escucha} ou \esp{escuchará} (\esp{creer que} affirmatif + indicatif, présent ou futur).
\end{corrige}

\begin{aretenir}[L'essentiel de la section]
Trois idées à retenir pour le Bac : (1) volonté, exigence, nécessité, doute → \textbf{subjonctif} ; (2) opinion affirmative et certitude → \textbf{indicatif} ; (3) même sujet → infinitif (\esp{quiero trabajar} / \esp{quiero que trabajes}). Avec ces trois réflexes et les connecteurs \esp{en primer lugar}, \esp{además}, \esp{sin embargo}, \esp{por lo tanto}, vous disposez de toute la machinerie grammaticale nécessaire pour rédiger la lettre ouverte et la pétition de la section suivante.
\end{aretenir}

\coursec{La lettre ouverte et la pétition : structure et formules}

Dans la section précédente, nous avons appris à exprimer l'exigence et l'opinion (\esp{exigimos que}, \esp{opinamos que}, suivi du subjonctif). Ces outils grammaticaux trouvent maintenant leur application naturelle dans deux genres de textes très codifiés et très fréquents à l'examen : la \cle{carta abierta} (lettre ouverte) et la \cle{petición} (pétition). Savoir les rédiger, c'est savoir transformer une revendication en texte efficace, respectueux des conventions hispaniques.

\courssub{Qu'est-ce qu'une lettre ouverte ? Observation d'abord}

\begin{definition}[La carta abierta]
Une \cle{lettre ouverte} est une lettre formelle adressée officiellement à un destinataire précis (un président, un ministre, un proviseur, un maire) mais \textbf{rendue publique} : elle est publiée dans la presse ou affichée pour informer l'opinion et faire pression sur le destinataire. Elle se distingue de la lettre privée par sa diffusion et de l'article d'opinion par son adresse directe au destinataire.
\end{definition}

Observons d'abord un extrait de lettre ouverte rédigée par des élèves de Terminale d'un lycée de Yaoundé :

\begin{textebox}[Extrait d'une carta abierta]
Yaoundé, 15 de marzo de 2025

Señor Ministro de Educación Secundaria :

Por medio de la presente, los alumnos de Terminale del Liceo Bilingüe de Yaoundé nos dirigimos a usted para denunciar la falta de laboratorios de ciencias en nuestra institución.

Llevamos tres años estudiando física y química sin poder realizar ningún experimento práctico. Esta situación perjudica gravemente nuestra formación y nos pone en desventaja frente a otros establecimientos.

Por estas razones, exigimos que se construya un laboratorio equipado antes del próximo curso escolar y solicitamos que se nombre a un técnico de laboratorio.

Hacemos un llamamiento a las autoridades para que actúen con rapidez. Esperamos una respuesta rápida y concreta.

Atentamente,

Los delegados de la clase de Terminale A
\end{textebox}

\begin{aretenir}[Glossaire de l'extrait]
\begin{itemize}
\item \esp{por medio de la presente} : par la présente (formule administrative).
\item \esp{denunciar} : dénoncer, signaler publiquement.
\item \esp{la falta de} : le manque de, l'absence de.
\item \esp{perjudicar} : nuire à, porter préjudice à.
\item \esp{en desventaja} : en situation de désavantage.
\item \esp{equipado} : équipé (participe passé comme adjectif).
\item \esp{el próximo curso escolar} : la prochaine année scolaire.
\item \esp{un llamamiento} : un appel (solennel).
\item \esp{nombrar a} : nommer, désigner (quelqu'un) ; \esp{se nombre} (passif réfléchi) : soit nommé.
\end{itemize}
\end{aretenir}

\begin{experience}[Activité d'observation guidée]
Répondez par écrit en français :
\begin{enumerate}
\item Qui écrit ? À qui ? Dans quel contexte (pays, ville) ?
\item Quel est le problème concret exposé (les faits) ?
\item Repérez la phrase qui contient la revendication principale. Quel verbe la introduit ? Quel mode suit ce verbe ?
\item Quels connecteurs structurent l'argumentation ?
\item Quelle formule de politesse clôt la lettre ? Est-elle suivie d'une signature individuelle ou collective ?
\end{enumerate}
\end{experience}

\textbf{Analyse attendue :} On remarque que chaque élément occupe une place précise : c'est toute la \textbf{structure} que nous allons maintenant étudier. La revendication (\esp{exigimos que se construya...}) utilise le subjonctif ; la signature est collective (\esp{Los delegados...}).

\courssub{La structure de la carta abierta : les huit étapes}

\begin{propriete}[Structure rigide de la lettre ouverte]
Une lettre ouverte espagnole suit un plan en \textbf{huit étapes obligatoires}, dans cet ordre immuable :
\begin{enumerate}
\item \textbf{Lieu et date} : \esp{Yaoundé, 15 de marzo de 2025} (en haut à droite ; le mois \textbf{sans majuscule}).
\item \textbf{En-tête au destinataire} : \esp{Señor Presidente :}, \esp{Señoras y señores :}, \esp{Estimado señor Director :} (suivi de \textbf{deux points}, jamais de virgule ; la ligne suivante commence par une majuscule).
\item \textbf{Introduction} : qui écrit et pourquoi (\esp{los abajo firmantes nos dirigimos a usted para...}).
\item \textbf{Exposé du problème} : les faits, présentés objectivement (indicatif).
\item \textbf{Arguments} : deux ou trois raisons justifiant la demande (\esp{en primer lugar...}, \esp{además...}, \esp{por último...}).
\item \textbf{Revendication explicite} : \esp{Exigimos que...} / \esp{Solicitamos que...} + \textbf{subjonctif}.
\item \textbf{Appel à l'action} : \esp{Hacemos un llamamiento para que...} + \textbf{subjonctif}.
\item \textbf{Formule de politesse et signature collective} : \esp{Atentamente,} puis le nom du groupe (\esp{Los delegados de clase}, \esp{El comité de padres de alumnos}).
\end{enumerate}
\end{propriete}

\begin{popfigure}
\begin{tikzpicture}[node distance=0.35cm, every node/.style={font=\small, align=center}]
\node[draw=popblue, fill=popblueL, rounded corners, minimum width=8cm] (e1) {1. Lugar y fecha : \esp{Yaoundé, 15 de marzo de 2025}};
\node[draw=popblue, fill=popblueL, rounded corners, minimum width=8cm, below=of e1] (e2) {2. Destinatario : \esp{Señor Presidente :}};
\node[draw=poppurple, fill=poppurpleL, rounded corners, minimum width=8cm, below=of e2] (e3) {3. Introducción : quién escribe y por qué};
\node[draw=poppurple, fill=poppurpleL, rounded corners, minimum width=8cm, below=of e3] (e4) {4. Exposición del problema (hechos objetivos)};
\node[draw=poporange, fill=poporangeL, rounded corners, minimum width=8cm, below=of e4] (e5) {5. Argumentos (2 o 3 razones con conectores)};
\node[draw=poporange, fill=poporangeL, rounded corners, minimum width=8cm, below=of e5] (e6) {6. Reivindicación : \esp{Exigimos que...} + subjuntivo};
\node[draw=popgreen, fill=popgreenL, rounded corners, minimum width=8cm, below=of e6] (e7) {7. Llamamiento a la acción : \esp{para que...} + subj.};
\node[draw=popgreen, fill=popgreenL, rounded corners, minimum width=8cm, below=of e7] (e8) {8. Despedida y firma colectiva : \esp{Atentamente,...}};
\foreach \a/\b in {e1/e2, e2/e3, e3/e4, e4/e5, e5/e6, e6/e7, e7/e8}
  \draw[-{Stealth}, thick, popdark] (\a) -- (\b);
\end{tikzpicture}
\legende{Organigramme : les huit étapes de la carta abierta}
\end{popfigure}

\courssub{Les formules : ouverture, revendication, clôture}

\begin{aretenir}[Formules d'ouverture (solennelles)]
Pour ouvrir la lettre, on annonce l'objet de l'écrit :
\begin{itemize}
\item \esp{Por medio de la presente, ...} : Par la présente, ... (très formel, administratif).
\item \esp{Nos dirigimos a ustedes para denunciar / solicitar / expresar...} : Nous nous adressons à vous pour dénoncer / demander / exprimer...
\item \esp{Los abajo firmantes queremos manifestar nuestra preocupación por...} : Les soussignés veulent manifester leur préoccupation au sujet de...
\item \esp{Por la presente, me dirijo a usted en mi calidad de...} : Par la présente, je m'adresse à vous en ma qualité de... (si signataire unique, rare pour une carta abierta).
\end{itemize}
\end{aretenir}

\begin{aretenir}[Formules de revendication et de clôture]
\textbf{Revendication} (toujours suivies de \esp{que} + \textbf{subjonctif}) :
\begin{itemize}
\item \esp{Exigimos que se respeten nuestros derechos.} : Nous exigeons que nos droits soient respectés.
\item \esp{Solicitamos que se nombre un profesor de informática.} : Nous demandons qu'un professeur d'informatique soit nommé.
\item \esp{Pedimos que se repare la carretera de Douala.} : Nous demandons que la route de Douala soit réparée.
\item \esp{Hacemos un llamamiento para que la población participe.} : Nous lançons un appel pour que la population participe.
\item \esp{Reclamamos que se cumpla la ley.} : Nous réclamons que la loi soit appliquée.
\end{itemize}
\textbf{Clôture} (formules fixes, sans subjonctif) :
\begin{itemize}
\item \esp{Esperamos una respuesta rápida.} : Nous espérons une réponse rapide.
\item \esp{En espera de sus noticias, le saludamos atentamente.} : Dans l'attente de vos nouvelles, nous vous saluons cordialement.
\item \esp{Atentamente,} : Cordialement, / Veuillez agréer... (formule finale unique, suivie de la signature).
\item \esp{Reciban un cordial saludo,} : Recevez un cordial salut, (alternative).
\end{itemize}
\end{aretenir}

\begin{exemplebox}[Exemples traduits et analysés]
\esp{Por medio de la presente, solicitamos que se construya un laboratorio.} \\
« Par la présente, nous demandons la construction d'un laboratoire. » (Passif réfléchi \esp{se construya} = \emph{que l'on construise / qu'il soit construit}, subjonctif présent).

\esp{Nos dirigimos a ustedes para protestar contra el aumento de las tasas escolares.} \\
« Nous nous adressons à vous pour protester contre l'augmentation des frais scolaires. » (\esp{protestar contra} + infinitif de but).

\esp{En espera de sus noticias, reciban un cordial saludo.} \\
« Dans l'attente de vos nouvelles, recevez un cordial salut. » (Impératif \emph{ustedes} \esp{reciban} comme formule de politesse).
\end{exemplebox}

\begin{attention}[Erreurs fréquentes des francophones — \textbf{Points de vigilance Bac}]
\begin{itemize}
\item \textbf{Salutation} : Ne commence \textbf{jamais} une lettre formelle par \esp{Hola} ni \esp{Querido} (réservé aux proches). Utilise \esp{Estimado señor :}, \esp{Señoras y señores :}, \esp{Muy señor nuestro :}, \esp{Señor Director :}.
\item \textbf{Ponctuation de l'en-tête} : En espagnol, l'en-tête se termine par \textbf{deux points} (\esp{Señor Director :}) et le corps du texte commence à la ligne suivante avec une \textbf{majuscule}. En français, on met une virgule.
\item \textbf{Faux ami \esp{demandar}} : \esp{Demandar} = « exiger avec force », « intenter un procès », « réclamer un droit ». Pour « demander poliment », utilise \esp{solicitar} ou \esp{pedir}.
\item \textbf{Formule de politesse} : \esp{Atentamente} (seul, avec une majuscule, suivi d'une virgule en espagnol \esp{Atentamente,}) suffit. N'invente pas de traduction littérale de « Veuillez agréer l'expression de mes sentiments distingués ».
\item \textbf{Subjonctif obligatoire} : Après \esp{exigimos que}, \esp{solicitamos que}, \esp{pedimos que}, \esp{reclamamos que}, \esp{hacemos un llamamiento para que} $\rightarrow$ \textbf{Subjonctif présent}. \textbf{Erreur} : \esp{exigimos que se construye} $\rightarrow$ \textbf{Correct} : \esp{exigimos que se construya}.
\item \textbf{Accord du participe passé passif} : \esp{se construya} (invariable avec \emph{se}), mais \esp{la carta escrita} (accord si adjectif). Ici \esp{se nombre}, \esp{se repare}, \esp{se construya} sont invariables.
\end{itemize}
\end{attention}

\courssub{La petición : un texte bref et collectif}

\begin{definition}[La petición]
La \cle{pétition} est un texte \textbf{bref} qui expose des faits et formule une demande, suivi d'une \textbf{liste de signatures} (nom + date). Contrairement à la lettre ouverte, elle n'a ni en-tête au destinataire ni formule de politesse développée : sa force vient du \textbf{nombre de signataires}. Elle peut être manuscrite ou numérique (Change.org, etc.).
\end{definition}

\begin{propriete}[Structure de la petición]
\begin{enumerate}
\item \textbf{Titre} : \esp{Petición para} + infinitif : \esp{Petición para mejorar el comedor escolar} (Pétition pour améliorer la cantine scolaire).
\item \textbf{Texte bref} : les faits en quelques phrases (indicatif), puis la demande introduite par \esp{Los abajo firmantes pedimos / solicitamos / exigimos que...} + \textbf{subjonctif}.
\item \textbf{Liste de signatures} : \esp{Firmas :} suivi des noms et dates (et éventuellement classe/fonction).
\end{enumerate}
\end{propriete}

\begin{textebox}[Mini-modèle de petición]
Petición para mejorar el comedor escolar

Los alumnos del Liceo de Bafoussam constatamos que las comidas del comedor son de mala calidad y que el local está sucio. Los abajo firmantes pedimos que se mejore la calidad de las comidas, que se limpie el local todos los días y que se aumente el número de mesas.

Esperamos que la dirección tome medidas urgentes.

Firmas :
1. Amina Mbarga — 12/03/2025
2. José Etoundi — 12/03/2025
3. María Ngo Bell — 13/03/2025
\end{textebox}

\begin{aretenir}[Glossaire de la petición]
\begin{itemize}
\item \esp{el comedor escolar} : la cantine scolaire.
\item \esp{los abajo firmantes} : les soussignés (invariable, toujours au pluriel).
\item \esp{constatamos} : nous constatons (indicatif présent, fait objectif).
\item \esp{sucio} : sale.
\item \esp{tomar medidas} : prendre des mesures.
\item \esp{urgentes} : urgentes.
\item \esp{se mejore, se limpie, se aumente} : subjonctifs présents passifs (impersonnels).
\item \esp{tome} : subjonctif présent de \esp{tomar} (après \esp{esperamos que}).
\end{itemize}
\end{aretenir}

\begin{popfigure}
\begin{tikzpicture}[node distance=0.4cm and 1.5cm, every node/.style={font=\small}]
\node[draw=popblue, fill=popblueL, rounded corners, align=center, minimum width=6cm] (t1) {\textbf{CARTA ABIERTA}};
\node[draw=popblue, rounded corners, align=left, minimum width=6cm, below=of t1, text width=5.6cm] (c1) {
-- Lugar y fecha\\
-- Destinatario (\esp{Señor... :})\\
-- Introducción (\esp{nos dirigimos a...})\\
-- Problema + argumentos\\
-- Reivindicación (\esp{Exigimos que...})\\
-- Llamamiento (\esp{para que...})\\
-- \esp{Atentamente} + firma colectiva};
\node[draw=poporange, fill=poporangeL, rounded corners, align=center, minimum width=6cm, right=of t1] (t2) {\textbf{PETICIÓN}};
\node[draw=poporange, rounded corners, align=left, minimum width=6cm, below=of t2, text width=5.6cm] (c2) {
-- Título : \esp{Petición para} + inf.\\
-- Texto breve (hechos)\\
-- Demanda : \esp{pedimos que...} + subj.\\
-- Sin fórmula de despedida\\
-- Lista de firmas + fechas\\
-- Fuerza = número de firmas};
\draw[{Stealth}-{Stealth}, very thick, poppurple] (c1.east) -- (c2.west) node[midway, above, font=\small\itshape, align=center, text width=1.5cm] {Mismo objetivo:\\ presión pública};
\end{tikzpicture}
\legende{Schéma comparatif : carta abierta vs petición}
\end{popfigure}

\courssub{Le registre : vouvoiement collectif et ton ferme mais courtois}

\begin{propriete}[Le registre de la revendication formelle]
\begin{itemize}
\item \textbf{Vouvoiement} : On s'adresse au destinataire par \esp{usted} / \esp{ustedes} (verbes à 3e personne), jamais par \esp{tú}.
\item \textbf{Voix collective} : Les auteurs s'expriment à la 1re personne du pluriel (\esp{nosotros}, \esp{nos dirigimos}, \esp{exigimos}, \esp{los abajo firmantes}). C'est une parole \textbf{collective}, porteuse de poids.
\item \textbf{Ton ferme mais courtois} : On revendique sans insulter. On préfère \esp{denunciamos una situación injusta} (nous dénonçons une situation injuste) aux attaques personnelles (\esp{usted es incompetente}). La fermeté réside dans la clarté de l'exigence, pas dans l'agressivité.
\item \textbf{Clarté syntaxique} : Une idée par phrase. Connecteurs visibles : \esp{en primer lugar}, \esp{en segundo lugar}, \esp{además}, \esp{por otra parte}, \esp{por lo tanto}, \esp{por consiguiente}.
\end{itemize}
\end{propriete}

\begin{center}
\begin{tabularx}{\linewidth}{|X|X|X|}
\hline
\rowcolor{popblueL} \textbf{Français} & \textbf{Español} & \textbf{Remarque / Contexte} \\
\hline
Par la présente & Por medio de la presente & Ouverture solennelle, administrative \\
\hline
Les soussignés & Los abajo firmantes & Invariable, collectif, très fréquent \\
\hline
Nous exigeons que... & Exigimos que... + subj. & Registre fort, droit, revendication \\
\hline
Nous demandons que... & Solicitamos que... + subj. & Registre formel, poli, administratif \\
\hline
Nous demandons (familier) & Pedimos que... + subj. & Registre standard, direct \\
\hline
Nous lançons un appel & Hacemos un llamamiento para que... + subj. & Mobilisation, appel à l'action \\
\hline
Dans l'attente de vos nouvelles & En espera de sus noticias & Transition vers la clôture \\
\hline
Cordialement / Respectueusement & Atentamente / Un cordial saludo & Clôture sobre, standard \\
\hline
\end{tabularx}
\end{center}

\coursec{Présentation orale d'un pays hispanophone}

Le programme de Terminale A (Module 4) exige la capacité de présenter un pays hispanophone de manière structurée, en mobilisant géographie, histoire, culture et économie. Cet exercice oral (souvent évalué en contrôle continu ou au Bac oral) suit des codes précis.

\courssub{Structure type d'un exposé oral (3 à 5 minutes)}

\begin{methode}[Préparer et réussir sa présentation orale]
\begin{enumerate}
\item \textbf{Introduction (30 s)} : Accroche (chiffre, anecdote, citation), annonce du pays, localisation, capitale, population. \textit{Ex: « ¿Sabían que la Guinea Ecuatorial es el único país africano donde el español es lengua oficial? »}
\item \textbf{Géographie et climat (45 s)} : Situation, relief, climat, villes principales, ressources naturelles.
\item \textbf{Histoire clé (45 s)} : Période précolombienne / précoloniale, colonisation, indépendance, époque contemporaine. Utiliser \esp{el pretérito indefinido} (actions ponctuées) et \esp{el pretérito imperfecto} (descriptions, contextes).
\item \textbf{Culture et société (45 s)} : Langues coofficielles, littérature, musique, fêtes, gastronomie, religion. Mentionner un auteur/artiste célèbre.
\item \textbf{Économie et défis actuels (45 s)} : Secteurs principaux (pétrole, cacao, tourisme, technologie), défis (chômage, inégalités, environnement), liens avec le Cameroun si pertinent.
\item \textbf{Conclusion et ouverture (30 s)} : Bilan personnel, pourquoi ce pays vous intéresse, question au jury / invitation au dialogue.
\end{enumerate}
\end{methode}

\begin{aretenir}[Lexique indispensable pour l'exposé]
\begin{center}
\begin{tabularx}{\linewidth}{|X|X|X|}
\hline
\rowcolor{popblueL} \textbf{Thème} & \textbf{Español} & \textbf{Français} \\
\hline
Géographie & \esp{limitar con, bañado por, la cordillera, el altiplano, la selva, la costa} & Border, être baigné par, chaîne de montagnes, haut plateau, jungle, côte \\
\hline
Histoire & \esp{la conquista, la independencia, el virreinato, la dictadura, la transición, proclamar} & La conquête, l'indépendance, la vice-royauté, la dictature, la transition, proclamer \\
\hline
Culture & \esp{la lengua cooficial, el folklore, la gastronomía, patrimonio de la humanidad, identidad} & Langue coofficielle, folklore, gastronomie, patrimoine de l'humanité, identité \\
\hline
Économie & \esp{el PIB, la renta per cápita, sector primario/secundario/terciario, inversión extranjera, desigualdad} & PIB, revenu par tête, secteur primaire/secondaire/tertiaire, investissement étranger, inégalité \\
\hline
Connecteurs oral & \esp{en primer lugar, cabe destacar, por un lado... por otro, en definitiva, para concluir} & En premier lieu, il faut souligner, d'un côté... de l'autre, en définitive, pour conclure \\
\hline
\end{tabularx}
\end{center}
\end{aretenir}

\begin{propriete}[Outils grammaticaux clés pour l'exposé]
\begin{itemize}
\item \textbf{Passif réfléchi (\emph{se} + verbe 3e pers.)} : Impersonnel, objectif. \esp{En Guinea Ecuatorial \textbf{se habla} español y francés.} / \esp{\textbf{Se exporta} petróleo y cacao.}
\item \textbf{Impersonnel avec \emph{se}} : \esp{\textbf{Se vive} bien en Málaga.} / \esp{\textbf{Se come} muy bien en México.}
\item \textbf{Passif analytique (\emph{ser} + participe)} : Pour insister sur l'agent. \esp{La independencia \textbf{fue proclamada} por...}
\item \textbf{Temps du récit historique} : \esp{El país \textbf{logró} su independencia en 1968 (indefinido). \textbf{Tenía} una economía basada en el cacao (imperfecto).}
\item \textbf{Subjonctif dans les clauses de but/opinion} : \esp{Es importante \textbf{que} los jóvenes \textbf{conozcan} su historia.} / \esp{El gobierno quiere \textbf{que se reduzca} la pobreza.}
\end{itemize}
\end{propriete}

\begin{exemplebox}[Fiche de préparation : Guinea Ecuatorial (modèle pour le Bac)]
\textbf{País :} Guinea Ecuatorial \quad \textbf{Capital :} Malabo (isla de Bioko) / Oyala (futura, en el continente) \\
\textbf{Población :} ~1,7 millones \quad \textbf{Idiomas oficiales :} Español, francés, portugués \\
\textbf{Geografía :} Parte insular (Bioko, Annobón) + parte continental (Río Muni). Clima ecuatorial, selva tropical. \\
\textbf{Historia :} Colonizada por España (1778--1968). Independencia 12 octubre 1968. Macías Nguema (dictadura), luego Obiang (desde 1979). \\
\textbf{Cultura :} Literatura en español (Juan Tomás Ávila Laurel, Donato Ndongo). Música: \emph{balélé}, \emph{ibanga}. Fiestas: 12 octubre (Día de la Independencia). \\
\textbf{Economía :} Petróleo (90\% exportaciones), gas, maderas, cacao. Desafío: diversificación, distribución de la riqueza, derechos humanos. \\
\textbf{Vínculo Camerún :} Vecino fronterizo (Río Muni), comunidad camerunesa en Malabo, cooperación CEMAC.
\end{exemplebox}

\begin{experience}[Simulation d'examen oral : « Présentez un pays hispanophone »]
\textbf{Consigne :} Choisissez un pays (Espagne, Mexique, Argentine, Colombie, Pérou, Chili, Guinée équatoriale...). Préparez une fiche bristol (mots-clés seulement, pas de phrases entières). Présentez 4 minutes devant la classe. La classe note sur une grille : Structure / Vocabulaire spécifique / Grammaire (temps, subjonctif, passif) / Prononciation / Intérêt du contenu.
\textbf{Astuce :} Intégrez une comparaison avec le Cameroun (climat, économie, langue) pour ancrer votre exposé dans votre réalité.
\end{experience}

\begin{savaistu}[La Guinée équatoriale : l'hispanité au cœur de l'Afrique]
Seul pays d'Afrique où l'espagnol est langue officielle (depuis 1844, confirmé à l'indépendance en 1968), la Guinée équatoriale (capital Malabo sur l'île de Bioko) est membre de l'Union africaine, de la CEMAC et de l'Association des académies de la langue espagnole. Sa littérature en espagnol est vive (Juan Tomás Ávila Laurel, María Nsue Angüe). Le français et le portugais sont aussi officiels depuis 1998 et 2010. Pour un élève camerounais, c'est le pays hispanophone le plus proche géographiquement et culturellement : un pont naturel vers l'hispanité.

\begin{center}
\begin{tikzpicture}[scale=0.8, every node/.style={font=\tiny}]
\draw[popblue, thick] (0,0) rectangle (6,4);
\draw[popblue, thick] (0,3.5) -- (6,3.5);
\node[align=center, popdark] at (3,3.75) {\textbf{GUINÉE ÉQUATORIALE : Repères clés}};
\node[align=left, anchor=north west, popdark] at (0.2,3.3) {
\textbullet\ Capital: Malabo (Bioko) / Oyala (continente)\\
\textbullet\ Superficie: 28 051 km²\\
\textbullet\ Pop: ~1,7 M hab.\\
\textbullet\ Langues: Español, Français, Portugués\\
\textbullet\ Monnaie: Franc CFA (XAF)\\
\textbullet\ Ressources: Pétrole, Gaz, Bois, Cacao\\
\textbullet\ Indépendance: 12 oct. 1968 (d'Espagne)\\
\textbullet\ Auteur clé: Juan Tomás Ávila Laurel};
\end{tikzpicture}
\end{center}
\legende{Carte mentale simplifiée : Guinée équatoriale (à reproduire au tableau)}
\end{savaistu}

\coursec{Production guidée et entraînement au Bac}

\courssub{Méthode : commenter un texte revendicatif}

\begin{methode}[Le commentaire de texte revendicatif en trois parties (Bac écrit)]
\begin{enumerate}
\item \textbf{Introduction} : Présente le texte (auteur ou collectif, nature : lettre ouverte, pétition, manifeste ; thème ; contexte de publication). Formule une \textbf{problématique} précise : \textit{« ¿Cómo intenta el colectivo autor convencer al destinatario y movilizar la opinión pública ? »} ou \textit{« ¿Qué estrategias argumentativas y lingüísticas se emplean para legitimar la reivindicación ? »}. Annonce le plan (2 ou 3 axes).
\item \textbf{Développement (2 ou 3 parties)} : Chaque partie analyse un aspect :
\begin{itemize}
\item \textbf{Les arguments} : Factuels (chiffres, situations), axiologiques (valeurs : justice, égalité, droit à l'éducation), émotionnels.
\item \textbf{Les procédés linguistiques} : Subjonctif d'exigence (\esp{exigimos que...}), « nous » collectif (\esp{los abajo firmantes}), connecteurs logiques, répétitions (anaphore), questions rhétoriques, champs lexicaux (droits, devoirs, urgence), modalisateurs (\esp{es urgente}, \esp{es inadmisible}).
\item \textbf{Le ton et la visée} : Ferme, courtois, indigné, constructif ? Visée : obtenir une réponse, faire pression, mobiliser.
\end{itemize}
\textbf{Citez le texte en espagnol} (guillemets \esp{« ... »}) pour justifier chaque remarque.
\item \textbf{Conclusion} : Bilan synthétique (le texte est-il efficace ?). Appréciation personnelle justifiée. Ouverture sur une question plus large (droits des élèves, conditions de travail, citoyenneté active, rôle de la jeunesse).
\end{enumerate}
\end{methode}

\begin{exoresolu}[Rédiger le début d'une lettre ouverte]
Énoncé : Les élèves de Terminale du lycée de Garoua écrivent au proviseur pour demander la création d'un club de théâtre hispanophone. Rédigez le lieu et la date, l'en-tête, l'introduction et la revendication.
\tcblower
Solution détaillée :

\esp{Garoua, 20 de mayo de 2025}

\esp{Señor Director :}

\esp{Por medio de la presente, los alumnos de Terminale nos dirigimos a usted para expresar nuestro deseo de crear un club de teatro hispanohablante en el liceo. El teatro nos permitiría practicar el español de forma activa y descubrir la cultura de los países hispanos.}

\esp{Por estas razones, solicitamos que se autorice la creación del club y que se ponga una sala a nuestra disposición los miércoles por la tarde.}

\esp{En espera de sus noticias, le saludamos atentamente.}

\esp{Los delegados de Terminale A}

\textbf{Commentaire de la correction :} On respecte les huit étapes (ici condensées pour l'exercice), l'en-tête se termine par deux points, la revendication emploie \esp{solicitamos que} + subjonctif (\esp{se autorice}, \esp{se ponga} -- passifs réfléchis), et la clôture reste sobre et courtoise (\esp{En espera de sus noticias... Atentamente}).
\end{exoresolu}

\begin{exercice}[Entraînement Bac : Pétition complète]
Rédige une pétition complète (titre, texte de cinq à huit lignes, trois signatures) intitulée \esp{Petición para crear una biblioteca en el barrio}, destinée au maire de ta ville. Utilise au moins deux formules de revendication vues dans cette leçon (\esp{pedimos que}, \esp{solicitamos que}, \esp{exigimos que}, \esp{hacemos un llamamiento para que}) et veille scrupuleusement à l'emploi du \textbf{subjonctif}. Respecte la structure en 3 parties (Titre, Texte, Firmas).
\end{exercice}

\begin{corrige}[Exercice : Pétition pour une bibliothèque de quartier]
\textbf{Proposition de rédaction :}

\esp{Petición para crear una biblioteca en el barrio}

\esp{Los vecinos del barrio Mvog-Ada constatamos la falta de un espacio cultural y educativo accesible para los jóvenes y los estudiantes. No existe ninguna biblioteca pública en un radio de cinco kilómetros, lo que dificulta el estudio y la lectura.}

\esp{Los abajo firmantes \textbf{solicitamos que se construya} una biblioteca municipal equipada con salas de lectura y conexión a internet, y \textbf{pedimos que se asigne} un presupuesto anual para la compra de libros en español y en francés.}

\esp{\textbf{Hacemos un llamamiento para que} el Ayuntamiento priorice este proyecto en el próximo plan de inversiones.}

\esp{Firmas :} \\
\esp{1. Paul Atangana — 20/05/2025} \\
\esp{2. Fatoumata Diallo — 20/05/2025} \\
\esp{3. Carlos Nguema — 21/05/2025}

\textbf{Analyse grammaticale :}
\begin{itemize}
\item \esp{constatamos} (indicatif : fait objectif), \esp{dificulta} (indicatif).
\item \esp{solicitamos que se construya} : \esp{construya} = subjonctif présent 3e sg de \esp{construir} (passif réfléchi).
\item \esp{pedimos que se asigne} : \esp{asigne} = subjonctif présent 3e sg de \esp{asignar} (passif réfléchi, changement g $\rightarrow$ gu devant e).
\item \esp{hacemos un llamamiento para que... priorice} : \esp{priorice} = subjonctif présent 3e sg de \esp{priorizar} (changement z $\rightarrow$ c devant e).
\end{itemize}
\end{corrige}

\begin{exercice}[Entraînement Bac : Thème / Version ciblé]
\textbf{Thème (Français $\rightarrow$ Espagnol) :} Traduisez en espagnol formel.
\begin{enumerate}
\item Par la présente, les parents d'élèves exigent que soit réparé le toit de l'école avant la saison des pluies.
\item Nous nous adressons à vous pour dénoncer le manque de manuels scolaires en espagnol.
\item Les soussignés demandent qu'un professeur de musique soit nommé.
\item Nous lançons un appel pour que les autorités écoutent la voix des jeunes.
\item Dans l'attente de votre réponse, nous vous saluons cordialement.
\end{enumerate}
\textbf{Version (Espagnol $\rightarrow$ Français) :} Traduisez en français.
\begin{enumerate}
\item \esp{Señor Alcalde: Por medio de la presente, los vecinos nos dirigimos a usted para reclamar que se limpie el mercado municipal.}
\item \esp{Exigimos que se respete el derecho a la educación y que se construyan más aulas.}
\item \esp{Los abajo firmantes hacemos un llamamiento para que no se corte el agua potable.}
\item \esp{Esperamos una solución inmediata. Atentamente, el Comité de Barrio.}
\end{enumerate}
\end{exercice}

\begin{corrige}[Exercice Thème / Version]
\textbf{Thème :}
\begin{enumerate}
\item \esp{Por medio de la presente, los padres de alumnos exigen que se repare el techo de la escuela antes de la temporada de lluvias.}
\item \esp{Nos dirigimos a ustedes para denunciar la falta de manuales escolares de español.}
\item \esp{Los abajo firmantes solicitamos que se nombre a un profesor de música.}
\item \esp{Hacemos un llamamiento para que las autoridades escuchen la voz de los jóvenes.}
\item \esp{En espera de su respuesta, les saludamos atentamente.}
\end{enumerate}
\textbf{Version :}
\begin{enumerate}
\item Monsieur le Maire, par la présente, nous, les habitants, nous nous adressons à vous pour réclamer que le marché municipal soit nettoyé.
\item Nous exigeons que le droit à l'éducation soit respecté et que davantage de salles de classe soient construites.
\item Les soussignés lançons un appel pour que l'eau potable ne soit pas coupée.
\item Nous espérons une solution immédiate. Cordialement, le Comité de Quartier.
\end{enumerate}
\end{corrige}

\coursec{Présentation orale d'un pays hispanophone}

Après avoir appris à défendre une cause par écrit grâce à la lettre ouverte et à la pétition, nous allons maintenant prendre la parole à l'oral. Présenter un pays hispanophone est un exercice classique de l'expression orale en Terminale : il combine la maîtrise des données factuelles (chiffres, dates, lieux), un lexique précis et une structure claire. C'est aussi une compétence très concrète : demain, dans le monde du travail, vous devrez peut-être présenter un pays partenaire à des collègues ou à des clients.

\courssub{Observer d'abord : une présentation modèle du Mexique}

Avant d'énoncer la moindre règle, lisons et écoutons attentivement cette présentation complète. Observez comment l'orateur annonce chaque partie, enchaîne les données et conclut.

\begin{textebox}[Presentación oral : México]
Buenos días a todos. Hoy voy a presentarles México, un país fascinante de América del Norte.

En primer lugar, hablaré de la geografía. México está situado en el sur de América del Norte. Limita al norte con los Estados Unidos, al sur con Guatemala y Belice, al este con el golfo de México y al oeste con el océano Pacífico. Su capital es Ciudad de México, una de las ciudades más grandes del mundo. El país cuenta con más de 125 millones de habitantes y tiene una superficie de casi dos millones de kilómetros cuadrados.

En segundo lugar, pasemos a la historia. México fue el corazón de grandes civilizaciones como los aztecas y los mayas. En 1519, los españoles, dirigidos por Hernán Cortés, llegaron al territorio. Después de tres siglos de dominación colonial, México inició su proceso de independencia el 16 de septiembre de 1810, fecha que se celebra cada año con grandes fiestas populares.

En tercer lugar, hablemos de la cultura. La lengua oficial es el español, pero también se hablan lenguas indígenas como el náhuatl. México es conocido por su gastronomía: los tacos, las enchiladas y el guacamole son famosos en todo el mundo. Entre sus personajes célebres, podemos mencionar a la pintora Frida Kahlo y al escritor Octavio Paz, premio Nobel de Literatura.

Por último, en el plano económico, su moneda es el peso mexicano. El país cuenta con importantes recursos: petróleo, plata y productos agrícolas como el aguacate. Exporta coches y productos electrónicos, sobre todo hacia los Estados Unidos. Además, el turismo es una fuente esencial de ingresos: millones de visitantes descubren cada año las playas de Cancún y las pirámides de Teotihuacán.

Para concluir, México es un país rico por su historia, su cultura y su economía dinámica. Muchas gracias por su atención. ¿Tienen preguntas?
\end{textebox}

\textbf{Glossaire :}
\begin{itemize}
\item \esp{limita al norte con} : il est limité au nord par / il borde au nord ;
\item \esp{cuenta con} : il compte, il dispose de ;
\item \esp{la superficie} : la superficie ;
\item \esp{inició su proceso de independencia} : il a entamé son processus d'indépendance (le 16/09/1810 marque le \emph{Grito de Dolores}, début de la guerre ; l'indépendance effective date de 1821) ;
\item \esp{el aguacate} : l'avocat ;
\item \esp{una fuente de ingresos} : une source de revenus ;
\item \esp{¿Tienen preguntas?} : Avez-vous des questions ?
\end{itemize}

\textbf{Questions de compréhension :}
\begin{enumerate}
\item ¿Con qué países limita México al norte y al sur ?
\item ¿Qué evento histórico ocurrió el 16 de septiembre de 1810 y cómo se celebra ?
\item ¿Qué recursos económicos y qué produits touristiques se mencionan ?
\end{enumerate}

\courssub{Le plan en quatre volets}

Toute présentation orale d'un pays suit le même squelette, facile à mémoriser et rassurant pour l'orateur comme pour l'auditoire.

\begin{definition}[Los cuatro ejes de la presentación]
\begin{enumerate}
\item \textbf{La geografía} : \esp{situación} (situation), \esp{capital}, \esp{población} (population), \esp{superficie}, pays frontaliers, reliefs, climat.
\item \textbf{La historia} : époque précolombienne, colonisation, \esp{independencia}, grandes dates, personnages historiques.
\item \textbf{La cultura} : \esp{lengua oficial}, traditions et fêtes, \esp{gastronomía}, musique, personnalités (artistes, écrivains, sportifs).
\item \textbf{La economía} : \esp{moneda} (monnaie), \esp{recursos naturales}, \esp{exportaciones}, industrie, \esp{turismo}.
\end{enumerate}
\end{definition}

\begin{center}
\begin{tabularx}{\linewidth}{|l|X|X|}
\hline
\rowcolor{popblueL} \textbf{Volet} & \textbf{Expression espagnole} & \textbf{Français} \\
\hline
Géographie & \esp{Está situado en...} & Il est situé en/au... \\
\hline
Géographie & \esp{Limita al norte con...} & Il est limité au nord par... \\
\hline
Géographie & \esp{Cuenta con... habitantes.} & Il compte... habitants. \\
\hline
Géographie & \esp{Tiene una superficie de...} & Il a une superficie de... \\
\hline
Histoire & \esp{Proclamó / Consiguió su independencia en...} & Il proclama / obtient son indépendance en... \\
\hline
Histoire & \esp{Fue colonizado por...} & Il fut colonisé par... \\
\hline
Culture & \esp{La lengua oficial es...} & La langue officielle est... \\
\hline
Culture & \esp{Es conocido por...} & Il est connu pour... \\
\hline
Culture & \esp{Entre sus personajes célebres...} & Parmi ses personnalités célèbres... \\
\hline
Économie & \esp{Su moneda es...} & Sa monnaie est... \\
\hline
Économie & \esp{Exporta... / Importa...} & Il exporte... / Il importe... \\
\hline
Économie & \esp{El turismo es una fuente de ingresos.} & Le tourisme est une source de revenus. \\
\hline
\end{tabularx}
\end{center}

\begin{exemplebox}[Ejemplos traducidos]
\begin{itemize}
\item \esp{México cuenta con más de 125 millones de habitantes.} = Le Mexique compte plus de 125 millions d'habitants.
\item \esp{Su moneda es el peso.} = Sa monnaie est le peso.
\item \esp{Perú limita al norte con Ecuador y Colombia.} = Le Pérou est limité au nord par l'Équateur et la Colombie.
\item \esp{España es conocida por el flamenco y la paella.} = L'Espagne est connue pour le flamenco et la paella.
\item \esp{Argentina declaró su independencia en 1816.} = L'Argentine a déclaré son indépendance en 1816.
\end{itemize}
\end{exemplebox}

\begin{popfigure}
\begin{center}
\begin{tikzpicture}[mindmap, grow cyclic, every node/.style={concept, text=white, execute at begin node=\hskip0pt}, root concept/.append style={concept color=popblue, font=\Large\bfseries}, level 1/.append style={level distance=3.2cm, sibling angle=90}, concept color=popblue]
\node {México}
  child[concept color=popgreen] { node {Geografía} 
    child[concept color=popgreen!70, level distance=2.4cm] { node[align=center, font=\small] {situación\\ capital\\ población} } }
  child[concept color=poporange] { node {Historia} 
    child[concept color=poporange!70, level distance=2.4cm] { node[align=center, font=\small] {civilizaciones\\ 1810 / 1821\\ independencia} } }
  child[concept color=poppurple] { node {Cultura} 
    child[concept color=poppurple!70, level distance=2.4cm] { node[align=center, font=\small] {lengua\\ gastronomía\\ Frida Kahlo} } }
  child[concept color=popgold] { node {Economía} 
    child[concept color=popgold!70, level distance=2.4cm] { node[align=center, font=\small] {peso\\ petróleo\\ turismo} } };
\end{tikzpicture}
\end{center}
\legende{Carte mentale de la présentation : quatre branches, trois mots-clés par branche — c'est tout ce qu'il faut sur votre fiche d'orateur.}
\end{popfigure}

\courssub{Dire les grands nombres en espagnol}

La présentation d'un pays exige de manier les chiffres avec aisance. Observons d'abord :

\esp{México tiene más de 125 millones de habitantes.} \esp{La superficie es de 1.964.375 kilómetros cuadrados.} Remarquez : le mot \esp{millón} prend un \textbf{de} devant un nom (\esp{dos millones de habitantes}), et le point sert de séparateur de milliers en espagnol, là où le français utilise une espace insécable.

\begin{propriete}[Los números grandes]
\begin{itemize}
\item \esp{cien / ciento} : cent ; \esp{quinientos millones} : cinq cents millions.
\item \esp{mil} : mille (invariable : \esp{tres mil habitantes}).
\item \esp{un millón (de + nom)} : un million ; pluriel \esp{millones} : \esp{130 millones de habitantes}.
\item \esp{mil millones} : un milliard (jamais \esp{billón} au sens français !).
\item \esp{un billón} : un million de millions, soit mille milliards (10\textsuperscript{12}).
\end{itemize}
\end{propriete}

\begin{center}
\begin{tabularx}{\linewidth}{|X|X|X|}
\hline
\rowcolor{popblueL} Français & Español & Remarque \\
\hline
1 000 & \esp{mil} & invariable \\
\hline
1 000 000 & \esp{un millón} & \esp{millón de + nom} \\
\hline
125 000 000 & \esp{125 millones} & \textbf{Attention :} \emph{millions} prend un \textbf{s} en français comme en espagnol (125 millions / 125 millones). \\
\hline
1 000 000 000 & \esp{mil millones} & un milliard (pas \esp{un billón}) \\
\hline
1 000 000 000 000 & \esp{un billón} & mille milliards (10\textsuperscript{12}) \\
\hline
\end{tabularx}
\end{center}

\begin{attention}[¡Cuidado con el billón!]
Le faux ami le plus dangereux des chiffres : \esp{billón} espagnol signifie \textbf{mille milliards} (un million de millions, 10\textsuperscript{12}), et non « billion » anglais (10\textsuperscript{9}). Pour dire « un milliard » (10\textsuperscript{9}), il faut absolument dire \esp{mil millones}. Autre piège : ne dites jamais \esp{un millón habitantes} ; il faut la préposition \esp{de} : \esp{un millón de habitantes}. Enfin, le séparateur des milliers est le point : \esp{2.500.000} = 2 500 000 ; la virgule sert pour les décimales : \esp{3,14}.
\end{attention}

\begin{exoresolu}[Dire les chiffres]
Énoncé : écrivez en toutes lettres en espagnol : a) 3 000 000 d'habitants ; b) un milliard de pesos ; c) 47 000 000 d'habitants (Espagne) ; d) mille milliards de dollars.
\tcblower
Solution : a) \esp{tres millones de habitantes} ; b) \esp{mil millones de pesos} (et non \esp{un billón} !) ; c) \esp{cuarenta y siete millones de habitantes} ; d) \esp{un billón de dólares}. Vérifiez la présence du \esp{de} après \esp{millón/millones} et du pluriel \esp{millones} dès que le nombre dépasse un.
\end{exoresolu}

\courssub{Structurer l'exposé : annoncer, enchaîner, conclure}

Un bon orateur guide son auditoire. Pour cela, il annonce chaque étape avec des connecteurs d'organisation.

\begin{methode}[Anunciar el plan y enlazar las partes]
\begin{enumerate}
\item \textbf{Introduction} : saluez et annoncez le sujet : \esp{Buenos días a todos. Hoy voy a presentarles...}
\item \textbf{Annonce du plan} : \esp{En primer lugar, hablaré de la geografía; después, de la historia; en tercer lugar, de la cultura; y por último, de la economía.}
\item \textbf{Transitions} entre les parties : \esp{Pasemos ahora a...} (passons maintenant à...), \esp{En el plano económico...} (sur le plan économique...), \esp{Otro aspecto importante es...} (un autre aspect important est...).
\item \textbf{Conclusion} : \esp{Para concluir... / En resumen...} puis remerciez : \esp{Muchas gracias por su atención. ¿Tienen preguntas?}
\end{enumerate}
\end{methode}

\begin{propriete}[Futur ou présent pour annoncer ?]
Pour annoncer le plan, on emploie le futur simple : \esp{hablaré de...}, \esp{veremos...}, ou la tournure \esp{voy a + infinitif} : \esp{voy a presentarles...}. Les deux sont corrects ; le futur simple est un peu plus soutenu, le futur proche plus naturel à l'oral.
\end{propriete}

\courssub{Conseils pratiques pour l'oral}

\begin{itemize}
\item \textbf{Débit} : parlez lentement et articulez. En espagnol, chaque voyelle se prononce : \esp{geografía} se dit \emph{gé-o-gra-fí-a}, sans manger les syllabes.
\item \textbf{Regard} : ne lisez pas vos notes. Regardez l'auditoire ; ne consultez votre fiche que pour les chiffres.
\item \textbf{Support visuel} : apportez une carte, quelques images (drapeau, monuments, plats typiques). Montrez-les au bon moment : \esp{Como pueden ver en el mapa...} (comme vous pouvez le voir sur la carte...).
\item \textbf{Fiche d'orateur} : uniquement des mots-clés et des chiffres, jamais de phrases complètes. La carte mentale ci-dessus est un modèle idéal.
\item \textbf{Introduction et conclusion} : apprenez-les par cœur ; ce sont les moments où le regard doit être le plus assuré.
\item \textbf{Prononciation} : soignez les sons difficiles pour un francophone : la \emph{j} (\esp{México} se prononce \emph{Méjico}), le \emph{rr} roulé (\esp{perro}, \esp{ferrocarril}), et l'accent tonique (\esp{economía}, \esp{histórico}, \esp{México}).
\end{itemize}

\begin{savaistu}[Guinea Ecuatorial, el puente hispano-africano]
La Guinée équatoriale est le seul pays africain dont l'espagnol est langue officielle (avec le français et le portugais). Ancienne colonie espagnole, indépendante depuis le 12 octobre 1968, elle partage une frontière avec le Cameroun : c'est un véritable pont entre notre pays et le monde hispanique. Sa capitale est Malabo (sur l'île de Bioko), sa monnaie est le franc CFA (zone CEMAC, comme au Cameroun !) et son économie repose principalement sur le pétrole et le gaz. Présenter la Guinée équatoriale à l'oral est un excellent choix : original, proche de nous, et toujours apprécié des examinateurs pour sa pertinence géopolitique.
\end{savaistu}

\courssub{Critères d'évaluation et entraînement}

\begin{definition}[¿Cómo se evalúa una presentación oral?]
\begin{enumerate}
\item \textbf{Exactitud de los hechos} : les faits sont exacts (capitale, dates, chiffres vérifiés).
\item \textbf{Corrección lingüística} : grammaire, conjugaison, prononciation, accents toniques.
\item \textbf{Fluidez} : débit régulier, peu d'hésitations, phrases complètes.
\item \textbf{Claridad de la estructura} : plan annoncé, transitions explicites, conclusion nette.
\item \textbf{Comunicación} : regard, gestuelle, utilisation pertinente du support visuel.
\end{enumerate}
\end{definition}

\begin{exoresolu}[Reconstruire une présentation]
Énoncé : remettez dans l'ordre les phrases suivantes d'une présentation du Pérou, puis traduisez la troisième phrase de la présentation reconstituée.
\begin{itemize}
\item[a)] \esp{Por último, su moneda es el sol y exporta cobre.}
\item[b)] \esp{Buenos días. Hoy voy a presentarles Perú.}
\item[c)] \esp{En primer lugar, está situado en América del Sur y su capital es Lima.}
\item[d)] \esp{Para concluir, Perú es un país rico en historia y cultura.}
\item[e)] \esp{En segundo lugar, fue el centro del imperio inca.}
\end{itemize}
\tcblower
Solution : ordre correct \textbf{b -- c -- e -- a -- d}.
Traduction de la 3\textsuperscript{e} phrase (phrase \textbf{e}) : « \emph{En second lieu, il fut le centre de l'empire inca.} » On reconnaît les marqueurs d'organisation : salutation et sujet (b), géographie (c), histoire (e), économie (a), conclusion (d).
\end{exoresolu}

\begin{exercice}[Prepara tu presentación]
Choisissez un pays hispanophone (Espagne, Argentine, Cuba, Guinée équatoriale, Chili, Colombie...). Remplissez une fiche en quatre volets (\esp{geografía, historia, cultura, economía}) avec au moins trois informations par volet, dont deux chiffres écrits en toutes lettres (ex. : \esp{cuarenta y siete millones}, \esp{mil millones}). Rédigez ensuite l'introduction et la conclusion complètes, puis présentez votre pays en deux minutes devant la classe avec une carte comme support visuel.
\end{exercice}

\begin{aretenir}[L'essentiel de la présentation orale]
Un plan en quatre volets (géographie, histoire, culture, économie), annoncé avec \esp{en primer lugar... después... por último} ; les expressions clés \esp{está situado en}, \esp{limita con}, \esp{cuenta con}, \esp{su moneda es}, \esp{es conocido por} ; les grands nombres avec \esp{millón de} et \esp{mil millones} pour « un milliard » ; une fiche de mots-clés, un débit lent, le regard vers l'auditoire, et une conclusion apprise par cœur : \esp{Muchas gracias por su atención.}
\end{aretenir}

\coursec{Production guidée et entraînement au Bac}

Après avoir appris à présenter oralement un pays hispanophone, il est temps de passer à la \cle{pratique} : c'est ici que tout le travail de la leçon se concrétise. Au baccalauréat, on ne vous demandera pas de réciter des règles, mais de \emph{produire} : commenter un texte, rédiger une lettre ouverte ou une pétition, traduire dans les deux sens, parler. Cette section vous entraîne exactement à ces tâches, pas à pas, avec des modèles corrigés.

\courssub{Commentaire modèle du texte support}

Reprenons la lettre ouverte étudiée dans la première section et rédigeons ensemble le commentaire attendu au Bac. La méthode est toujours la même : introduction (présentation du texte et problématique), développement (analyse des arguments et du ton), conclusion (bilan et ouverture).

\begin{methode}[Commenter un texte revendicatif en trois étapes]
\begin{enumerate}
\item \textbf{Introducción} : présenter le type de texte (\esp{carta abierta}), son auteur, son destinataire, son thème. Formuler une problématique simple : \esp{¿Cómo intenta el autor convencer a su destinatario?}
\item \textbf{Desarrollo} : analyser d'abord les \emph{arguments} (faits, chiffres, exemples, valeurs invoquées), puis le \emph{ton} (formules de politesse, subjonctifs d'exigence, connecteurs, registre soutenu ou émotionnel).
\item \textbf{Conclusión} : évaluer l'efficacité du texte et ouvrir sur une question plus large (les droits des travailleurs, le rôle des lettres ouvertes dans la démocratie).
\end{enumerate}
\end{methode}

\begin{exemplebox}[Commentaire modèle (extrait rédigé)]
Introducción : El texto es una carta abierta dirigida al director de una empresa por sus empleados, que reclaman mejores condiciones de trabajo. ¿Cómo consigue el autor que su reivindicación sea respetuosa pero firme?

Desarrollo : En primer lugar, el autor expone los hechos con objetividad: los salarios son bajos y los horarios, excesivos. Después, utiliza el subjuntivo para formular sus exigencias: «exigimos que se respeten nuestros derechos». El tono es cortés gracias a fórmulas como «le rogamos que», pero firme gracias a conectores como «por lo tanto» y «sin embargo».

Conclusión : La carta es eficaz porque equilibra la cortesía y la firmeza. Nos hace pensar en el papel de la palabra escrita en la defensa de los derechos.
\end{exemplebox}

Traduction de l'introduction : « Le texte est une lettre ouverte adressée au directeur d'une entreprise par ses employés, qui réclament de meilleures conditions de travail. Comment l'auteur parvient-il à ce que sa revendication soit respectueuse mais ferme ? »

\courssub{Production écrite guidée : une carta abierta au proviseur}

\textbf{Sujet} : Vous êtes délégué(e) de classe. Rédigez une lettre ouverte au proviseur de votre lycée pour réclamer une salle informatique (15 à 20 lignes).

\begin{methode}[Plan à suivre]
\begin{enumerate}
\item En-tête : lieu, date, destinataire (\esp{Al señor director del Liceo de Yaoundé}).
\item Salutation : \esp{Señor director :}
\item Introduction : qui vous êtes et pourquoi vous écrivez (\esp{Me dirijo a usted en nombre de los alumnos de Terminale para...}).
\item Arguments : 2 ou 3 arguments (importance de l'informatique pour les études, l'examen, l'avenir professionnel), reliés par des connecteurs.
\item Exigence polie : \esp{Le rogamos que considere nuestra petición} ; subjonctif après \esp{es necesario que}.
\item Clôture : \esp{Le agradecemos de antemano su atención. Atentamente,} + signature.
\end{enumerate}
\end{methode}

\begin{exoresolu}[Rédaction guidée : la carta abierta]
Énoncé : rédigez la lettre en suivant le plan (15-20 lignes).
\tcblower
\textbf{Solution détaillée (modèle)} :

Yaoundé, 15 de marzo

Al señor director del Liceo de Yaoundé

Señor director:

Me dirijo a usted en nombre de los alumnos de Terminale para solicitar la creación de una sala de informática en nuestro liceo.

En primer lugar, la informática es indispensable para nuestros estudios: necesitamos buscar información, redactar trabajos y preparar el bachillerato. Además, muchos alumnos no tienen ordenador en casa y no pueden practicar. Por último, es necesario que nuestro liceo se adapte al mundo moderno, porque la informática es esencial para el mundo del trabajo.

Por lo tanto, le pedimos que instale una sala con veinte ordenadores y conexión a Internet. Le rogamos que estudie nuestra petición con atención.

Le agradecemos de antemano su comprensión.

Atentamente,

Amina Mbarga, delegada de Terminale A

\emph{Points de méthode} : vouvoiement constant (\esp{usted}, \esp{le}), subjonctif après \esp{es necesario que} (\esp{se adapte}) et après \esp{le pedimos que} (\esp{instale}), connecteurs variés (\esp{en primer lugar}, \esp{además}, \esp{por último}, \esp{por lo tanto}), formules d'ouverture et de clôture respectées.
\end{exoresolu}

\courssub{Production écrite : la petición pour un club d'espagnol}

\textbf{Sujet} : Rédigez une pétition demandant la création d'un club d'espagnol dans votre lycée.

La \cle{petición} est plus courte et plus collective que la lettre ouverte : un titre, un exposé des motifs (\esp{exposición de motivos}), la demande précise, puis les signatures.

\begin{exemplebox}[Modèle de petición]
PETICIÓN: CREACIÓN DE UN CLUB DE ESPAÑOL

Los abajo firmantes, alumnos del Liceo de Douala, solicitamos la creación de un club de español.

Exposición de motivos: el español es la segunda lengua más hablada del mundo; nuestro país tiene relaciones con Guinea Ecuatorial, país hispanohablante; un club nos permitiría practicar la lengua oralmente y descubrir la cultura hispánica.

Pedimos que la dirección apruebe este proyecto y que un profesor acompañe al club.

Firmas: ...
\end{exemplebox}

\begin{attention}[Petición : les pièges]
N'écrivez pas \esp{los abajo firmantes} comme « les signataires ci-dessous » mot à mot dans une lettre classique : c'est une formule propre à la pétition. Attention aussi au subjonctif après \esp{pedimos que} : \esp{pedimos que la dirección apruebe} (et non \esp{aprueba}).
\end{attention}

\courssub{Thème : traduire des phrases de revendication}

\begin{exemplebox}[Modèles traduits]
\begin{itemize}
\item « Nous exigeons que le gouvernement agisse vite. » = \esp{Exigimos que el gobierno actúe rápidamente.}
\item « Il est nécessaire que l'école achète des ordinateurs. » = \esp{Es necesario que la escuela compre ordenadores.}
\item « Nous demandons que nos droits soient respectés. » = \esp{Pedimos que se respeten nuestros derechos.}
\item « Il faut que les élèves aient accès à Internet. » = \esp{Es preciso que los alumnos tengan acceso a Internet.}
\item « Nous souhaitons qu'un professeur dirige le club. » = \esp{Deseamos que un profesor dirija el club.}
\end{itemize}
\end{exemplebox}

\begin{center}
\begin{tabularx}{\linewidth}{|X|X|X|}
\hline
\rowcolor{popblueL} Français & Español & Remarque \\
\hline
exiger que + subjonctif & exigir que + subjuntivo & même construction \\
\hline
il est nécessaire que & es necesario que & tournure impersonnelle + subjonctif \\
\hline
il faut que & hay que / es preciso que & \esp{hay que} + infinitif (impersonnel) ; \esp{es preciso que} + subjonctif (sujet exprimé dans la subordonnée) \\
\hline
demander que & pedir que & \esp{pedir} se conjugue : pido, pides, pide... \\
\hline
souhaiter que & desear que & subjonctif obligatoire \\
\hline
\end{tabularx}
\end{center}

\begin{attention}[Le piège du subjonctif]
En français, le subjonctif est souvent invisible (« qu'il agisse » ressemble à l'indicatif). En espagnol, il est toujours marqué : \esp{actúe}, \esp{compre}, \esp{tengan}, \esp{dirija}. Avant de traduire, demandez-vous : la principale exprime-t-elle une exigence, un souhait ou une nécessité ? Si oui, subjonctif obligatoire dans la subordonnée.
\end{attention}

\courssub{Version : extrait de lettre ouverte}

\begin{textebox}[Version à traduire]
Señor alcalde:

Los vecinos del barrio le escribimos para denunciar el mal estado de las calles. Desde hace dos años, las carreteras están llenas de agujeros y los niños no pueden ir a la escuela sin peligro cuando llueve.

Es urgente que el ayuntamiento repare las calles antes de la próxima temporada de lluvias. También pedimos que se instale alumbrado público, porque la oscuridad favorece la inseguridad.

Confiamos en que usted escuchará nuestra voz y actuará con rapidez.

Atentamente,

La asociación de vecinos
\end{textebox}

Glossaire : \esp{vecinos} = voisins ; \esp{agujeros} = trous ; \esp{ayuntamiento} = mairie ; \esp{alumbrado público} = éclairage public ; \esp{confiamos en que} = nous sommes convaincus que.

\begin{exoresolu}[Version corrigée]
Énoncé : traduisez le texte ci-dessus en français.
\tcblower
\textbf{Solution détaillée} :

« Monsieur le maire,

Les habitants du quartier vous écrivent pour dénoncer le mauvais état des rues. Depuis deux ans, les routes sont pleines de trous et les enfants ne peuvent pas aller à l'école sans danger quand il pleut.

Il est urgent que la mairie répare les rues avant la prochaine saison des pluies. Nous demandons aussi que soit installé un éclairage public, car l'obscurité favorise l'insécurité.

Nous sommes convaincus que vous écouterez notre voix et que vous agirez rapidement.

Cordialement,

L'association des habitants. »

\emph{Points de méthode} : \esp{es urgente que... repare} → subjonctif espagnol = subjonctif français (« répare ») ; \esp{se instale} → tournure passive en français (« soit installé ») ; \esp{confiamos en que... actuará} → indicatif futur après une expression de certitude.
\end{exoresolu}

\begin{attention}[En version : le mode d'abord]
Avant de traduire un verbe espagnol, repérez son \cle{mode}. Si c'est un subjonctif (\esp{repare}, \esp{instale}, \esp{escuche}), cherchez dans la principale l'expression qui le justifie (exigence, souhait, émotion) et rendez-le par un subjonctif français. Traduire \esp{pedimos que se instale} par « nous demandons qu'on installe » (indicatif) est une erreur de grammaire sanctionnée au Bac.
\end{attention}

\courssub{Préparation orale : fiche de présentation d'un pays}

Pour l'oral, préparez une fiche par pays. Modèle à compléter pour le pays de votre choix (Mexique, Pérou, Chili, Guinée équatoriale...) :

\begin{center}
\begin{tabularx}{\linewidth}{|l|X|}
\hline
\rowcolor{popblueL} Rubrique & Contenu à préparer \\
\hline
\esp{Nombre y capital} & nom officiel, capitale, situation géographique \\
\hline
\esp{Geografía} & relief, climat, pays voisins \\
\hline
\esp{Historia} & deux ou trois dates clés (indépendance, époque coloniale) \\
\hline
\esp{Cultura} & langues, fêtes, gastronomie, musique, personnalités \\
\hline
\esp{Economía} & ressources principales (pétrole, cuivre, tourisme, cacao...) \\
\hline
\esp{Conclusión} & pourquoi ce pays vous intéresse \\
\hline
\end{tabularx}
\end{center}

\begin{experience}[Simulation d'oral]
Par groupes de trois : un élève présente son pays en deux minutes à partir de sa fiche, les deux autres posent chacun une question (\esp{¿Cuál es la capital?}, \esp{¿Qué productos exporta?}). Puis on échange les rôles. Critères : fluidité, exactitude des faits, prononciation, regard vers l'auditoire.
\end{experience}

\courssub{De la rédaction à la copie parfaite : les étapes}

Toute production écrite réussie suit le même chemin. Ne rédigez jamais directement au propre.

\begin{popfigure}
\begin{tikzpicture}[node distance=0.6cm]
\node[draw=popblue, fill=popblueL, rounded corners, minimum width=2.6cm, minimum height=1cm, align=center] (a) {\textbf{1. Brouillon}\\ \esp{ideas y léxico}};
\node[draw=poporange, fill=poporangeL, rounded corners, minimum width=2.6cm, minimum height=1cm, align=center, right=of a] (b) {\textbf{2. Plan}\\ \esp{estructura}};
\node[draw=popgreen, fill=popgreenL, rounded corners, minimum width=2.6cm, minimum height=1cm, align=center, right=of b] (c) {\textbf{3. Redacción}\\ \esp{escritura}};
\node[draw=poppurple, fill=poppurpleL, rounded corners, minimum width=2.6cm, minimum height=1cm, align=center, right=of c] (d) {\textbf{4. Relectura}\\ \esp{corrección}};
\draw[-{Stealth}, very thick, popdark] (a) -- (b);
\draw[-{Stealth}, very thick, popdark] (b) -- (c);
\draw[-{Stealth}, very thick, popdark] (c) -- (d);
\end{tikzpicture}
\legende{Les quatre étapes de la production écrite : du brouillon à la relecture.}
\end{popfigure}

\begin{methode}[Checklist de relecture avant de rendre la copie]
\begin{enumerate}
\item \textbf{Subjonctif} : après \esp{exigir que}, \esp{pedir que}, \esp{es necesario que}, \esp{es urgente que} — le verbe de la subordonnée est-il bien au subjonctif ?
\item \textbf{Formules} : ouverture (\esp{Me dirijo a usted...}, \esp{Señor director :}) et clôture (\esp{Atentamente}, \esp{Le agradecemos de antemano}) présentes et correctes ?
\item \textbf{Connecteurs} : au moins trois connecteurs variés (\esp{en primer lugar}, \esp{además}, \esp{sin embargo}, \esp{por lo tanto}) ?
\item \textbf{Vouvoiement} : \esp{usted} et \esp{le} partout, jamais de \esp{tú} dans une lettre officielle ?
\item \textbf{Accents} : \esp{información}, \esp{petición}, \esp{además}, \esp{rápidamente} — chaque accent vérifié ?
\end{enumerate}
\end{methode}

\begin{exercice}[Entraînement Bac]
\begin{enumerate}
\item Traduisez en espagnol : « Il est urgent que le proviseur réponde à notre lettre. »
\item Traduisez en espagnol : « Nous demandons que le club se réunisse chaque semaine. »
\item Rédigez la clôture d'une lettre ouverte au maire de Douala (2 lignes).
\item Préparez oralement la fiche du pays de votre choix et présentez-le en deux minutes.
\end{enumerate}
\end{exercice}

\begin{corrige}[Exercice d'entraînement]
\begin{enumerate}
\item \esp{Es urgente que el director responda a nuestra carta.} (subjonctif \esp{responda} après \esp{es urgente que}).
\item \esp{Pedimos que el club se reúna cada semana.} (subjonctif \esp{se reúna}, verbe \esp{reunirse}).
\item \esp{Le agradecemos de antemano su atención y esperamos su respuesta. Atentamente,} + nom et qualité.
\item Oral : vérifier la présence des six rubriques de la fiche et la correction des faits.
\end{enumerate}
\end{corrige}

\begin{savaistu}[Comment est évaluée votre production au Bac ?]
Au baccalauréat, la production écrite est notée sur trois critères : la \emph{pertinence du contenu} (vous répondez au sujet, vos arguments sont logiques), la \emph{correction grammaticale} (subjonctifs, accords, conjugaisons) et la \emph{richesse du lexique} (vocabulaire varié, connecteurs, formules idiomatiques). Une copie simple mais correcte et bien structurée obtient toujours une meilleure note qu'une copie ambitieuse pleine de fautes. Relisez-vous : c'est là que se gagnent les derniers points !
\end{savaistu}

\newpage

\coursec{Activité d'intégration}

\courssub{Objectifs de l'activité}

À l'issue de cette activité d'intégration, tu seras capable de :

\begin{itemize}
\item rédiger, en groupe, une \cle{carta abierta} ou une \cle{petición} adressée à une autorité scolaire, en respectant la structure et les formules étudiées ;
\item argumenter une revendication avec au moins \cle{trois arguments} hiérarchisés et deux \cle{subjonctifs d'exigence} correctement employés ;
\item préparer et présenter oralement, en trois minutes, un \cle{pays hispanophone} (géographie, histoire, culture, économie) avec un support visuel ;
\item évaluer la production des autres groupes à l'aide d'une \cle{grille d'évaluation} simple et bienveillante.
\end{itemize}

Cette activité mobilise les quatre compétences du programme : compréhension de l'écrit, expression écrite argumentative, expression orale informative et interaction orale.

\courssub{Situation}

Tu es élève de Terminale A au Lycée bilingue de Yaoundé. Ton établissement accueille, à la fin du trimestre, une délégation d'élèves et de professeurs d'espagnol venus de Malabo (Guinée équatoriale) dans le cadre d'un échange scolaire. Pour préparer cette visite, le club d'espagnol lance le projet \esp{« Nuestra voz cuenta »} (« Notre voix compte »), qui comporte deux volets :

\begin{enumerate}
\item les élèves rédigent une lettre ouverte ou une pétition adressée au proviseur pour demander une amélioration concrète du lycée ; les meilleures lettres seront affichées au tableau d'annonces du club et remises officiellement à l'administration ;
\item chaque groupe prépare un exposé oral sur un pays hispanophone, qui sera présenté devant les invités équato-guinéens lors de la journée culturelle.
\end{enumerate}

Voici l'appel à contributions publié par le club d'espagnol :

\begin{textebox}[Appel à contributions du club d'espagnol]
\textbf{PROYECTO « NUESTRA VOZ CUENTA »}\par
\medskip
Queridos compañeros:\par
\medskip
Con motivo de la visita de nuestros amigos de Malabo, el club de español os invita a participar en dos actividades.\par
\medskip
Primero, redactad por grupos una \emph{carta abierta} o una \emph{petición} dirigida al señor director para pedir una mejora concreta de nuestro liceo: el comedor, la biblioteca, el campo de deportes, las aulas... Vuestra carta debe contener al menos tres argumentos sólidos y respetar las fórmulas de cortesía.\par
\medskip
Segundo, preparad una presentación oral de tres minutos sobre un país hispanohablante: su geografía, su historia, su cultura y su economía. Podéis utilizar carteles, mapas o dibujos.\par
\medskip
Las mejores cartas se publicarán en el tablón de anuncios y las presentaciones se harán el día de la fiesta cultural.\par
\medskip
¡Vuestra voz cuenta! ¡Participad!\par
\medskip
\emph{El equipo del club de español}
\end{textebox}

\textbf{Glossaire :} \esp{con motivo de} : à l'occasion de ; \esp{el comedor} : la cantine ; \esp{las aulas} : les salles de classe ; \esp{el tablón de anuncios} : le tableau d'affichage ; \esp{redactad} : rédigez (impératif, 2\up{e} pers. du pluriel) ; \esp{las fórmulas de cortesía} : les formules de politesse.

\courssub{Consignes}

\textbf{Tâche 1 — Compréhension et repérage (15 min).} Lis l'appel à contributions en groupe. Relevez : a) les deux productions attendues ; b) les destinataires de chaque production ; c) les contraintes de forme (nombre d'arguments, durée de l'exposé). Soulignez ensuite dans le texte deux formules de politesse et un impératif collectif.

\textbf{Tâche 2 — Choix du sujet et remue-méninges (10 min).} Choisissez \emph{un seul} problème concret de votre lycée : la cantine, la bibliothèque, le terrain de sport, les toilettes, l'éclairage des salles... Notez au brouillon trois arguments : un argument de \cle{santé ou sécurité}, un argument de \cle{réussite scolaire}, un argument de \cle{vie collective}. Classez-les du plus faible au plus fort.

\textbf{Tâche 3 — Rédaction de la carta abierta ou de la petición (35 min).} Rédigez votre texte en respectant la structure suivante : en-tête (lieu, date, destinataire), formule d'appel, introduction (qui vous êtes et pourquoi vous écrivez), développement (trois arguments avec connecteurs), conclusion (demande précise et formule de remerciement), signature collective. Votre texte doit contenir \cle{au moins deux subjonctifs d'exigence} (par exemple après \esp{exigimos que}, \esp{pedimos que}, \esp{es necesario que}) et \cle{un connecteur de chaque type} : addition, cause, conséquence.

\textbf{Tâche 4 — Préparation de l'exposé de pays (25 min).} Le professeur attribue à chaque groupe un pays hispanophone (par exemple : España, México, Argentina, Colombia, Perú, Cuba, Guinea Ecuatorial, Chile). Répartissez les rôles : un élève présente la \cle{géographie} (situation, capitale, population, relief), un autre l'\cle{histoire} (un ou deux événements majeurs), un troisième la \cle{culture} (langue, musique, gastronomie, fêtes) et le quatrième l'\cle{économie} (ressources, monnaie, activités principales). Fabriquez un support visuel simple : carte dessinée, drapeau, schéma.

\textbf{Tâche 5 — Passation et évaluation par les pairs (30 min).} Chaque groupe lit sa lettre à voix haute puis présente son exposé en trois minutes maximum. Les autres groupes remplissent la grille d'évaluation ci-dessous et formulent un point fort et un conseil d'amélioration.

\begin{center}
\begin{tabularx}{\linewidth}{|X|l|l|l|}\hline
\rowcolor{popblueL} \textbf{Critère} & \textbf{Bien (2)} & \textbf{Moyen (1)} & \textbf{À revoir (0)} \\ \hline
Contenu : arguments clairs et pertinents & & & \\ \hline
Langue : subjonctifs, connecteurs, formules & & & \\ \hline
Fluidité : débit, prononciation, regard & & & \\ \hline
Respect du temps et du support visuel & & & \\ \hline
\end{tabularx}
\end{center}

\courssub{Ressources à mobiliser}

\begin{aretenir}[Boîte à outils de la revendication]
\begin{itemize}
\item \textbf{Formules d'appel :} \esp{Señor director:} / \esp{Estimado señor director:}
\item \textbf{Présentation du groupe :} \esp{Los alumnos de Terminale A, abajo firmantes, ...} (« les élèves de Terminale A, soussignés »)
\item \textbf{Exigence + subjonctif :} \esp{exigimos que / pedimos que / solicitamos que / es necesario que / es urgente que} + \emph{subjuntivo}
\item \textbf{Connecteurs :} \esp{en primer lugar, además, por otra parte} (addition) ; \esp{porque, ya que, puesto que} (cause) ; \esp{por lo tanto, por eso, en consecuencia} (conséquence)
\item \textbf{Conclusion :} \esp{Por todo ello, esperamos que nuestra petición sea escuchada.} \esp{Agradecemos de antemano su atención.}
\item \textbf{Signature :} \esp{Atentamente, / Le saludan atentamente,} + noms du groupe.
\end{itemize}
\end{aretenir}

\begin{propriete}[Rappel : le subjonctif après les verbes d'exigence]
Après les verbes et expressions de volonté, de demande ou de nécessité suivis de \esp{que} (\esp{pedir, exigir, solicitar, querer, es necesario, es urgente, es importante}), le verbe de la subordonnée se met au \cle{subjonctif} lorsque le sujet des deux verbes est différent : \esp{Pedimos que el director \textbf{mejore} el comedor.} (« Nous demandons que le proviseur améliore la cantine. »). Si le sujet est le même, on utilise l'infinitif : \esp{Queremos \textbf{mejorar} el campo de deportes.} (« Nous voulons améliorer le terrain de sport. »). Rappel des formes du présent de subjonctif : \esp{mejore, mejores, mejore, mejoremos, mejoréis, mejoren} ; \esp{construya, construyas, construya, construyamos, construyáis, construyan} ; \esp{sea, seas, sea, seamos, seáis, sean}.
\end{propriete}

\begin{aretenir}[Lexique de la présentation d'un pays]
\esp{está situado en... / limita con... / su capital es... / tiene una población de... habitantes / la moneda es... / los principales recursos son... / se cultiva... / la lengua oficial es... / una fiesta típica es... / la gastronomía incluye...}. Verbes utiles : \esp{exportar, producir, celebrar, hablar, cultivar}.
\end{aretenir}

\begin{attention}[Pièges à éviter]
\begin{itemize}
\item Ne dis pas \esp{exigimos que el director \emph{mejora}} : après un verbe d'exigence, le subjonctif est obligatoire : \esp{exigimos que el director \emph{mejore} el comedor}.
\item Faux ami : \esp{una petición} signifie « une demande, une pétition », pas « une petite question ». De même, \esp{exigir} signifie « exiger, réclamer », et non « exiger » au sens de « nécessiter ».
\item Attention à la formule \esp{abajo firmantes} (« soussignés »), invariable dans l'usage administratif, et à la formule de politesse finale : on écrit \esp{Atentamente} (un seul mot) ; ne calque pas « Veuillez agréer... », qui n'a pas d'équivalent littéral en espagnol scolaire.
\item Noms de lieux : en espagnol, Yaoundé s'écrit \esp{Yaundé} et un lycée se dit \esp{un liceo}. Évite le mélange \esp{Lycée bilingüe}.
\item Pour les chiffres de population, reste prudent : dis \esp{unos veinte millones de habitantes} (« environ vingt millions ») plutôt qu'un chiffre précis non vérifié.
\end{itemize}
\end{attention}

\courssub{Proposition de production et corrigé commenté}

\begin{exoresolu}[Modèle de carta abierta et plan d'exposé]
Rédige la carta abierta du groupe 3, qui revendique la rénovation du terrain de sport du lycée, et donne le plan détaillé de l'exposé sur la Guinée équatoriale.
\tcblower
\textbf{1. Carta abierta (modèle)}\par
\medskip
\esp{Yaundé, 15 de marzo}\par
\medskip
\esp{Señor director del Liceo bilingüe de Yaundé:}\par
\medskip
\esp{Los alumnos de Terminale A, abajo firmantes, le escribimos para pedir la renovación del campo de deportes de nuestro liceo.}\par
\medskip
\esp{En primer lugar, el terreno está lleno de piedras y de agujeros, y varios alumnos han resultado heridos durante los partidos de fútbol. Además, cuando llueve, el campo se convierte en barro y es imposible practicar deporte durante semanas. Por otra parte, el deporte es esencial para nuestra salud y nuestra concentración en clase: un alumno sano aprende mejor.}\par
\medskip
\esp{Por eso, pedimos que se nivele el terreno y que se instalen porterías nuevas. También exigimos que se construya una pequeña tribuna para los espectadores. Creemos que es urgente que el liceo invierta en el deporte, porque nuestra salud y nuestro orgullo escolar están en juego.}\par
\medskip
\esp{Por todo ello, esperamos que nuestra petición sea escuchada y le agradecemos de antemano su atención.}\par
\medskip
\esp{Le saludan atentamente,}\par
\esp{El grupo 3 de Terminale A: Amina, Boris, Carine y Dieudonné.}
\medskip

\textbf{Traduction :} « Yaoundé, le 15 mars. Monsieur le Proviseur du Lycée bilingue de Yaoundé : Les élèves de Terminale A, soussignés, vous écrivons pour demander la rénovation du terrain de sport de notre lycée. D'abord, le terrain est plein de pierres et de trous, et plusieurs élèves ont été blessés pendant les matchs de football. De plus, quand il pleut, le terrain se transforme en boue et il est impossible de pratiquer un sport pendant des semaines. Par ailleurs, le sport est essentiel pour notre santé et notre concentration en classe : un élève en bonne santé apprend mieux. C'est pourquoi nous demandons que le terrain soit nivelé et que de nouvelles cages de but soient installées. Nous exigeons aussi qu'une petite tribune soit construite pour les spectateurs. Nous pensons qu'il est urgent que le lycée investisse dans le sport, car notre santé et notre fierté scolaire sont en jeu. Pour toutes ces raisons, nous espérons que notre demande sera entendue et nous vous remercions par avance de votre attention. Salutations distinguées, le groupe 3 de Terminale A. »
\medskip

\textbf{Commentaire des choix de langue :}
\begin{itemize}
\item \textbf{Structure respectée :} lieu et date, appel (\esp{Señor director}), présentation des signataires, trois arguments, demande précise, remerciement, formule de sortie.
\item \textbf{Subjonctifs d'exigence :} \esp{pedimos que se nivele}, \esp{exigimos que se construya}, \esp{es urgente que el liceo invierta} : le \emph{se} passif + subjonctif (\esp{se nivele, se instalen, se construya}) est une tournure élégante et impersonnelle, idéale pour une revendication polie mais ferme ; le verbe s'accorde avec le nom qui suit (\esp{se instalen porterías}, pluriel).
\item \textbf{Connecteurs variés :} \esp{en primer lugar, además, por otra parte} (addition) ; \esp{porque} (cause) ; \esp{por eso, por todo ello} (conséquence et bilan).
\item \textbf{ser / estar :} \esp{el terreno está lleno} (état résultant, \esp{estar} + participe) ; \esp{el deporte es esencial} (caractéristique, \esp{ser}).
\item \textbf{Ton :} ferme (\esp{exigimos}) mais courtois (\esp{le agradecemos de antemano}) : c'est l'équilibre attendu d'une lettre ouverte scolaire.
\end{itemize}
\medskip

\textbf{2. Plan d'exposé : Guinea Ecuatorial}\par
\medskip
\esp{Geografía: Guinea Ecuatorial está situada en África Central; tiene una parte continental, Río Muni, y varias islas, como Bioko, donde está la capital, Malabo. Limita con Camerún y Gabón.}\par
\esp{Historia: fue colonia española y accedió a la independencia en 1968; es el único país africano donde el español es lengua oficial.}\par
\esp{Cultura: conviven las lenguas española, francesa y portuguesa con lenguas locales como el fang y el bubi; la música y la danza ocupan un lugar central en las fiestas.}\par
\esp{Economía: el petróleo es el principal recurso; también se cultiva cacao y café. La moneda es el franco CFA, como en Camerún.}\par
\medskip
\textbf{Commentaire :} chaque partie est brève (trois ou quatre phrases), factuelle et prudente (pas de chiffres invérifiables). Le lien avec le Cameroun (\esp{limita con Camerún}, \esp{como en Camerún}) capte l'attention du public : c'est une excellente stratégie orale. Note l'emploi correct de \esp{donde} dans \esp{Bioko, donde está la capital} et l'accord de \esp{situada} avec \esp{Guinea Ecuatorial} (nom de pays féminin avec article : \esp{la Guinea Ecuatorial} est aussi admis).
\end{exoresolu}

\courssub{Bilan de l'activité}

Cette activité t'a permis de mettre en pratique, dans une tâche authentique et motivante, l'ensemble des acquis de la leçon :

\begin{itemize}
\item tu as \cle{compris} un appel à contributions rédigé en espagnol authentique ;
\item tu as \cle{rédigé} une lettre ouverte argumentée, en respectant la structure, les formules de politesse et la grammaire de l'exigence (subjonctif, connecteurs) ;
\item tu as \cle{présenté oralement} un pays hispanophone en un temps limité, avec un support visuel, ce qui entraîne directement à l'épreuve orale du Baccalauréat ;
\item tu as \cle{évalué} tes camarades avec une grille objective, ce qui développe ton esprit critique et ta coopération.
\end{itemize}

\begin{methode}[Pour aller plus loin]
\begin{enumerate}
\item Garde ta lettre dans ton portfolio : elle constitue un modèle réutilisable à l'examen (sujet de production écrite de type revendicatif).
\item Enregistre-toi en train de présenter ton exposé et écoute-toi : repère deux points à améliorer (prononciation, débit, regard).
\item Pour le Bac, entraîne-toi à rédiger une pétition en 25 minutes maximum, avec chronomètre, en suivant toujours le même plan : appel, présentation, arguments, demande, remerciement, signature.
\end{enumerate}
\end{methode}

\newpage

\coursec{Méthodes et exercices résolus}

\courssub{Texte support : une lettre ouverte revendicative}

\begin{textebox}[Carta abierta de los estudiantes de un liceo de Yaoundé]
Yaundé, 3 de marzo de 2025

Señor Ministro de Educación:

Nos dirigimos a usted, en nombre de los estudiantes de los liceos de la capital, para denunciar la falta de profesores de ciencias en muchos establecimientos públicos.

En primer lugar, desde hace dos años, numerosas clases no tienen profesor de física ni de matemáticas. Además, los laboratorios carecen de material y los alumnos no pueden hacer ningún experimento. Sin embargo, el bachillerato exige conocimientos científicos sólidos.

Por lo tanto, exigimos que el gobierno contrate más profesores y que repare los laboratorios antes del próximo año escolar. No pedimos un favor : reclamamos nuestro derecho a una educación de calidad.

Esperando una respuesta favorable, le saludan atentamente,

Los delegados de las clases de Terminale.
\end{textebox}

\textbf{Glossaire.} \esp{Se dirigir a usted} : s'adresser à vous (formel) ; \esp{denunciar} : dénoncer ; \esp{la falta de} : le manque de ; \esp{desde hace dos años} : depuis deux ans (durée qui continue) ; \esp{carecer de} : manquer de ; \esp{ningún experimento} : aucune expérience (scientifique) ; \esp{contratar} : embaucher, recruter ; \esp{reclamar} : réclamer ; \esp{un favor} : une faveur ; \esp{el derecho a} : le droit à.

\textbf{Questions de repérage.} 1. Quelle est la nature du texte ? 2. Qui écrit, à qui, pour quoi ? 3. Relevez deux verbes de revendication. 4. Relevez trois connecteurs d'argumentation. Ces questions préparent le commentaire (méthode 5).

\courssub{Lexique thématique : revendication, droits et monde du travail}

\begin{definition}[Vocabulaire de la revendication et du travail]
\begin{center}
\begin{tabularx}{\linewidth}{|X|X|}\hline
\rowcolor{popblueL} \textbf{Español} & \textbf{Français} \\ \hline
la carta abierta & la lettre ouverte \\ \hline
la petición / la reivindicación & la pétition / la revendication \\ \hline
exigir / pedir / reclamar & exiger / demander / réclamer \\ \hline
denunciar & dénoncer \\ \hline
los derechos humanos & les droits de l'homme \\ \hline
el derecho a la educación & le droit à l'éducation \\ \hline
el trabajador / el obrero & le travailleur / l'ouvrier \\ \hline
el salario (mínimo) & le salaire (minimum) \\ \hline
la jornada laboral & la journée de travail \\ \hline
el aumento de sueldo & l'augmentation de salaire \\ \hline
la huelga & la grève \\ \hline
el sindicato / el portavoz & le syndicat / le porte-parole \\ \hline
el desempleo / el empleo & le chômage / l'emploi \\ \hline
contratar / despedir & embaucher / licencier \\ \hline
las condiciones de trabajo & les conditions de travail \\ \hline
la justicia / la injusticia & la justice / l'injustice \\ \hline
firmar una petición & signer une pétition \\ \hline
\end{tabularx}
\end{center}
\end{definition}

\begin{attention}[Faux amis du monde du travail]
\begin{itemize}
\item \esp{la jornada} = la journée (de travail), jamais « le voyage » (\esp{el viaje}).
\item \esp{una demanda} = une plainte en justice ou la demande (économique) ; une revendication sociale se dit \esp{una reivindicación} ou \esp{una petición}.
\item \esp{asistir a} = être présent à, assister à ; « aider » se dit \esp{ayudar}.
\item \esp{actualmente} = actuellement, en ce moment ; « en fait, en réalité » se dit \esp{en realidad}.
\item \esp{un descanso} = une pause, un repos (pas « un mécontentement »).
\end{itemize}
\end{attention}

\courssub{Grammaire de l'argumentation : exprimer l'exigence et l'opinion}

\begin{propriete}[Verbes de volonté + subjonctif]
Après les verbes qui expriment la \cle{volonté}, l'\cle{exigence} ou la \cle{nécessité} suivis de \esp{que}, le verbe subordonné est au \textbf{subjonctif} : \esp{exigir que}, \esp{pedir que}, \esp{querer que}, \esp{es necesario que}, \esp{es importante que}.

\esp{Exigimos que el gobierno actúe.} « Nous exigeons que le gouvernement agisse. »

\esp{Es necesario que haya más profesores.} « Il faut qu'il y ait plus de professeurs. »

\textbf{Exception.} Si le sujet est le même dans les deux propositions, on utilise l'\textbf{infinitif} sans \esp{que} : \esp{Queremos trabajar.} « Nous voulons travailler. » (et non \esp{queremos que trabajemos}). De même, devant un nom, pas de subjonctif : \esp{Exigimos la reapertura.} « Nous exigeons la réouverture. »
\end{propriete}

\begin{propriete}[Subjonctif présent : rappel de formation]
Radical de la 1\iere{} personne du présent de l'indicatif + terminaisons inversées :
\begin{center}
\begin{tabular}{|l|l|l|l|}\hline
\rowcolor{popblueL} & \textbf{hablar} & \textbf{comer} & \textbf{abrir} \\ \hline
yo & hable & coma & abra \\ \hline
tú & hables & comas & abras \\ \hline
él/usted & hable & coma & abra \\ \hline
nosotros & hablemos & comamos & abramos \\ \hline
vosotros & habléis & comáis & abráis \\ \hline
ellos/ustedes & hablen & coman & abran \\ \hline
\end{tabular}
\end{center}
Irréguliers fréquents : \esp{ser} $\rightarrow$ sea ; \esp{estar} $\rightarrow$ esté ; \esp{haber} $\rightarrow$ haya ; \esp{ir} $\rightarrow$ vaya ; \esp{tener} $\rightarrow$ tenga ; \esp{hacer} $\rightarrow$ haga ; \esp{poder} $\rightarrow$ pueda ; \esp{crear} $\rightarrow$ cree.
\end{propriete}

\begin{propriete}[Verbes d'opinion : indicatif ou subjonctif]
\begin{itemize}
\item Opinion \textbf{affirmée} (certitude) + \textbf{indicatif} : \esp{Creo que la situación es grave.} « Je crois que la situation est grave. »
\item Opinion \textbf{niée} (doute) + \textbf{subjonctif} : \esp{No creo que sea justo.} « Je ne crois pas que ce soit juste. »
\end{itemize}
Même logique avec \esp{pensar que}, \esp{opinar que}, \esp{es cierto que} / \esp{no es cierto que}.
\end{propriete}

\begin{popfigure}
\begin{tikzpicture}[node distance=0.55cm and 1.1cm]
\node[draw=popblue, fill=popblueL, rounded corners, align=center, font=\small\bfseries] (q) {Verbe principal + \esp{que}};
\node[draw=popgreen, fill=popgreenL, rounded corners, align=center, font=\small, below left=of q] (subj) {Volonté, exigence,\\ nécessité, opinion niée\\ \textbf{SUBJUNTIVO}};
\node[draw=poporange, fill=poporangeL, rounded corners, align=center, font=\small, below right=of q] (ind) {Certitude, fait,\\ opinion affirmée\\ \textbf{INDICATIVO}};
\draw[-{Stealth}, thick, popgreen] (q) -- (subj);
\draw[-{Stealth}, thick, poporange] (q) -- (ind);
\end{tikzpicture}
\legende{Choisir le mode après \esp{que} : souhait ou doute $\rightarrow$ subjonctif ; certitude $\rightarrow$ indicatif.}
\end{popfigure}

\begin{exoresolu}[Thème : phrases à traduire]
\textbf{Énoncé.} Traduisez en espagnol : 1. Nous demandons que le directeur ouvre la bibliothèque. 2. Je crois que les élèves ont raison. 3. Je ne pense pas que cette solution soit suffisante. 4. Il faut que le gouvernement crée des emplois pour les jeunes.
\tcblower
\textbf{Raisonnement.} Phrases 1 et 4 : volonté et nécessité $\rightarrow$ subjonctif. Phrase 2 : opinion affirmative $\rightarrow$ indicatif. Phrase 3 : opinion niée $\rightarrow$ subjonctif.

\textbf{Solutions.}
\begin{enumerate}
\item \esp{Pedimos que el director abra la biblioteca.} — \esp{abra} : subjonctif présent d'\esp{abrir} après \esp{pedir que} ; sujets différents, donc \esp{que} + subjonctif.
\item \esp{Creo que los alumnos tienen razón.} — indicatif après \esp{creer que} affirmatif ; attention : \esp{tener razón} (jamais de calque avec \esp{haber}).
\item \esp{No pienso que esta solución sea suficiente.} — \esp{sea} : subjonctif de \esp{ser} après opinion niée.
\item \esp{Es necesario que el gobierno cree empleos para los jóvenes.} — \esp{cree} : subjonctif de \esp{crear} ; \esp{para} exprime le bénéficiaire (« pour les jeunes »).
\end{enumerate}
\textbf{Conclusion.} Le choix du mode dépend du verbe principal : volonté et nécessité appellent le subjonctif, certitude l'indicatif.
\end{exoresolu}

\courssub{Méthode 1 : rédiger une lettre ouverte}

\begin{methode}[Rédiger une lettre ouverte revendicative]
Une \cle{carta abierta} est une lettre adressée à une personnalité ou à une institution, mais destinée à être lue par le public (presse, réseaux sociaux). Démarche en six étapes :
\begin{enumerate}
\item \textbf{Identifier} le destinataire officiel et le public visé : cela détermine le registre (formel, \esp{usted}).
\item \textbf{Rédiger l'en-tête} : lieu et date à droite (\esp{Yaundé, 15 de marzo de 2025}), puis la formule d'appel : \esp{Señor Ministro:}, \esp{Estimado Director:}, \esp{A la atención de...}. En espagnol, la formule d'appel est suivie de \textbf{deux-points}.
\item \textbf{Introduire} : se présenter brièvement et annoncer l'objet de la lettre : \esp{Me dirijo a usted en nombre de los estudiantes para denunciar...}.
\item \textbf{Développer} en 2 ou 3 paragraphes : constat (\esp{La situación es insostenible}), arguments avec connecteurs (\esp{en primer lugar, además, por último}), exemples concrets.
\item \textbf{Formuler la revendication} clairement : \esp{Exigimos que...} + subjonctif, \esp{Solicitamos la construcción de...}, \esp{Pedimos medidas urgentes}.
\item \textbf{Conclure et signer} : \esp{Esperando una respuesta favorable, le saluda atentamente,} + nom, prénom, qualité (\esp{delegado de la clase de Terminale A}).
\end{enumerate}
Réflexes : un paragraphe = une idée ; toujours \esp{usted} (jamais \esp{tú}) ; ton ferme mais courtois ; vérifier les accents avant de rendre la copie.
\end{methode}

\begin{exoresolu}[Lettre ouverte au proviseur]
\textbf{Énoncé.} Vous êtes délégué(e) de classe. La bibliothèque de votre lycée à Douala est fermée depuis six mois. Rédigez une lettre ouverte au proviseur (80--100 mots) pour demander sa réouverture, en publiant cette lettre dans le journal du lycée.
\tcblower
\textbf{Raisonnement.} Destinataire officiel : le proviseur (\esp{el director}) ; registre formel, \esp{usted}. Plan : appel, présentation + objet, constat, arguments (2), revendication, formule de politesse.

\textbf{Réponse rédigée.}
\esp{Douala, 10 de febrero de 2025}

\esp{Señor Director:}

\esp{Me dirijo a usted en nombre de los alumnos de Terminale para denunciar el cierre de la biblioteca del liceo desde hace seis meses.}

\esp{En primer lugar, sin biblioteca, los alumnos no pueden consultar libros ni preparar el bachillerato en buenas condiciones. Además, muchos estudiantes no tienen libros en casa y la biblioteca era su único recurso.}

\esp{Por eso, exigimos la reapertura inmediata de la biblioteca y solicitamos la compra de libros nuevos.}

\esp{Esperando una respuesta favorable, le saluda atentamente,}

\esp{Marie Ngo Bell, delegada de la clase de Terminale A.}

\textbf{Justifications.} \esp{Me dirijo a usted} : formule d'introduction formelle ; \esp{desde hace seis meses} : présent + \esp{desde hace} pour une durée qui continue ; \esp{exigimos la reapertura} : nom après \esp{exigir} (pas de subjonctif sans \esp{que}) ; \esp{le saluda} : pronom \esp{le} pour \esp{usted} (complément indirect).

\textbf{Conclusion.} La lettre respecte les six étapes : environ 95 mots, ton ferme et courtois, revendication explicite.
\end{exoresolu}

\courssub{Méthode 2 : rédiger une pétition}

\begin{methode}[Rédiger une pétition]
Une \cle{petición} est un texte court, collectif, qui demande des signatures. Structure à retenir :
\begin{enumerate}
\item \textbf{Titre} en majuscules : \esp{PETICIÓN : POR UNA CARRETERA ASFALTADA EN NUESTRO BARRIO}.
\item \textbf{Destinataire} : \esp{Dirigida a: el Alcalde de la ciudad de Yaundé}.
\item \textbf{Exposé du problème} en 3--5 lignes : faits, chiffres prudents, conséquences : \esp{Los habitantes del barrio sufren cada día...}.
\item \textbf{Revendication unique et précise} introduite par \esp{Por lo tanto, pedimos que...} + subjonctif : \esp{pedimos que se asfalte la carretera principal}.
\item \textbf{Appel à signer} : \esp{¡Firma esta petición y compártela!}, \esp{Tu firma puede cambiar las cosas.}
\item \textbf{Tableau des signatures} : nom, prénom, signature, date.
\end{enumerate}
Réflexes : phrases courtes et percutantes ; une seule demande ; utiliser \esp{nosotros} (la communauté) ; subjonctif obligatoire après \esp{pedir que}, \esp{exigir que}.
\end{methode}

\begin{exoresolu}[Pétition pour l'eau potable]
\textbf{Énoncé.} Dans votre quartier de Yaoundé, les coupures d'eau sont quotidiennes. Rédigez une pétition adressée au maire (60--80 mots) avec un titre, un exposé du problème, une revendication et un appel à signer.
\tcblower
\textbf{Raisonnement.} Pétition = texte collectif (\esp{nosotros}), une seule revendication, subjonctif après \esp{pedimos que}.

\textbf{Réponse rédigée.}
\esp{PETICIÓN : POR AGUA POTABLE EN NUESTRO BARRIO}

\esp{Dirigida al Señor Alcalde de Yaundé.}

\esp{Los habitantes del barrio sufren cortes de agua todos los días. Los niños no pueden lavarse antes de ir al colegio y las familias compran agua cara en el mercado. Esta situación es peligrosa para la salud de todos.}

\esp{Por lo tanto, pedimos que se reparen las tuberías y que el barrio reciba agua potable cada día.}

\esp{¡Firma esta petición! Tu firma puede cambiar las cosas.}

\textbf{Justifications.} \esp{pedimos que se reparen} : subjonctif présent après \esp{pedir que} ; \esp{se reparen} : tournure passive réfléchie ; \esp{reciba} : subjonctif dans la seconde demande coordonnée ; \esp{Firma} : impératif affirmatif de \esp{firmar} (tutoiement du lecteur-citoyen, acceptable dans une pétition).

\textbf{Conclusion.} Texte de 75 mots environ, structure complète, revendication unique et réaliste.
\end{exoresolu}

\courssub{Méthode 3 : présenter oralement un pays hispanophone}

\begin{methode}[Présenter un pays en 5 parties]
Pour un exposé oral de 3--5 minutes, suivre un plan fixe :
\begin{enumerate}
\item \textbf{Presentación} : \esp{Voy a presentar Guinea Ecuatorial / España / México...} Nommer la capitale et la situation : \esp{Es un país situado en...}.
\item \textbf{Geografía} : superficie, relief, climat, pays voisins : \esp{Limita al norte con..., tiene un clima tropical...}.
\item \textbf{Historia} : 2 ou 3 dates clés au prétérit simple : \esp{Fue colonia española hasta 1968, año de su independencia.}
\item \textbf{Cultura} : langues, musique, gastronomie, fêtes : \esp{Se habla español, francés y portugués ; la música típica es...}.
\item \textbf{Economía} : ressources et activités : \esp{La economía depende del petróleo, del cacao y de la pesca.} Conclure : \esp{En conclusión, es un país rico en cultura...}.
\end{enumerate}
Réflexes oraux : parler lentement, regarder l'auditoire, utiliser le présent de l'indicatif pour les faits actuels et le prétérit simple pour les dates historiques, préparer une fiche avec des mots-clés (jamais un texte lu), donner les chiffres approximatifs sans les inventer avec précision.
\end{methode}

\begin{savaistu}[La Guinée équatoriale, voisine hispanophone du Cameroun]
La Guinée équatoriale est le seul pays d'Afrique subsaharienne où l'espagnol est langue officielle. Ancienne colonie espagnole indépendante depuis 1968, elle partage une frontière avec le Cameroun (région du Sud). Pour un élève camerounais, parler espagnol, c'est pouvoir communiquer avec un voisin direct : un atout pour le commerce, les études et la coopération régionale. C'est pourquoi l'espagnol est une langue stratégique au Cameroun, avec le français et l'anglais.
\end{savaistu}

\begin{exoresolu}[Questions sur une présentation]
\textbf{Énoncé.} Voici un extrait d'exposé : \esp{Guinea Ecuatorial es el único país africano donde el español es lengua oficial. Su capital es Malabo, situada en la isla de Bioko. El país fue colonia española hasta 1968. Su economía depende sobre todo del petróleo, pero el cacao y el café también son importantes.} Répondez en espagnol : 1. ¿Cuál es la capital de Guinea Ecuatorial? 2. ¿Hasta cuándo fue colonia española? 3. ¿De qué depende su economía? 4. Pourquoi ce pays est-il important pour un élève camerounais ?
\tcblower
\textbf{Solutions.}
\begin{enumerate}
\item \esp{La capital de Guinea Ecuatorial es Malabo, situada en la isla de Bioko.} — reprise exacte du texte, phrase complète.
\item \esp{Fue colonia española hasta 1968, año de su independencia.} — \esp{fue} : prétérit simple de \esp{ser} (fait historique ponctuel et daté).
\item \esp{Su economía depende sobre todo del petróleo, pero el cacao y el café también son importantes.} — \esp{depender de} : préposition \esp{de} obligatoire.
\item \esp{Es importante porque es un país vecino del Camerún y porque el español es su lengua oficial : hablar español facilita los intercambios económicos y culturales entre los dos países.} — argument personnel correct, \esp{porque} + indicatif.
\end{enumerate}
\textbf{Conclusion.} À l'oral, répondre par des phrases complètes en reprenant les mots de la question : c'est la règle d'or de la compréhension orale au Bac.
\end{exoresolu}

\courssub{Méthode 4 : commenter un texte revendicatif}

\begin{methode}[Lecture et commentaire d'un texte revendicatif]
\begin{enumerate}
\item \textbf{Lire deux fois} le texte ; souligner les marques de la revendication : verbes (\esp{exigir, denunciar, reclamar}), champs lexicaux (droits, travail, injustice).
\item \textbf{Présenter} le document : nature (lettre ouverte, article, pétition), auteur, destinataire, thème : \esp{El texto es una carta abierta escrita por... que denuncia...}.
\item \textbf{Dégager la thèse} : que demande l'auteur ? \esp{El autor exige que...}.
\item \textbf{Analyser les procédés} : ton (ferme, ironique), pronoms (\esp{nosotros} = union du groupe), champs lexicaux, connecteurs.
\item \textbf{Donner son avis} argumenté : \esp{En mi opinión, el autor tiene razón porque...} / \esp{Sin embargo, creo que...}.
\end{enumerate}
Réflexe : toujours justifier une réponse par une courte citation du texte entre guillemets.
\end{methode}

\begin{exoresolu}[Compréhension d'un texte revendicatif]
\textbf{Énoncé.} Texte : \esp{Los trabajadores de la fábrica de cacao de Douala denuncian los salarios bajos y las jornadas de doce horas. « No pedimos caridad, exigimos nuestros derechos », declara su portavoz. Los obreros quieren un aumento del salario mínimo y descansos regulares.} Questions : 1. ¿Qué denuncian los trabajadores? 2. ¿Qué significa « No pedimos caridad, exigimos nuestros derechos »? 3. ¿Qué quieren los obreros? 4. Donnez votre avis en deux phrases.
\tcblower
\textbf{Solutions.}
\begin{enumerate}
\item \esp{Denuncian los salarios bajos y las jornadas de doce horas.} — citation fidèle du texte ; \esp{jornada} = journée de travail (faux ami : pas « voyage »).
\item \esp{Significa que los trabajadores no quieren regalos : reclaman lo que les corresponde legalmente, es decir, salarios justos.} — reformulation de l'opposition \esp{pedir caridad} / \esp{exigir derechos}.
\item \esp{Quieren un aumento del salario mínimo y descansos regulares.} — \esp{aumento} = augmentation (nom d'\esp{aumentar}).
\item \esp{En mi opinión, los trabajadores tienen razón porque doce horas de trabajo es demasiado. Además, un salario justo permite vivir dignamente.} — opinion + justification ; \esp{tener razón} sans article.
\end{enumerate}
\textbf{Conclusion.} Chaque réponse s'appuie sur le texte ; l'avis personnel reste bref, argumenté et en espagnol correct.
\end{exoresolu}

\courssub{Production guidée et entraînement au Bac}

\begin{experience}[Débat et jeu de rôle : la réunion de crise]
Par groupes de quatre : deux élèves jouent les délégués des élèves, un joue le proviseur, un joue le journaliste du journal du lycée. Situation : la cantine est fermée depuis un mois. Les délégués présentent leurs revendications (\esp{exigimos que..., pedimos que...}), le proviseur répond (\esp{entiendo su petición, sin embargo...}), le journaliste pose des questions (\esp{¿Qué denuncian ustedes? ¿Qué solución proponen?}). Puis chaque groupe rédige la pétition correspondante en dix minutes et la présente à la classe.
\end{experience}

\begin{exercice}[Entraînement au Bac : thème et rédaction]
\textbf{A. Thème.} Traduisez en espagnol : 1. Les ouvriers exigent que l'entreprise augmente les salaires. 2. Nous ne croyons pas que cette décision soit juste. 3. Il est important que les jeunes trouvent du travail. 4. Je pense que le ministre va répondre à notre lettre.

\textbf{B. Rédaction.} Votre quartier de Douala n'a pas d'éclairage public. Rédigez une pétition adressée au maire (70--90 mots) : titre, exposé du problème, revendication unique, appel à signer.
\end{exercice}

\begin{corrige}[Entraînement au Bac]
\textbf{A. Thème.}
\begin{enumerate}
\item \esp{Los obreros exigen que la empresa aumente los salarios.} — \esp{aumente} : subjonctif après \esp{exigir que}.
\item \esp{No creemos que esta decisión sea justa.} — opinion niée $\rightarrow$ subjonctif \esp{sea}.
\item \esp{Es importante que los jóvenes encuentren trabajo.} — \esp{encuentren} : subjonctif d'\esp{encontrar} (radical \esp{encuentr-}) ; \esp{trabajo} sans article.
\item \esp{Pienso que el ministro va a responder a nuestra carta.} — opinion affirmative $\rightarrow$ indicatif ; futur proche \esp{va a responder} ; \esp{responder a}.
\end{enumerate}
\textbf{B. Rédaction (modèle).}
\esp{PETICIÓN : POR LUZ ELÉCTRICA EN LAS CALLES DE NUESTRO BARRIO}

\esp{Dirigida al Señor Alcalde de Douala.}

\esp{Los habitantes del barrio viven sin luz en las calles desde hace un año. Por la noche, los alumnos no pueden estudiar fuera y los vecinos tienen miedo de salir. Esta situación es peligrosa para la seguridad de todos.}

\esp{Por lo tanto, pedimos que se instalen farolas en las calles principales del barrio.}

\esp{¡Firma esta petición! Juntos podemos cambiar nuestro barrio.}

Points de méthode : titre en majuscules, une seule revendication, \esp{pedimos que se instalen} (subjonctif + passive réfléchie), appel à signer à l'impératif.
\end{corrige}

\begin{attention}[Les 5 erreurs qui coûtent des points au Bac]
\begin{enumerate}
\item \textbf{Oublier le subjonctif} après les verbes de volonté : écrire \esp{exigimos que el gobierno \emph{actúa}} au lieu de \esp{actúe}. Toute revendication avec \esp{que} exige le subjonctif.
\item \textbf{Faux amis} : \esp{asistir a} signifie « être présent à » (pas « aider » = \esp{ayudar}) ; \esp{una demanda} est une plainte en justice, une revendication se dit \esp{una reivindicación} ou \esp{una petición} ; \esp{actualmente} = « actuellement, en ce moment », et non « en fait » (\esp{en realidad}).
\item \textbf{Calques du français} : « avoir raison » = \esp{tener razón} (jamais \esp{haber razón}) ; « être d'accord » = \esp{estar de acuerdo} ; « demander à quelqu'un de faire » = \esp{pedir a alguien que} + subjonctif, sans infinitif calqué du français.
\item \textbf{Accents oubliés} : \esp{petición, situación, opinión, económico, geografía, política, América, México}. Un accent oublié change parfois le temps : \esp{hablo} (je parle) / \esp{habló} (il parla).
\item \textbf{Mauvais choix de temps} : pour un fait historique daté, prétérit simple : \esp{El país fue colonia hasta 1968} (pas d'imparfait) ; pour une durée qui continue, présent + \esp{desde hace} : \esp{La biblioteca está cerrada desde hace seis meses} (pas de passé composé calqué du français).
\end{enumerate}
\end{attention}

\begin{aretenir}[L'essentiel de la leçon]
\begin{itemize}
\item Lettre ouverte : en-tête, appel formel avec deux-points, \esp{usted}, arguments connectés, revendication explicite, formule de politesse.
\item Pétition : titre en majuscules, une seule revendication, \esp{nosotros}, appel à signer à l'impératif.
\item Grammaire clé : volonté, exigence, nécessité, opinion niée + \esp{que} $\rightarrow$ \textbf{subjonctif} ; certitude et opinion affirmée $\rightarrow$ \textbf{indicatif}.
\item Présentation d'un pays : 5 parties (présentation, géographie, histoire, culture, économie), présent pour les faits actuels, prétérit simple pour les dates.
\item La Guinée équatoriale, voisine du Cameroun, est le seul pays africain hispanophone officiel : un argument fort pour apprendre l'espagnol.
\end{itemize}
\end{aretenir}

\newpage

\coursec{Exercices}

\courssub{Je vérifie mes connaissances}

\begin{exercice}[QCM : la lettre ouverte]
Pour chaque question, choisis la bonne réponse.
\begin{enumerate}
\item Une \esp{carta abierta} s'adresse :
\begin{enumerate}
\item uniquement à un ami intime
\item à un destinataire précis mais est rendue publique
\item uniquement à un tribunal
\end{enumerate}
\item Pour exprimer une exigence en espagnol, on utilise :
\begin{enumerate}
\item \esp{exigir que} + subjonctif
\item \esp{exigir que} + indicatif
\item \esp{exigir de} + infinitif
\end{enumerate}
\item La formule correcte d'ouverture d'une lettre ouverte au ministre est :
\begin{enumerate}
\item \esp{Querido Ministro, ¿qué tal?}
\item \esp{Excelentísimo Señor Ministro:}
\item \esp{Hola Señor Ministro,}
\end{enumerate}
\item Dans une pétition, la formule \esp{Los abajo firmantes} signifie :
\begin{enumerate}
\item les signataires ci-dessous
\item les absents
\item les responsables
\end{enumerate}
\end{enumerate}
\end{exercice}

\begin{exercice}[Vrai ou faux ? Justifie]
Réponds par vrai ou faux et justifie ta réponse en une phrase.
\begin{enumerate}
\item Après \esp{es urgente que}, on emploie l'indicatif.
\item \esp{No creemos que sea posible} contient un subjonctif correctement employé.
\item Une pétition doit toujours commencer par une insulte pour être efficace.
\item La Guinée équatoriale est un pays hispanophone d'Afrique.
\end{enumerate}
\end{exercice}

\begin{exercice}[Vocabulaire de la revendication]
Complète chaque phrase avec le mot approprié de la liste : \esp{huelga, derechos, salario, manifestación, reclamar, sindicato}.
\begin{enumerate}
\item Los trabajadores están en \dots\ desde el lunes : no trabajan.
\item Los ciudadanos salieron a la calle en una gran \dots\ pacífica.
\item Los empleados \dots\ un aumento de \dots\ porque no ganan bastante.
\item El \dots\ defiende los intereses de los obreros.
\item Todos los seres humanos tienen \dots\ fundamentales.
\end{enumerate}
\end{exercice}

\begin{exercice}[Conjugaison : le subjonctif de l'exigence]
Complète avec le présent du subjonctif du verbe entre parenthèses.
\begin{enumerate}
\item Exigimos que el Gobierno \dots\ (construir) más escuelas.
\item Es urgente que el Estado \dots\ (respetar) nuestros derechos.
\item Pedimos que los profesores \dots\ (ser) mejor pagados.
\item No creemos que esta situación \dots\ (poder) continuar.
\item Queremos que todos los alumnos \dots\ (tener) acceso a la educación.
\item Es importante que ustedes \dots\ (venir) a la reunión.
\end{enumerate}
\end{exercice}

\courssub{Je m'entraîne}

\begin{exercice}[Thème : phrases à traduire]
Traduis en espagnol les phrases suivantes.
\begin{enumerate}
\item Il est urgent que l'État construise des écoles.
\item Nous ne croyons pas que ce soit impossible.
\item Les élèves demandent que leurs droits soient respectés.
\item Il faut que le gouvernement écoute les jeunes.
\end{enumerate}
\end{exercice}

\begin{exercice}[Transformations : de l'indicatif au subjonctif]
Transforme chaque phrase selon le modèle.\\[4pt]
\emph{Modèle :} \esp{El Gobierno ayuda a los jóvenes. (Es necesario que)} $\rightarrow$ \esp{Es necesario que el Gobierno ayude a los jóvenes.}
\begin{enumerate}
\item El Estado construye más liceos. (Es urgente que)
\item Los empresarios respetan a los trabajadores. (Exigimos que)
\item Hay más profesores de español. (No creemos que)
\item Todos los niños van a la escuela. (Es importante que)
\item Usted firma esta petición. (Le pedimos que)
\end{enumerate}
\end{exercice}

\begin{exercice}[Texte à trous : une petición]
Complète cette pétition avec les mots et formes de la liste : \esp{Excelentísimo, abajo, solicitamos, respeten, atentamente, construya}.

\begin{textebox}[Petición]
\dots\ Señor Gobernador:

Los \dots\ firmantes, habitantes del barrio de Nkolbisson, \dots\ que el Ayuntamiento \dots\ una escuela primaria en nuestro barrio. Es urgente que las autoridades \dots\ los derechos de nuestros hijos.

Le saludamos \dots.

Los vecinos de Nkolbisson
\end{textebox}
\end{exercice}

\begin{exercice}[Appariements : formules de la lettre ouverte]
Relie chaque formule espagnole à sa fonction.
\begin{center}
\begin{tabularx}{\linewidth}{|X|X|}
\hline
\rowcolor{popblueL} \textbf{Formule} & \textbf{Fonction} \\
\hline
\esp{Excelentísimo Señor Ministro:} & 1. Formule de politesse finale \\
\hline
\esp{Me dirijo a usted para expresarle nuestra preocupación.} & 2. En-tête / destinataire \\
\hline
\esp{Exigimos que se tomen medidas urgentes.} & 3. Annonce de l'objet de la lettre \\
\hline
\esp{Confiando en su comprensión, le saludamos atentamente.} & 4. Formulation de la revendication \\
\hline
\end{tabularx}
\end{center}
\end{exercice}

\courssub{Je me prépare au Bac}

\begin{exercice}[Compréhension de l'écrit et questions de langue]
Lis le texte suivant, puis réponds aux questions \textbf{en espagnol}, avec des phrases complètes.

\begin{textebox}[Carta abierta de los trabajadores de Douala]
A la opinión pública y a las autoridades del país:

Nosotros, los trabajadores del puerto de Douala, escribimos esta carta abierta para denunciar una situación que ya no podemos soportar. Trabajamos diez horas al día, bajo un sol terrible, pero nuestros salarios no nos permiten alimentar dignamente a nuestras familias.

No creemos que sea justo que los obreros más pobres paguen siempre los errores de los más ricos. Exigimos que la empresa aumente nuestros salarios y que respete nuestras condiciones de trabajo. Pedimos también que el Gobierno controle a las empresas que maltratan a sus empleados.

No somos enemigos de nadie : somos padres y madres que quieren un futuro mejor para sus hijos. Esperamos que nuestra voz sea escuchada antes de que sea demasiado tarde.

Los trabajadores del puerto de Douala
\end{textebox}

\textbf{Questions de compréhension}
\begin{enumerate}
\item ¿Quiénes escriben esta carta y a quiénes se dirige?
\item ¿Cuáles son las dos quejas principales de los trabajadores?
\item ¿Qué exigen los trabajadores a la empresa? ¿Y qué piden al Gobierno?
\item ¿Cuál es el tono de esta carta : violento, irónico o firme y digno? Justifica tu respuesta con una frase del texto.
\end{enumerate}

\textbf{Questions de langue}
\begin{enumerate}
\item Releva dans le texte deux verbes au subjonctif et explique pourquoi le subjonctif est employé.
\item Conjugue le verbe \esp{exigir} au présent de l'indicatif (toutes les personnes).
\item Trouve dans le texte un synonyme de \esp{protestar} et un contraire de \esp{amigos}.
\end{enumerate}
\end{exercice}

\begin{exercice}[Thème et version]
\textbf{Thème.} Traduis en espagnol (attention au subjonctif et à la voix passive) :
\begin{enumerate}
\item Les ouvriers exigent que leurs salaires soient augmentés.
\item Il est important que le gouvernement écoute les revendications des jeunes.
\item Nous ne pensons pas que cette situation soit normale.
\item Les habitants du quartier demandent qu'on construise un hôpital.
\end{enumerate}

\textbf{Version.} Traduis en français l'extrait suivant :

\begin{textebox}[Extracto de una carta abierta]
Señor Director:

Llevamos tres meses en huelga y nadie nos escucha. No pedimos la luna : pedimos simplemente que se respeten nuestros derechos más básicos. Es injusto que trabajemos sin contrato y sin seguridad social. Nuestros hijos nos miran y nos preguntan por qué luchamos. Les respondemos que luchamos por ellos, para que un día tengan un trabajo digno y una vida mejor. No nos rendiremos jamás.
\end{textebox}
\end{exercice}

\begin{exercice}[Production écrite et présentation orale]
\textbf{A. Production écrite (type Bac).} Tu es président(e) de l'association des élèves hispanistes de ton lycée. Rédige une \esp{carta abierta} (150 à 180 mots) au ministre de l'Éducation pour réclamer plus de professeurs d'espagnol dans les lycées du Cameroun. Tu respecteras la structure de la lettre ouverte : en-tête, formule d'appel, présentation du problème, arguments, revendications avec \esp{exigir que} / \esp{pedir que} + subjonctif, formule de politesse finale et signature.

\textbf{B. Présentation orale (3 minutes).} Prépare un exposé oral de trois minutes pour présenter \textbf{la Guinée équatoriale} ou \textbf{le Pérou}. Ton exposé doit comporter :
\begin{enumerate}
\item une introduction (\esp{Voy a presentarles...});
\item la géographie (situation, capitale, population);
\item un élément d'histoire (colonisation, indépendance);
\item un aspect culturel (langue, musique, gastronomie);
\item l'économie (ressources principales);
\item une conclusion personnelle (\esp{Me gustaría visitar este país porque...}).
\end{enumerate}
\end{exercice}

\newpage

\coursec{Corrigés des exercices}

\begin{corrige}[Exercice 1 — QCM : la lettre ouverte]
\begin{enumerate}
\item Réponse \textbf{b}. Une \esp{carta abierta} s'adresse à un destinataire précis (ministre, directeur, président) mais elle est \cle{rendue publique} (presse, affichage, réseaux sociaux) pour toucher l'opinion.
\item Réponse \textbf{a}. \esp{Exigir que} est un verbe de volonté : il est suivi du \cle{subjonctif}. \esp{Exigimos que el Gobierno actúe.}
\item Réponse \textbf{b}. La formule protocolaire correcte est \esp{Excelentísimo Señor Ministro:} (avec deux-points, jamais de virgule dans une lettre formelle). Les formules \esp{¿qué tal?} et \esp{Hola} relèvent du registre familier.
\item Réponse \textbf{a}. \esp{Los abajo firmantes} = « les soussignés », littéralement « les signataires ci-dessous ».
\end{enumerate}
\end{corrige}

\begin{corrige}[Exercice 2 — Vrai ou faux ? Justifie]
\begin{enumerate}
\item \textbf{Faux.} Après \esp{es urgente que}, on emploie le \cle{subjonctif}, car il s'agit d'une expression impersonnelle de nécessité : \esp{Es urgente que el Estado actúe.}
\item \textbf{Vrai.} Après \esp{no creer que} (opinion niée), le subjonctif est obligatoire : \esp{No creemos que sea posible} (\esp{sea} = subjonctif présent de \esp{ser}).
\item \textbf{Faux.} Une pétition efficace utilise un ton \cle{ferme mais respectueux} ; l'insulte dessert la cause et décrédibilise les signataires.
\item \textbf{Vrai.} La Guinée équatoriale, ancienne colonie espagnole, a l'espagnol pour langue officielle : c'est le seul pays hispanophone d'Afrique subsaharienne.
\end{enumerate}
\end{corrige}

\begin{corrige}[Exercice 3 — Vocabulaire de la revendication]
\begin{enumerate}
\item Los trabajadores están en \textbf{huelga} desde el lunes : no trabajan. (« Les travailleurs sont en grève depuis lundi. »)
\item Los ciudadanos salieron a la calle en una gran \textbf{manifestación} pacífica. (« Les citoyens sont descendus dans la rue dans une grande manifestation pacifique. »)
\item Los empleados \textbf{reclaman} un aumento de \textbf{salario} porque no ganan bastante. (« Les employés réclament une augmentation de salaire. »)
\item El \textbf{sindicato} defiende los intereses de los obreros. (« Le syndicat défend les intérêts des ouvriers. »)
\item Todos los seres humanos tienen \textbf{derechos} fundamentales. (« Tous les êtres humains ont des droits fondamentaux. »)
\end{enumerate}
\end{corrige}

\begin{attention}[Faux amis et pièges lexicaux]
\begin{itemize}
\item « Être en grève » = \esp{estar en huelga} (avec \esp{estar}) ; « faire grève » = \esp{hacer huelga}. Ne pas traduire « grève » par un faux ami : \esp{huelga} est correct.
\item \esp{Salario} et \esp{sueldo} = « salaire » ; attention, \esp{salario} est masculin : \esp{un aumento de salario}.
\item \esp{Reclamar} = « réclamer » (transitif direct, sans préposition) : \esp{reclaman un aumento}, jamais \esp{reclaman por un aumento}.
\item \esp{Fundamental} a la même forme au masculin et au féminin ; au pluriel : \esp{los derechos fundamentales}, \esp{las libertades fundamentales}.
\end{itemize}
\end{attention}


\begin{corrige}[Exercice 4 — Conjugaison : le subjonctif de l'exigence]
\begin{enumerate}
\item Exigimos que el Gobierno \textbf{construya} más escuelas. (radical \esp{constru-} + désinence \esp{-ya}, verbe en \esp{-uir})
\item Es urgente que el Estado \textbf{respete} nuestros derechos.
\item Pedimos que los profesores \textbf{sean} mejor pagados. (\esp{ser}, subjonctif irrégulier : \esp{sea, seas, sea, seamos, seáis, sean})
\item No creemos que esta situación \textbf{pueda} continuar. (\esp{poder} : changement de voyelle \esp{o $\rightarrow$ ue})
\item Queremos que todos los alumnos \textbf{tengan} acceso a la educación. (\esp{tener} : \esp{tenga, tengas, tenga, tengamos, tengáis, tengan})
\item Es importante que ustedes \textbf{vengan} a la reunión. (\esp{venir} : \esp{venga, vengas, venga, vengamos, vengáis, vengan})
\end{enumerate}
\end{corrige}

\begin{aretenir}[Règle]
Après \esp{exigir que, pedir que, querer que, es urgente que, es importante que} et \esp{no creer que}, le verbe de la subordonnée est au \cle{subjonctif présent} : radical de la 1\iere{} personne de l'indicatif + désinences inversées (\esp{-e/-es/-e...} pour \esp{-ar} ; \esp{-a/-as/-a...} pour \esp{-er/-ir}).
\end{aretenir}


\begin{corrige}[Exercice 5 — Thème : phrases à traduire]
\begin{enumerate}
\item \esp{Es urgente que el Estado construya escuelas.}
\item \esp{No creemos que sea imposible.} (ou \esp{No creemos que esto sea imposible.})
\item \esp{Los alumnos piden que sus derechos sean respetados.} (voix passive avec \esp{ser} + participe accordé : \esp{respetados})
\item \esp{Es necesario que el Gobierno escuche a los jóvenes.} (ou \esp{Hay que} + infinitif si le sujet n'est pas précisé ; ici sujet exprimé, donc \esp{que} + subjonctif)
\end{enumerate}
\end{corrige}

\begin{attention}[Points de vigilance]
\begin{itemize}
\item « L'État » = \esp{el Estado} (majuscule) ; « l'état » (condition) = \esp{el estado}.
\item Après \esp{escuchar} devant une personne, la préposition \esp{a} est obligatoire : \esp{escuchar a los jóvenes}.
\item Ne pas calquer « ce soit » par \esp{esto sea} systématiquement : \esp{sea} seul suffit souvent.
\end{itemize}
\end{attention}


\begin{corrige}[Exercice 6 — Transformations : de l'indicatif au subjonctif]
\begin{enumerate}
\item \esp{Es urgente que el Estado construya más liceos.}
\item \esp{Exigimos que los empresarios respeten a los trabajadores.}
\item \esp{No creemos que haya más profesores de español.} (\esp{haber} impersonnel : \esp{haya})
\item \esp{Es importante que todos los niños vayan a la escuela.} (\esp{ir} : \esp{vaya, vayas, vaya, vayamos, vayáis, vayan})
\item \esp{Le pedimos que usted firme esta petición.} (\esp{firmar} : \esp{firmes} pour \esp{tú}, \esp{firme} pour \esp{usted} — ici \esp{usted}, donc \esp{firme})
\end{enumerate}
\end{corrige}

\begin{attention}[Piège fréquent]
Le verbe \esp{ir} au subjonctif n'est pas \esp{ia} mais \esp{vaya}. De même, \esp{haber} impersonnel donne \esp{haya} (et non \esp{haba}).
\end{attention}


\begin{corrige}[Exercice 7 — Texte à trous : una petición]
Texte complété :

\esp{Excelentísimo Señor Gobernador:}

\esp{Los abajo firmantes, habitantes del barrio de Nkolbisson, solicitamos que el Ayuntamiento construya una escuela primaria en nuestro barrio. Es urgente que las autoridades respeten los derechos de nuestros hijos.}

\esp{Le saludamos atentamente.}

\esp{Los vecinos de Nkolbisson}

Traduction : « Excellence Monsieur le Gouverneur : Les soussignés, habitants du quartier de Nkolbisson, demandons que la mairie construise une école primaire dans notre quartier. Il est urgent que les autorités respectent les droits de nos enfants. Nous vous saluons respectueusement. Les habitants de Nkolbisson. »

Justifications : \esp{construya} et \esp{respeten} sont au subjonctif après \esp{solicitamos que} et \esp{es urgente que} ; \esp{Excelentísimo} s'accorde avec \esp{Señor Gobernador} (masculin singulier).
\end{corrige}

\begin{corrige}[Exercice 8 — Appariements : formules de la lettre ouverte]
\begin{center}
\begin{tabularx}{\linewidth}{|X|X|}
\hline
\rowcolor{popblueL} \textbf{Formule} & \textbf{Fonction} \\
\hline
\esp{Excelentísimo Señor Ministro:} & 2. En-tête / destinataire \\
\hline
\esp{Me dirijo a usted para expresarle nuestra preocupación.} & 3. Annonce de l'objet de la lettre \\
\hline
\esp{Exigimos que se tomen medidas urgentes.} & 4. Formulation de la revendication \\
\hline
\esp{Confiando en su comprensión, le saludamos atentamente.} & 1. Formule de politesse finale \\
\hline
\end{tabularx}
\end{center}
\end{corrige}

\begin{aretenir}[Ordre logique]
La lettre ouverte suit toujours l'ordre : en-tête $\rightarrow$ annonce de l'objet $\rightarrow$ arguments $\rightarrow$ revendication $\rightarrow$ politesse finale $\rightarrow$ signature.
\end{aretenir}


\begin{corrige}[Exercice 9 — Compréhension de l'écrit et questions de langue]
\textbf{Questions de compréhension} (réponses modèles en espagnol) :
\begin{enumerate}
\item \esp{Esta carta la escriben los trabajadores del puerto de Douala y se dirige a la opinión pública y a las autoridades del país.}
\item \esp{Las dos quejas principales son que trabajan diez horas al día bajo un sol terrible y que sus salarios no les permiten alimentar dignamente a sus familias.}
\item \esp{Exigen a la empresa que aumente sus salarios y que respete sus condiciones de trabajo. Piden al Gobierno que controle a las empresas que maltratan a sus empleados.}
\item \esp{El tono es firme y digno: los trabajadores dicen «No somos enemigos de nadie: somos padres y madres que quieren un futuro mejor para sus hijos», lo que demuestra que protestan con respeto y sin violencia.}
\end{enumerate}

\textbf{Questions de langue} :
\begin{enumerate}
\item Subjonctifs possibles : \esp{sea} (\esp{No creemos que sea justo} : subjonctif après une opinion niée) ; \esp{aumente} et \esp{respete} (\esp{Exigimos que...} : subjonctif après un verbe d'exigence) ; \esp{sea escuchada} (\esp{Esperamos que...} : subjonctif après un verbe d'espoir).
\item Présent de l'indicatif de \esp{exigir} : \esp{exijo, exiges, exige, exigimos, exigís, exigen}. Attention : changement \esp{g $\rightarrow$ j} à la 1\iere{} personne (\esp{exijo}) pour conserver le son.
\item Synonyme de \esp{protestar} : \esp{denunciar}. Contraire de \esp{amigos} : \esp{enemigos}.
\end{enumerate}
\textbf{Barème indicatif (type Bac)} : compréhension 8 pts (2 pts par question) ; questions de langue 6 pts (2 pts par question).
\end{corrige}

\begin{corrige}[Exercice 10 — Thème et version]
\textbf{Thème} (traductions proposées) :
\begin{enumerate}
\item \esp{Los obreros exigen que sus salarios sean aumentados.} (ou tournure pronominale : \esp{exigen que se aumenten sus salarios})
\item \esp{Es importante que el Gobierno escuche las reivindicaciones de los jóvenes.}
\item \esp{No pensamos que esta situación sea normal.}
\item \esp{Los habitantes del barrio piden que se construya un hospital.} (ou \esp{que construyan un hospital})
\end{enumerate}
Points de vigilance : « revendications » se dit \esp{reivindicaciones} (avec \esp{ei}) ; la voix passive espagnole exige l'accord du participe (\esp{sean aumentados}) ; après \esp{no pensar que}, subjonctif obligatoire.

\textbf{Version} (traduction proposée) :

« Monsieur le Directeur :

Cela fait trois mois que nous sommes en grève et personne ne nous écoute. Nous ne demandons pas la lune : nous demandons simplement que soient respectés nos droits les plus élémentaires. Il est injuste que nous travaillions sans contrat et sans sécurité sociale. Nos enfants nous regardent et nous demandent pourquoi nous luttons. Nous leur répondons que nous luttons pour eux, afin qu'un jour ils aient un travail digne et une vie meilleure. Nous ne nous rendrons jamais. »

Points de vigilance : \esp{Llevamos tres meses en huelga} = « cela fait trois mois que... » (présent + \esp{llevar}) ; \esp{no pedir la luna} = « ne pas demander la lune » (même image en français) ; \esp{para que tengan} = « afin qu'ils aient » (subjonctif de but) ; \esp{no rendirse} = « ne pas se rendre / ne pas abandonner ».

\textbf{Barème indicatif} : thème 8 pts (2 pts par phrase) ; version 8 pts (fidélité 4 pts, qualité du français 4 pts).
\end{corrige}

\begin{corrige}[Exercice 11 — Production écrite et présentation orale]
\textbf{A. Production écrite — proposition de rédaction modèle} :

\begin{textebox}[Carta abierta al Ministro de Educación]
Excelentísimo Señor Ministro de Educación:

Me dirijo a usted, en nombre de la asociación de los alumnos hispanistas de nuestro liceo, para expresarle nuestra profunda preocupación. En muchos liceos de Camerún, el español se enseña sin profesores suficientes: a veces una sola persona enseña a cientos de alumnos.

Sin embargo, el español es una lengua internacional hablada por millones de personas, y nuestro país tiene fronteras con Guinea Ecuatorial, país hispanohablante. Aprender español es una oportunidad económica y cultural para la juventud camerunesa.

Por eso, exigimos que el Ministerio contrate más profesores de español y pedimos que se creen centros de formación para los futuros docentes. Es urgente que cada liceo tenga al menos dos profesores de esta lengua.

Confiando en su comprensión, le saludamos atentamente.

El presidente de la asociación de hispanistas
\end{textebox}

Traduction : « Excellence Monsieur le Ministre de l'Éducation : Je m'adresse à vous, au nom de l'association des élèves hispanistes de notre lycée, pour vous exprimer notre profonde inquiétude. Dans beaucoup de lycées du Cameroun, l'espagnol est enseigné sans professeurs suffisants : parfois une seule personne enseigne à des centaines d'élèves. Pourtant, l'espagnol est une langue internationale parlée par des millions de personnes, et notre pays a des frontières avec la Guinée équatoriale, pays hispanophone. Apprendre l'espagnol est une opportunité économique et culturelle pour la jeunesse camerounaise. C'est pourquoi nous exigeons que le ministère recrute plus de professeurs d'espagnol et demandons que soient créés des centres de formation pour les futurs enseignants. Il est urgent que chaque lycée ait au moins deux professeurs de cette langue. Espérant votre compréhension, nous vous saluons respectueusement. Le président de l'association des hispanistes. »

\textbf{Barème indicatif (type Bac, 20 pts)} : respect de la structure de la lettre ouverte 4 pts ; correction grammaticale (subjonctifs, accords) 6 pts ; richesse du lexique de la revendication 4 pts ; cohérence et force de l'argumentation 4 pts ; présentation et longueur respectée 2 pts.

\textbf{Points de vigilance} : deux-points après la formule d'appel ; \esp{exigir que} / \esp{pedir que} + subjonctif (\esp{contrate, se creen, tenga}) ; en espagnol, « le Cameroun » se dit \esp{Camerún} avec accent (et \esp{de Camerún} sans article) ; « réclamer » = \esp{reclamar} (transitif direct, sans \esp{por}).

\textbf{B. Présentation orale — plan détaillé attendu} (exemple avec la Guinée équatoriale) :
\begin{enumerate}
\item Introduction : \esp{Voy a presentarles Guinea Ecuatorial, un país africano hispanohablante.}
\item Géographie : \esp{Está situada en África Central, en el golfo de Guinea. Su capital es Malabo, en la isla de Bioko, y su ciudad principal del continente es Bata. Tiene fronteras con Camerún y Gabón.}
\item Histoire : \esp{Fue una colonia española y se independizó en 1968.}
\item Culture : \esp{El español es lengua oficial, junto con el francés y el portugués. Su música y sus danzas tradicionales son muy ricas.}
\item Économie : \esp{Su economía depende sobre todo del petróleo y de la madera.}
\item Conclusion : \esp{Me gustaría visitar este país porque es el único país hispanohablante de África subsahariana y quiero conocer su cultura.}
\end{enumerate}
\textbf{Barème indicatif (oral, 10 pts)} : respect du plan 2 pts ; correction linguistique 3 pts ; fluidité et prononciation 3 pts ; durée respectée et aisance 2 pts.
\end{corrige}

\newpage

\coursec{Fiche bilan}

\begin{popfigure}
\begin{tikzpicture}[mindmap, grow cyclic, every node/.style={concept, align=center, text width=2.4cm, font=\small},
root concept/.append style={concept color=popblue, text=white, font=\bfseries\small, text width=3cm},
level 1 concept/.append style={level distance=4.4cm, sibling angle=51.5},
level 2 concept/.append style={level distance=2.8cm}]
\node[root concept]{Carta abierta, petición y presentación de un país}
child[concept color=poporange]{ node[align=center]{Texto reivindicativo\\ \footnotesize lectura y comentario}
child[concept color=poporange!60]{ node[text width=2cm]{\footnotesize tesis + argumentos}}
child[concept color=poporange!60]{ node[text width=2cm]{\footnotesize tono y destinatario}}
}
child[concept color=popgreen]{ node[align=center]{Léxico\\ \footnotesize derechos y trabajo}
child[concept color=popgreen!60]{ node[text width=2cm]{\footnotesize la huelga, el salario}}
child[concept color=popgreen!60]{ node[text width=2cm]{\footnotesize exigir, reclamar}}
}
child[concept color=popgold]{ node[align=center]{Gramática\\ \footnotesize exigencia y opinión}
child[concept color=popgold!60]{ node[text width=2cm]{\footnotesize subjuntivo tras exigir}}
child[concept color=popgold!60]{ node[text width=2cm]{\footnotesize es necesario que...}}
}
child[concept color=poppink]{ node[align=center]{Carta abierta\\ \footnotesize estructura}
child[concept color=poppink!60]{ node[text width=2cm]{\footnotesize encabezado, cuerpo, firma}}
}
child[concept color=poppurple]{ node[align=center]{Petición\\ \footnotesize fórmulas}
child[concept color=poppurple!60]{ node[text width=2cm]{\footnotesize solicitamos, pedimos}}
}
child[concept color=popblue!70]{ node[align=center]{País hispanohablante\\ \footnotesize exposición oral}
child[concept color=popblue!40]{ node[text width=2cm]{\footnotesize geografía, historia}}
child[concept color=popblue!40]{ node[text width=2cm]{\footnotesize cultura, economía}}
};
\end{tikzpicture}
\legende{Carte mentale de la leçon : lettre ouverte, pétition et présentation d'un pays hispanophone}
\end{popfigure}

\begin{bilan}
\textbf{1. Lexique clé : revendication, droits et monde du travail}
\begin{itemize}
\item \esp{la carta abierta} : la lettre ouverte ; \esp{la petición / la solicitud} : la pétition / la demande
\item \esp{reclamar / exigir / reivindicar} : réclamer / exiger / revendiquer ; \esp{la reivindicación} : la revendication
\item \esp{los derechos} : les droits ; \esp{el derecho al trabajo / a la educación} : le droit au travail / à l'éducation
\item \esp{el salario / el sueldo} : le salaire ; \esp{el aumento salarial} : l'augmentation de salaire
\item \esp{la huelga} : la grève ; \esp{el sindicato} : le syndicat ; \esp{los trabajadores} : les travailleurs
\item \esp{las condiciones de trabajo} : les conditions de travail ; \esp{el desempleo / el paro} : le chômage
\item \esp{la igualdad / la desigualdad} : l'égalité / l'inégalité ; \esp{la justicia social} : la justice sociale
\item \esp{dirigirse a las autoridades} : s'adresser aux autorités ; \esp{firmar una petición} : signer une pétition
\end{itemize}

\textbf{2. Grammaire à connaître par c\oe ur}
\begin{center}
\begin{tabularx}{\linewidth}{|X|X|X|}
\hline
\rowcolor{popblueL} Fonction & Español & Français \\
\hline
Exigence + subjonctif & \esp{Exigimos que se respeten los derechos.} & Nous exigeons que les droits soient respectés. \\
\hline
Demande + subjonctif & \esp{Le pedimos que actúe con rapidez.} & Nous vous demandons d'agir vite. \\
\hline
Opinion affirmative + indicatif & \esp{Creemos que la situación es injusta.} & Nous croyons que la situation est injuste. \\
\hline
Opinion niée + subjonctif & \esp{No creemos que sea suficiente.} & Nous ne croyons pas que ce soit suffisant. \\
\hline
Nécessité impersonnelle + subjonctif & \esp{Es necesario que haya un diálogo.} & Il faut qu'il y ait un dialogue. \\
\hline
But : \esp{para que} + subjonctif & \esp{Firmamos para que nos escuchen.} & Nous signons pour qu'on nous écoute. \\
\hline
\end{tabularx}
\end{center}
Verbes déclencheurs du subjonctif : \esp{exigir, pedir, solicitar, querer, desear, esperar, prohibir} + \esp{que}. Après \esp{creer / pensar / opinar} affirmatifs : indicatif ; niés : subjonctif.

\textbf{3. Structure de la lettre ouverte}
\begin{enumerate}
\item Lieu et date : \esp{Yaoundé, 15 de mayo de 2025}
\item Destinataire : \esp{Al señor Ministro de Trabajo} / \esp{A la opinión pública}
\item Titre annonçant la revendication : \esp{Carta abierta por un salario digno}
\item Corps : présentation du problème, arguments, exigences précises
\item Formule finale : \esp{Esperamos una respuesta favorable. Atentamente,}
\item Signature : nom, qualité, collectif
\end{enumerate}

\textbf{4. Structure de la pétition} : titre clair, exposé des faits (\esp{Considerando que...}), demande formulée (\esp{Solicitamos que...}), liste de signatures.

\textbf{5. Présentation orale d'un pays} : saluer, annoncer le plan (\esp{Voy a hablar de... En primer lugar...}), géographie (situation, capitale, population), histoire (indépendance, repères), culture (langue, fêtes, gastronomie), économie (ressources, monnaie), conclusion personnelle. Connecteurs : \esp{además, sin embargo, por último, en conclusión}.
\end{bilan}

\begin{aretenir}[Checklist avant le Bac]
\begin{itemize}
\item[$\square$] Je connais le lexique de la revendication et du monde du travail (\esp{huelga, salario, derechos, exigir}).
\item[$\square$] Je sais employer le subjonctif après \esp{exigir que, pedir que, es necesario que}.
\item[$\square$] Je distingue opinion affirmative (indicatif) et opinion niée (subjonctif).
\item[$\square$] Je maîtrise la structure complète d'une lettre ouverte (destinataire, corps, formule finale, signature).
\item[$\square$] Je sais rédiger une pétition avec \esp{Considerando que...} et \esp{Solicitamos que...}.
\item[$\square$] Je connais les formules de politesse administrative : \esp{Estimado señor Ministro}, \esp{Atentamente}.
\item[$\square$] Je sais présenter oralement un pays hispanophone en quatre parties (géographie, histoire, culture, économie).
\item[$\square$] J'utilise des connecteurs pour organiser mon propos : \esp{en primer lugar, además, por último}.
\item[$\square$] Je n'oublie pas les accents (\esp{petición, sindicato, opinión}) ni les signes ¿ ¡.
\item[$\square$] J'évite les calques : « réclamer » se dit \esp{reclamar}, pas \esp{reclamar por} ; « exiger que » + subjonctif.
\item[$\square$] Je peux argumenter une revendication avec au moins trois arguments reliés logiquement.
\item[$\square$] Je relis ma production : accords, conjugaisons, ponctuation espagnole.
\end{itemize}
\end{aretenir}

\findecours

\end{document}
